% Autrui — Philosophie 1re Tech — Cours signé OnBuch+
% Programme officiel MINESEC. Compiler avec : tectonic N06-autrui.tex
\newcommand{\DOCMATIERE}{Philosophie}
\newcommand{\DOCNIVEAU}{1re Tech}
\newcommand{\DOCMODULE}{Philosophie – classe de Première technique}
\newcommand{\DOCLECON}{N06}
\newcommand{\DOCTITRE}{Autrui}
% OnBuch+ — préambule « cours signés » ENRICHI (compilé avec tectonic / XeLaTeX)
% Système visuel « Pop » d'OnBuch+ : fond crème (jamais blanc), bordures encre
% épaisses, ombres « dures » décalées, titres Archivo Black en capitales.
% Version MATHÉMATIQUES : graphes (pgfplots/tikz), boîtes pédagogiques
% (définition, propriété, méthode, attention, le savais-tu, exercice résolu,
% exercice, TP, bilan), page de garde et signature OnBuch+ ancrée en bas.
%
% Macros attendues dans main.tex AVANT \input{preamble} :
%   \DOCMATIERE  \DOCNIVEAU  \DOCMODULE  \DOCLECON (numéro)  \DOCTITRE
\documentclass[11pt]{article}

\usepackage{fontspec}
\usepackage[french]{babel}
\usepackage[a4paper,margin=2cm,top=3.4cm,bottom=2.8cm,headheight=1.4cm,footskip=1.3cm]{geometry}
\usepackage{amsmath,amssymb,amsfonts}
\usepackage{mathtools} % \xmapsto, \coloneqq... souvent utilisés par les modèles
\usepackage{cancel} % simplifications barrées (bilans de forces)
% \abs{z} (module, valeur absolue) et \norm{u} (norme), utilisés par les modèles
\providecommand{\abs}{}\let\abs\relax\DeclarePairedDelimiter{\abs}{\lvert}{\rvert}
\providecommand{\norm}{}\let\norm\relax\DeclarePairedDelimiter{\norm}{\lVert}{\rVert}
\usepackage[table]{xcolor} % [table] : fournit aussi \rowcolors (alternance de lignes)
\usepackage{pagecolor}
\usepackage{tikz}
\usetikzlibrary{arrows.meta,positioning,calc,shapes.geometric,shapes.misc,shapes.arrows,decorations.pathmorphing,decorations.pathreplacing,decorations.markings,patterns,shadows,mindmap,trees,fit,backgrounds,matrix,intersections,angles,quotes}
\usepackage{pgfplots}
\usepackage[version=4]{mhchem} % équations nucléaires \ce{^{235}_{92}U}
\usepackage[european,siunitx]{circuitikz} % schémas électriques
\pgfplotsset{compat=1.18}
\usepgfplotslibrary{fillbetween,groupplots} % aires entre courbes (calcul intégral)
\usepackage{enumitem}
\usepackage{fancyhdr}
% [most] charge toutes les bibliothèques tcolorbox (listings, minted, raster,
% documentation...) et gonfle inutilement la mémoire TeX sur un document long
% (~300 boîtes) : on ne charge que ce qui est réellement utilisé.
\usepackage[skins,breakable]{tcolorbox}
\usepackage{graphicx}
\usepackage{colortbl}
\usepackage{booktabs}
\usepackage{tabularx}
\usepackage{multirow}
\usepackage{diagbox} % case d'en-tête diagonale des tableaux à double entrée
% Colonnes extensibles centrée / gauche / droite (C, L, R), souvent utilisées par les modèles
\newcolumntype{C}{>{\centering\arraybackslash}X}
\newcolumntype{L}{>{\raggedright\arraybackslash}X}
\newcolumntype{R}{>{\raggedleft\arraybackslash}X}
\usepackage{array}
\usepackage{multicol}
\usepackage{siunitx}
% parse-numbers=false : accepte aussi \SI{2,5 \times 10^{-3}}{...}
\sisetup{locale=FR,output-decimal-marker={,},group-minimum-digits=4,parse-numbers=false}
\DeclareSIUnit\ohm{\ensuremath{\symup{\Omega}}} % U+2126 absent de la police
\DeclareSIUnit\division{div} % réglages d'oscilloscope (ms/div, V/div)
\DeclareSIUnit\tr{tr} % tours (vitesse de rotation, ex. \SI{1500}{\tr\per\minute})
\DeclareSIUnit\molar{M} % concentration molaire (ex. \SI{3}{\molar}, \SI{1}{\micro\molar})
\DeclareSIUnit\an{an} % année (le modèle confond parfois avec \year, un compteur TeX) — ex. \SI{5}{\kilo\meter\per\an}
\DeclareSIUnit\FCFA{FCFA} % franc CFA (ex. \SI{500}{\FCFA\per\kilo\gram})
\DeclareSIUnit\degreeBrix{\textdegree Brix} % teneur en sucre d'un sirop/jus (réfractométrie)
\DeclareSIUnit\month{mois} % le modèle utilise parfois \month, un compteur TeX, comme unité de durée
\DeclareSIUnit\microliter{\micro\liter} % le modèle écrit parfois \microliter en un seul token
\DeclareSIUnit\million{millions} % dénombrement (ex. \SI{15}{\million\per\mL} spermatozoïdes)
\usepackage{qrcode}
% \degree : commande fréquemment utilisée par les modèles pour un angle ou une
% température en dehors de \SI{}{} (non fournie par les paquets chargés).
\providecommand{\degree}{^{\circ}}
% \qed (fin de démonstration), utilisé par les modèles sans amsthm
\providecommand{\qed}{\hfill$\square$}
% \barre : double barre des valeurs interdites dans un tableau de variations (tkz-tab)
\providecommand{\barre}{\ensuremath{\|}}
% environnement proof (amsthm non chargé) utilisé par les modèles
\ifdefined\proof\else\newenvironment{proof}[1][Démonstration]{\par\noindent\textit{#1.}\ }{\qed\par}\fi
\usepackage{lastpage}
\usepackage{hyperref}
\hypersetup{hidelinks}

% ============================================================
% Palette « Pop » OnBuch+ — mêmes hex que la classe P (plus_ui.dart)
% ============================================================
\definecolor{poporange}{HTML}{F59321}
\definecolor{popdark}{HTML}{E07A0C}
\definecolor{popcream}{HTML}{FFF6E8}
\definecolor{popink}{HTML}{1C1714}
\definecolor{poppurple}{HTML}{7A5AE0}
\definecolor{popgreen}{HTML}{2E9E6A}
\definecolor{poppink}{HTML}{E0457B}
\definecolor{popblue}{HTML}{2F7DE1}
\definecolor{popgold}{HTML}{C98A12}
\definecolor{popmuted}{HTML}{8A7F76}
\definecolor{popline}{HTML}{EADFD2}
\definecolor{popgray}{HTML}{8A7F76}
\colorlet{popgrey}{popgray}
% Alias tolérants (le modèle nomme parfois les couleurs différemment)
\colorlet{popwhite}{white}
\colorlet{popblack}{popink}
\colorlet{popred}{poppink}
\colorlet{popyellow}{popgold}
% Teintes claires (fonds de tableaux/encadrés) — le modèle nomme parfois
% les couleurs avec un suffixe L (ex. popblueL) sans les avoir définies.
\colorlet{poporangeL}{poporange!20}
\colorlet{popdarkL}{popdark!20}
\colorlet{poppurpleL}{poppurple!20}
\colorlet{popgreenL}{popgreen!20}
\colorlet{poppinkL}{poppink!20}
\colorlet{popblueL}{popblue!20}
\colorlet{popgoldL}{popgold!20}
\colorlet{popredL}{popred!20}
\colorlet{popgrayL}{popgray!20}
% Teintes claires pour fonds de tableaux / zones
\colorlet{popblueL}{popblue!10!white}
\colorlet{poporangeL}{poporange!14!white}
\colorlet{popgreenL}{popgreen!12!white}
\colorlet{poppurpleL}{poppurple!12!white}
\colorlet{poppinkL}{poppink!12!white}
\colorlet{popgoldL}{popgold!14!white}

\pagecolor{popcream}

% ============================================================
% Typographie
% ============================================================
\defaultfontfeatures{Ligatures=TeX}
\setmainfont{PlusJakartaSans-Medium.ttf}[Path=../../../fonts/,BoldFont=PlusJakartaSans-Bold.ttf]
% Maths en sans-serif (Fira Math) avec chiffres et lettres droites de Plus
% Jakarta, pour que les formules se fondent dans le texte.
\usepackage[math-style=ISO,bold-style=ISO]{unicode-math}
\setmathfont{FiraMath-Regular.otf}[Path=../../../fonts/]
\setmathfont{PlusJakartaSans-Medium.ttf}[Path=../../../fonts/,range={up/num,up/{Latin,latin},\mathup}]
\usepackage{icomma} % virgule décimale sans espace en mode math
% Caractères mathématiques Unicode absents des polices de texte (Plus Jakarta,
% Archivo Black) : rendus par leur équivalent mathématique, en texte comme en
% formule — sinon ils s'affichent en carré vide (ex. « Arithmétique dans ℤ »).
\usepackage{newunicodechar}
\newunicodechar{ℝ}{\ensuremath{\mathbb{R}}}
\newunicodechar{ℤ}{\ensuremath{\mathbb{Z}}}
\newunicodechar{ℂ}{\ensuremath{\mathbb{C}}}
\newunicodechar{ℕ}{\ensuremath{\mathbb{N}}}
\newunicodechar{ℚ}{\ensuremath{\mathbb{Q}}}
\newunicodechar{𝔻}{\ensuremath{\mathbb{D}}}
\newunicodechar{α}{\ensuremath{\alpha}}
\newunicodechar{β}{\ensuremath{\beta}}
\newunicodechar{γ}{\ensuremath{\gamma}}
\newunicodechar{δ}{\ensuremath{\delta}}
\newunicodechar{ε}{\ensuremath{\varepsilon}}
\newunicodechar{θ}{\ensuremath{\theta}}
\newunicodechar{λ}{\ensuremath{\lambda}}
\newunicodechar{μ}{\ensuremath{\mu}}
\newunicodechar{σ}{\ensuremath{\sigma}}
\newunicodechar{τ}{\ensuremath{\tau}}
\newunicodechar{φ}{\ensuremath{\varphi}}
\newunicodechar{ψ}{\ensuremath{\psi}}
\newunicodechar{ω}{\ensuremath{\omega}}
\newunicodechar{Σ}{\ensuremath{\Sigma}}
% majuscules produites par \MakeUppercase dans les titres (α-aminés -> Α)
\newunicodechar{Α}{\ensuremath{\alpha}}
\newunicodechar{Β}{\ensuremath{\beta}}
\newunicodechar{Γ}{\ensuremath{\Gamma}}
\newunicodechar{Δ}{\ensuremath{\Delta}}
\newunicodechar{Ε}{\ensuremath{\varepsilon}}
\newunicodechar{Θ}{\ensuremath{\Theta}}
\newunicodechar{Λ}{\ensuremath{\Lambda}}
\newunicodechar{Μ}{\ensuremath{\mu}}
\newunicodechar{Τ}{\ensuremath{\tau}}
\newunicodechar{Φ}{\ensuremath{\Phi}}
\newunicodechar{Ψ}{\ensuremath{\Psi}}
\newunicodechar{Ω}{\ensuremath{\Omega}}
\newunicodechar{↦}{\ensuremath{\mapsto}}
\newunicodechar{∈}{\ensuremath{\in}}
\newunicodechar{∉}{\ensuremath{\notin}}
\newunicodechar{∧}{\ensuremath{\wedge}}
\newunicodechar{⇔}{\ensuremath{\Leftrightarrow}}
\newunicodechar{⟺}{\ensuremath{\Longleftrightarrow}}
\newunicodechar{⇒}{\ensuremath{\Rightarrow}}
\newunicodechar{⟹}{\ensuremath{\Longrightarrow}}
\newunicodechar{∘}{\ensuremath{\circ}}
\newunicodechar{∪}{\ensuremath{\cup}}
\newunicodechar{∩}{\ensuremath{\cap}}
\newunicodechar{≡}{\ensuremath{\equiv}}
\newunicodechar{⊂}{\ensuremath{\subset}}
\newunicodechar{⊥}{\ensuremath{\perp}}
\newunicodechar{∃}{\ensuremath{\exists}}
\newunicodechar{∀}{\ensuremath{\forall}}
\newunicodechar{∛}{\ensuremath{\sqrt[3]{\,}}}
\newunicodechar{∜}{\ensuremath{\sqrt[4]{\,}}}
\newunicodechar{⃗}{\ensuremath{{}^{\to}}}
\newunicodechar{ⁿ}{\ifmmode^{n}\else\textsuperscript{\ensuremath{n}}\fi}
\newunicodechar{ˣ}{\ifmmode^{x}\else\textsuperscript{\ensuremath{x}}\fi}
\newunicodechar{ᵇ}{\ifmmode^{b}\else\textsuperscript{\ensuremath{b}}\fi}
\newunicodechar{ʳ}{\ifmmode^{r}\else\textsuperscript{\ensuremath{r}}\fi}
\newunicodechar{ᵘ}{\ifmmode^{u}\else\textsuperscript{\ensuremath{u}}\fi}
\newunicodechar{ʸ}{\ifmmode^{y}\else\textsuperscript{\ensuremath{y}}\fi}
\newunicodechar{ᵃ}{\ifmmode^{a}\else\textsuperscript{\ensuremath{a}}\fi}
\newunicodechar{ᵉ}{\ifmmode^{e}\else\textsuperscript{\ensuremath{e}}\fi}
\newunicodechar{ᵏ}{\ifmmode^{k}\else\textsuperscript{\ensuremath{k}}\fi}
\newunicodechar{ᵖ}{\ifmmode^{p}\else\textsuperscript{\ensuremath{p}}\fi}
\newunicodechar{ᴺ}{\ifmmode^{N}\else\textsuperscript{\ensuremath{N}}\fi}
\newunicodechar{ᶜ}{\ifmmode^{c}\else\textsuperscript{\ensuremath{c}}\fi}
\newunicodechar{ʰ}{\ifmmode^{h}\else\textsuperscript{\ensuremath{h}}\fi}
\newunicodechar{ᵠ}{\ifmmode^{q}\else\textsuperscript{\ensuremath{q}}\fi}
\newunicodechar{ₙ}{\ifmmode_{n}\else\textsubscript{\ensuremath{n}}\fi}
\newunicodechar{ₐ}{\ifmmode_{a}\else\textsubscript{\ensuremath{a}}\fi}
\newunicodechar{ᵢ}{\ifmmode_{i}\else\textsubscript{\ensuremath{i}}\fi}
\newunicodechar{ᵣ}{\ifmmode_{r}\else\textsubscript{\ensuremath{r}}\fi}
\newunicodechar{ₖ}{\ifmmode_{k}\else\textsubscript{\ensuremath{k}}\fi}
\newunicodechar{ᵦ}{\ifmmode_{\beta}\else\textsubscript{\ensuremath{\beta}}\fi}
\newunicodechar{ₓ}{\ifmmode_{x}\else\textsubscript{\ensuremath{x}}\fi}
\newfontfamily\popdisplay{ArchivoBlack-Regular.ttf}[Path=../../../fonts/]
\newfontfamily\poplogo{SpaceGrotesk-Bold.ttf}[Path=../../../fonts/]
\newfontfamily\popsemi{PlusJakartaSans-SemiBold.ttf}[Path=../../../fonts/]
\newfontfamily\popxbold{PlusJakartaSans-ExtraBold.ttf}[Path=../../../fonts/]
\newfontfamily\popxboldls{PlusJakartaSans-ExtraBold.ttf}[Path=../../../fonts/,LetterSpace=3.0]

\setlist[enumerate]{leftmargin=1.7em,itemsep=5pt,topsep=4pt}
\setlist[itemize]{leftmargin=1.7em,itemsep=4pt,topsep=4pt,label={\color{poporange}\textbullet}}
\setlength{\parindent}{0pt}
\setlength{\parskip}{5pt}
\renewcommand{\arraystretch}{1.25}

% pgfplots : style Pop par défaut
\pgfplotsset{
  every axis/.append style={axis line style={popink,line width=1pt},tick style={popink},
    grid style={popline,line width=0.5pt},label style={font=\small},tick label style={font=\footnotesize}},
  popcurve/.style={poporange,line width=1.8pt,no markers,smooth},
}
% Même style disponible dans un \draw TikZ simple (hors pgfplots)
\tikzset{popcurve/.style={poporange,line width=1.8pt}}
\tikzset{popgrid/.style={popline,very thin}}
\colorlet{popcurve}{poporange} % le nom du style est parfois utilisé comme couleur

% ============================================================
% Pill / squiggle
% ============================================================
\newcommand{\poppill}[3]{%
  \tikz[baseline=-0.55ex]\node[draw=popink,line width=1.2pt,rounded corners=2.6pt,%
    fill=#2,inner xsep=6.5pt,inner ysep=3pt]{%
    \popxboldls\fontsize{7}{7}\selectfont\color{#3}\MakeUppercase{#1}};%
}
\newcommand{\popsquiggle}[1]{%
  \tikz[baseline=-0.4ex]{
    \draw[#1,line width=2.6pt,line cap=round]
      plot [smooth] coordinates {(0,0.09) (0.34,0.01) (0.68,0.21) (1.02,0.01) (1.36,0.10)};
  }%
}
% Mot-clé surligné (marqueur orange sous le texte)
\newcommand{\cle}[1]{{\popxbold\color{popdark}#1}}

% ============================================================
% En-tête / pied
% ============================================================
\pagestyle{fancy}
\fancyhf{}
\renewcommand{\headrulewidth}{0pt}
\renewcommand{\footrulewidth}{0pt}
\lhead{\raisebox{-0.15ex}{\poplogo\fontsize{13}{13}\selectfont\color{popink}OnBuch\color{poporange}+}}
\rhead{\poppill{\DOCMATIERE\ \textbullet\ \DOCNIVEAU\ \textbullet\ Leçon \DOCLECON}{white}{popink}}
\lfoot{\popxbold\fontsize{7.6}{7.6}\selectfont\color{popmuted}\MakeUppercase{Cours signé OnBuch+}\ \textbullet\ {\popsemi\DOCTITRE}}
\rfoot{\tikz[baseline=-0.6ex]\node[rounded corners=8.5pt,fill=popink,minimum height=17pt,minimum width=17pt,inner xsep=5pt,inner sep=0]{\popxbold\fontsize{8}{8}\selectfont\color{popcream}\thepage\,/\,\pageref*{LastPage}};}

% ============================================================
% Titres
% ============================================================
\newcommand{\chaptitre}[2]{%
  \begin{tcolorbox}[enhanced,colback=poporange,colframe=popink,boxrule=2.6pt,arc=11pt,
      drop shadow=popink,left=15pt,right=15pt,top=12pt,bottom=11pt,before skip=3pt,after skip=18pt]
    \poppill{#2}{popink}{popcream}\par\vspace{9pt}
    {\raggedright\hyphenpenalty=10000\exhyphenpenalty=10000\popdisplay\fontsize{22}{24}\selectfont\color{popcream}\MakeUppercase{#1}\par}
  \end{tcolorbox}
}
% Section numérotée : pastille orange avec le numéro + titre Archivo
\newcounter{popsec}
\newcommand{\coursec}[1]{%
  \stepcounter{popsec}\setcounter{popsub}{0}%
  \par\vspace{16pt}\noindent
  \begin{minipage}{\linewidth}
  \tikz[baseline=-0.7ex]\node[draw=popink,line width=1.6pt,fill=poporange,rounded corners=4pt,minimum size=22pt,inner sep=0]{\popdisplay\fontsize{12}{12}\selectfont\color{popcream}\thepopsec};%
  \hspace{8pt}{\popdisplay\fontsize{15}{17}\selectfont\color{popink}\MakeUppercase{#1}}\par
  \vspace{4pt}\noindent{\color{poporange}\rule{36pt}{3.6pt}}
  \end{minipage}\par\nopagebreak\vspace{8pt}
}
\newcounter{popsub}
\newcommand{\courssub}[1]{%
  \stepcounter{popsub}%
  \par\vspace{9pt}\noindent{\popxbold\fontsize{11.5}{13}\selectfont\color{popink}{\color{poporange}\thepopsec.\thepopsub}\hspace{6pt}#1}\par\nopagebreak\vspace{4pt}
}

% ============================================================
% Boîtes pédagogiques — même recette : fond blanc, bordure épaisse colorée,
% ombre dure encre, badge plein. Titre optionnel [#1].
% ============================================================
\tcbset{popbase/.style={breakable,enhanced,colback=white,boxrule=2.3pt,arc=9pt,
  shadow={3pt}{-3pt}{0pt}{popink},left=14pt,right=14pt,top=11pt,bottom=11pt,before skip=11pt,after skip=15pt,
  fontupper=\popsemi\fontsize{10.6}{14.5}\selectfont\color{popink},
  fontlower=\popsemi\fontsize{10.6}{14.5}\selectfont\color{popink}}}
% Coche dessinée (le glyphe ✓ n'existe pas dans les polices du document)
\newcommand{\popcheck}[1]{\tikz[baseline=-0.5ex]\draw[#1,line width=1.6pt,line cap=round,line join=round](-0.13,0.0)--(-0.04,-0.09)--(0.13,0.1);}
\providecommand{\coche}{\popcheck{popgreen}}
\providecommand{\resultat}[1]{\par\smallskip\noindent\fcolorbox{popgreen}{popgreenL}{\parbox{\dimexpr\linewidth-2\fboxsep-2\fboxrule\relax}{\textbf{Résultat.} #1}}\par\smallskip}
\newcommand{\popboxhead}[3]{\poppill{#1}{#2}{white}\ifx\relax#3\relax\else\hspace{6pt}{\popxbold\fontsize{10.6}{12}\selectfont\color{popink}#3}\fi\par\vspace{7pt}}

\newtcolorbox{aretenir}[1][]{popbase,colframe=popgreen,before upper={\popboxhead{À retenir}{popgreen}{#1}}}
\newtcolorbox{exemplebox}[1][]{popbase,colframe=poporange,before upper={\popboxhead{Exemple}{poporange}{#1}}}
\newtcolorbox{definition}[1][]{popbase,colframe=popblue,before upper={\popboxhead{Définition}{popblue}{#1}}}
\newtcolorbox{propriete}[1][]{popbase,colframe=poppurple,before upper={\popboxhead{Propriété}{poppurple}{#1}}}
\newtcolorbox{methode}[1][]{popbase,colframe=popgold,before upper={\popboxhead{Méthode}{popgold}{#1}}}
\newtcolorbox{attention}[1][]{popbase,colframe=poppink,before upper={\popboxhead{Attention}{poppink}{#1}}}
\newtcolorbox{experience}[1][]{popbase,colframe=popblue,colback=popblueL,before upper={\popboxhead{Expérience}{popblue}{#1}}}
% « Le savais-tu ? » : carte violette pleine (contexte, culture, Cameroun)
\newtcolorbox{savaistu}[1][]{popbase,colback=poppurple,colframe=popink,
  fontupper=\popsemi\fontsize{10.4}{14.2}\selectfont\color{white},
  before upper={\poppill{Le savais-tu\,?}{white}{poppurple}\ifx\relax#1\relax\else\hspace{6pt}{\popxbold\color{white}#1}\fi\par\vspace{7pt}}}
% Exercice résolu : énoncé en haut, \tcblower puis la solution sur fond crème
\newcounter{popexo}\newcounter{popexr}
\newtcolorbox{exoresolu}[1][]{popbase,colframe=poporange,colbacklower=popcream,
  segmentation style={popink,line width=1.2pt,dashed},
  before upper={\stepcounter{popexr}\popboxhead{Exercice résolu \thepopexr}{poporange}{#1}},
  before lower={\poppill{Solution}{popink}{popcream}\par\vspace{6pt}}}
% Exercice (non résolu en place) — la correction est dans la partie Corrigés
\newtcolorbox{exercice}[1][]{popbase,colframe=popink,
  before upper={\stepcounter{popexo}\popboxhead{Exercice \thepopexo}{popink}{#1}}}
% Corrigé
\newtcolorbox{corrige}[1][]{popbase,colframe=popgreen,colback=popgreenL,
  before upper={\popboxhead{Corrigé}{popgreen}{#1}}}
% Objectifs / prérequis / bilan
\newtcolorbox{objectifs}{popbase,colframe=popink,colback=poporangeL,
  before upper={\popboxhead{Objectifs de la leçon}{popink}{}}}
\newtcolorbox{prerequis}{popbase,colframe=popink,colback=popblueL,
  before upper={\popboxhead{Prérequis}{popblue}{}}}
\newtcolorbox{bilan}{popbase,colframe=popink,colback=white,boxrule=2.8pt,
  before upper={\popboxhead{Fiche bilan}{poporange}{L'essentiel à connaître pour l'examen}}}
% Figure encadrée avec légende
\newtcolorbox{popfigure}[1][]{enhanced,breakable=false,colback=white,colframe=popline,boxrule=1.4pt,arc=8pt,
  left=10pt,right=10pt,top=10pt,bottom=8pt,before skip=10pt,after skip=12pt,halign=center,
  lower separated=false,halign lower=center,
  fontlower=\popsemi\fontsize{9}{11.5}\selectfont\color{popmuted},
  #1}
\newcommand{\legende}[1]{\par\vspace{6pt}{\popsemi\fontsize{9}{11.5}\selectfont\color{popmuted}#1}}

% ============================================================
% Signature OnBuch+
% ============================================================
\newcommand{\poplogoinline}[1]{{\poplogo\fontsize{#1}{#1}\selectfont\color{popink}OnBuch\color{poporange}+}}
\newcommand{\signatureonbuch}{%
  \begin{tcolorbox}[enhanced,colback=popink,colframe=popink,boxrule=2.2pt,arc=10pt,
      drop shadow=poporange,left=14pt,right=14pt,top=10pt,bottom=10pt,before skip=0pt,after skip=0pt]
    \tikz[baseline=-0.7ex]\node[circle,fill=popgreen,draw=popcream,line width=1.3pt,minimum size=22pt,inner sep=0]{\popcheck{white}};%
    \hspace{9pt}\begin{minipage}[c]{0.62\linewidth}
      {\popxbold\fontsize{12}{13}\selectfont\color{popcream}Signé\ \textbullet\ {\poplogo OnBuch\color{poporange}+}}\par\vspace{2pt}
      {\raggedright\popsemi\fontsize{8}{10.5}\selectfont\color{popline}\MakeUppercase{Équipe pédagogique OnBuch+}\\\MakeUppercase{Conforme au programme officiel MINESEC}\par}
    \end{minipage}\hfill
    {\poplogo\fontsize{22}{22}\selectfont\color{popcream}OnBuch\color{poporange}+}
  \end{tcolorbox}%
}

% ============================================================
% Page de garde
% #1 titre · #2 matière · #3 niveau · #4 module · #5 n° leçon · #6 durée ·
% #7 description / objectifs (texte ou itemize)
% ============================================================
\newcommand{\pagegarde}[7]{%
  \thispagestyle{empty}
  \newgeometry{a4paper,margin=1.7cm,top=1.6cm,bottom=1.4cm}%
  \begin{tikzpicture}[remember picture,overlay]
    \fill[poporange!18] ([xshift=-3.2cm,yshift=-3.6cm]current page.north east) circle (4.6cm);
    \fill[poppurple!14] ([xshift=2.4cm,yshift=7.6cm]current page.south west) circle (3.8cm);
  \end{tikzpicture}%
  {\poplogo\fontsize{30}{30}\selectfont\color{popink}OnBuch\color{poporange}+}\hfill
  \raisebox{0.6ex}{\poppill{Cours signé \textbullet\ Premium}{popink}{popcream}}\par
  \vspace{1pt}\popsquiggle{poppurple}\par
  \vspace{8pt}
  \begin{tcolorbox}[enhanced,colback=poporange,colframe=popink,boxrule=2.8pt,arc=13pt,
      drop shadow=popink,left=18pt,right=18pt,top=12pt,bottom=13pt,after skip=10pt]
    \poppill{#2 \textbullet\ #3}{popink}{popcream}\hspace{5pt}\poppill{#4}{white}{popink}\par\vspace{8pt}
    {\popdisplay\fontsize{14}{14}\selectfont\color{popink}LEÇON #5}\par\vspace{4pt}
    {\raggedright\hyphenpenalty=10000\exhyphenpenalty=10000\popdisplay\fontsize{25}{27}\selectfont\color{popcream}\MakeUppercase{#1}\par}
    \vspace{8pt}
    {\popxbold\fontsize{9.5}{9.5}\selectfont\color{popink}\MakeUppercase{Durée conseillée : #6}}
  \end{tcolorbox}
  \begin{tcolorbox}[enhanced,colback=white,colframe=popink,boxrule=2.2pt,arc=10pt,
      drop shadow=popink,left=16pt,right=16pt,top=10pt,bottom=10pt,after skip=8pt]
    \poppill{Dans cette leçon}{poppurple}{white}\par\vspace{5pt}
    \popsemi\fontsize{9.3}{12.6}\selectfont\color{popink}#7
  \end{tcolorbox}
  \noindent\begin{minipage}[t]{0.487\linewidth}
    \begin{tcolorbox}[enhanced,colback=popgreen,colframe=popink,boxrule=2.2pt,arc=10pt,
        drop shadow=popink,left=11pt,right=11pt,top=10pt,bottom=10pt,height=3.1cm,valign=center]
      \tcbox[colback=white,colframe=popink,boxrule=1.3pt,arc=5pt,boxsep=0pt,nobeforeafter,
        left=3pt,right=3pt,top=3pt,bottom=3pt,valign=center]{\qrcode[height=1.9cm]{https://play.google.com/store/apps/details?id=com.luvvix.onbuch&hl=fr}}%
      \hspace{8pt}\begin{minipage}[c]{0.5\linewidth}
        {\popdisplay\fontsize{11}{12.5}\selectfont\color{white}\MakeUppercase{Télécharge OnBuch}}\par\vspace{4pt}
        {\raggedright\popsemi\fontsize{8.4}{11}\selectfont\color{white}Gratuit \textbullet\ Google Play\\Tuteur IA, annales, concours, résultats\par}
      \end{minipage}
    \end{tcolorbox}
  \end{minipage}\hfill
  \begin{minipage}[t]{0.487\linewidth}
    \begin{tcolorbox}[enhanced,colback=poppurple,colframe=popink,boxrule=2.2pt,arc=10pt,
        drop shadow=popink,left=11pt,right=11pt,top=10pt,bottom=10pt,height=3.1cm,valign=center]
      \tikz[baseline=-0.7ex]\node[circle,fill=white,draw=popink,line width=1.3pt,minimum size=1.6cm,inner sep=0]{\popdisplay\fontsize{24}{24}\selectfont\color{poppurple}+};%
      \hspace{8pt}\begin{minipage}[c]{0.6\linewidth}
        {\popdisplay\fontsize{11}{12.5}\selectfont\color{white}\MakeUppercase{Passe à OnBuch+}}\par\vspace{4pt}
        {\raggedright\popsemi\fontsize{8.4}{11}\selectfont\color{white}Tous les cours signés, Tuteur IA illimité, fiches PDF — onglet OnBuch+ de l'app\par}
      \end{minipage}
    \end{tcolorbox}
  \end{minipage}
  \vspace{8pt}
  \popcontenu
  \vfill
  \signatureonbuch
  \restoregeometry
  \newpage
}

% Rangée « Ce cours contient » (page de garde)
\newcommand{\poptile}[3]{% #1 couleur #2 gros texte #3 légende
  \begin{tcolorbox}[enhanced,colback=#1,colframe=popink,boxrule=2pt,arc=9pt,shadow={2.5pt}{-2.5pt}{0pt}{popink},
      left=6pt,right=6pt,top=9pt,bottom=9pt,halign=center,valign=center,height=1.75cm,width=\linewidth,nobeforeafter]
    {\popdisplay\fontsize{12.5}{13}\selectfont\color{white}\MakeUppercase{#2}}\par\vspace{3pt}
    {\centering\popsemi\fontsize{7.8}{9.5}\selectfont\color{white}#3\par}
  \end{tcolorbox}}
\newcommand{\popcontenu}{%
  \par\noindent{\popxboldls\fontsize{7.5}{7.5}\selectfont\color{popmuted}CE COURS SIGNÉ CONTIENT}\par\vspace{7pt}
  \noindent\begin{minipage}[t]{0.235\linewidth}\poptile{popblue}{Cours}{complet, illustré, pas à pas}\end{minipage}\hfill
  \begin{minipage}[t]{0.235\linewidth}\poptile{poporange}{Méthodes}{et exercices résolus}\end{minipage}\hfill
  \begin{minipage}[t]{0.235\linewidth}\poptile{poppink}{Exercices}{type examen + corrigés}\end{minipage}\hfill
  \begin{minipage}[t]{0.235\linewidth}\poptile{popgreen}{Fiche bilan}{l'essentiel à retenir}\end{minipage}\par}

% Sommaire
\newtcolorbox{sommaire}{enhanced,colback=white,colframe=popink,boxrule=2.2pt,arc=10pt,
  drop shadow=popink,left=18pt,right=18pt,top=13pt,bottom=13pt,before skip=6pt,after skip=20pt,
  before upper={\poppill{Sommaire}{poppurple}{white}\par\vspace{9pt}\popsemi\fontsize{10.6}{14.5}\selectfont\color{popink}}}

% Signature de fin de cours — poussée tout en bas de la dernière page
\newcommand{\findecours}{%
  \par\vfill\nopagebreak
  \begin{center}\popsquiggle{poporange}\end{center}\vspace{4pt}
  \signatureonbuch
}
\AtBeginDocument{\def\llbracket{\mathopen{[\mkern-3.5mu[}}\def\rrbracket{\mathclose{]\mkern-3.5mu]}}} % intervalles d'entiers (glyphe ⟦ absent de la police)
% Glyphes absents de Fira Math : redéfinis à partir de glyphes disponibles
\AtBeginDocument{%
  \def\setminus{\mathbin{\backslash}}\let\smallsetminus\setminus
  \def\vdots{\mathord{\vbox{\baselineskip3.2pt\lineskiplimit0pt\kern4pt\hbox{.}\hbox{.}\hbox{.}}}}%
  \def\ddots{\mathinner{\mkern1mu\raise6.5pt\hbox{.}\mkern2mu\raise3.6pt\hbox{.}\mkern2mu\raise0.7pt\hbox{.}\mkern1mu}}%
  \def\triangle{\mathord{\scalebox{0.9}{$\symup{\Delta}$}}}%
  \def\nabla{\mathord{\raisebox{\depth}{\scalebox{1}[-1]{$\symup{\Delta}$}}}}%
  \def\checkmark{\popcheck{popdark}}%
  \def\star{\mathbin{\ast}}\def\bigstar{\mathord{\ast}}%
  \def\diamond{\mathbin{\scalebox{0.6}{$\Diamond$}}}%
  \def\bigcup{\mathop{\mathchoice{\raisebox{-0.3em}{\scalebox{1.7}{$\cup$}}}{\scalebox{1.25}{$\cup$}}{\cup}{\cup}}\displaylimits}%
  \def\bigcap{\mathop{\mathchoice{\raisebox{-0.3em}{\scalebox{1.7}{$\cap$}}}{\scalebox{1.25}{$\cap$}}{\cap}{\cap}}\displaylimits}%
  \def\lesseqqgtr{\mathrel{\lessgtr}}\def\bigoplus{\mathop{\oplus}\displaylimits}%
}
\providecommand{\ssi}{si et seulement si} % abréviation fréquente
\AtBeginDocument{\providecommand{\wideparen}{\overparen}} % arc de cercle

%% FIGFIX-BEGIN v4 — correctifs globaux des figures (docs/onbuch-generation/tools/figfix)
\usepackage{adjustbox}\usepackage{etoolbox}
% toute figure TikZ/pgfplots plus large que la ligne est réduite à la largeur disponible
\BeforeBeginEnvironment{tikzpicture}{\begin{adjustbox}{max width=\linewidth}}
\AfterEndEnvironment{tikzpicture}{\end{adjustbox}}
% légendes pgfplots lisibles par-dessus les courbes
\pgfplotsset{every axis/.append style={legend style={fill=white,fill opacity=0.92,text opacity=1,draw=black!15,inner sep=3pt}}}
% étiquettes d'arêtes (syntaxe edge["..."]) avec fond blanc
\tikzset{every edge quotes/.append style={fill=white,inner sep=1.2pt}}
%% FIGFIX-END
\begin{document}

\pagegarde{Autrui}{Philosophie}{1re Tech}{Philosophie – classe de Première technique}{N06}{2 séances (fiche de progression harmonisée)}{Cette leçon examine comment nous connaissons autrui (perception, analogie, empathie, langage) et analyse l'intersubjectivité, c'est-à-dire la nature des rapports entre les consciences : conflit, reconnaissance, amour, justice. Elle conduit l'élève à appliquer ces notions aux situations concrètes de la vie sociale camerounaise et à construire une argumentation philosophique rigoureuse.
\vspace{4pt}
\begin{multicols}{2}\small
\begin{itemize}
\item À la fin de cette leçon, je sais définir autrui et le distinguer de l'objet et de la chose
\item Identifier et expliquer les différents modes de connaissance d'autrui (perception, raisonnement analogique, empathie, langage)
\item Analyser le problème philosophique de l'accès à la conscience d'autrui
\item Définir l'intersubjectivité et décrire ses formes (communication, reconnaissance, conflit, amour)
\item Confronter les thèses de Sartre (conflit) et de Hegel (reconnaissance) sur les rapports avec autrui
\item Appliquer les notions de la leçon à des situations concrètes de la vie courante au Cameroun
\end{itemize}\end{multicols}}

\begin{sommaire}
\begin{multicols}{2}
\begin{enumerate}
\item Introduction : problématique et définitions
\item La connaissance d'autrui par la perception et l'analogie
\item Empathie, langage et communication
\item L'intersubjectivité : définition et formes
\item Autrui : obstacle ou condition de ma liberté ? (Hegel, Sartre)
\item Applications, méthode et entraînement au Bac
\item Activité d'intégration
\item Méthodes et exercices résolus
\item Exercices
\item Corrigés des exercices
\item Fiche bilan
\end{enumerate}
\end{multicols}
\end{sommaire}

\begin{objectifs}
\begin{itemize}
\item À la fin de cette leçon, je sais définir autrui et le distinguer de l'objet et de la chose
\item Identifier et expliquer les différents modes de connaissance d'autrui (perception, raisonnement analogique, empathie, langage)
\item Analyser le problème philosophique de l'accès à la conscience d'autrui
\item Définir l'intersubjectivité et décrire ses formes (communication, reconnaissance, conflit, amour)
\item Confronter les thèses de Sartre (conflit) et de Hegel (reconnaissance) sur les rapports avec autrui
\item Appliquer les notions de la leçon à des situations concrètes de la vie courante au Cameroun
\item Rédiger et justifier une démarche argumentative, une réponse ou une conclusion en dissertation
\item Mobiliser des références philosophiques précises (Husserl, Hegel, Sartre, Levinas) dans un devoir
\end{itemize}
\end{objectifs}

\begin{prerequis}
\begin{itemize}
\item Notion de conscience et de subjectivité (leçon sur la conscience)
\item Distinction entre sujet et objet
\item Notion de perception (leçon sur la perception)
\item Notions de liberté et de devoir (classes de Seconde)
\item Méthode de la dissertation philosophique : introduction, problématique, plan
\end{itemize}
\end{prerequis}

\newpage

\coursec{Introduction : problématique et définitions}

\courssub{Situation d'entrée : la scène du marché}

Au marché Mokolo à Yaoundé, un commerçant devine, sans qu'on le lui dise, que le client qui s'approche est pressé et mécontent : il lit cela sur son visage et sa démarche. Il n'a entendu aucune parole, il n'a reçu aucune explication ; pourtant, en une fraction de seconde, il a \emph{compris} quelque chose de la vie intérieure d'un inconnu. Comment est-ce possible ? Ce que le commerçant perçoit réellement, c'est un corps : des sourcils froncés, des pas rapides, une mâchoire serrée. Mais ce qu'il prétend connaître, c'est autre chose : une humeur, une intention, un état d'esprit, c'est-à-dire une \cle{conscience}. Or une conscience ne se voit pas, ne se touche pas, ne se pèse pas. Il y a donc un écart mystérieux entre ce que je perçois d'autrui (son corps) et ce que je crois savoir de lui (sa vie intérieure).

Cette scène ordinaire soulève l'un des problèmes les plus profonds de la philosophie. Autrui peut-il être un obstacle à ma liberté, comme le suggèrent les conflits quotidiens dans nos quartiers, ou est-il au contraire la condition de mon humanité ? Cette leçon interroge les \cle{modes de connaissance d'autrui} et la nature de nos \cle{rapports avec autrui}, conformément aux deux savoirs de l'unité : les modes de connaissance d'autrui, et l'intersubjectivité ou question des rapports avec autrui.

\courssub{Définition d'autrui : un autre moi-même}

Commençons par l'intuition la plus simple. Parmi tout ce qui m'entoure, je distingue spontanément trois sortes de réalités : les choses (une pierre, une table, une moto), les êtres vivants non humains (un chien, un arbre), et les autres hommes. Ces derniers occupent une place à part : je ne les traite jamais exactement comme des choses, ni exactement comme des animaux. Le mot \cle{autrui} désigne précisément cette réalité singulière : l'autre homme.

\begin{definition}[Autrui]
\cle{Autrui} est l'\emph{autre homme} : un sujet conscient \textbf{semblable à moi}, mais \textbf{autre que moi}. Je le perçois comme un objet parmi les objets du monde (un corps visible dans l'espace), mais il m'échappe comme sujet : il possède, comme moi, une vie intérieure, des pensées, des émotions et une liberté auxquelles je n'ai pas d'accès direct. La philosophie le désigne parfois par l'expression latine \emph{alter ego}, littéralement « un autre moi-même ».
\end{definition}

La définition contient deux termes en tension, et toute la difficulté du problème tient dans cette tension :

\begin{itemize}
\item Autrui est \textbf{semblable à moi} : il a un corps comme le mien, il parle, il rit, il souffre, il pense. Je le reconnais comme un être de même nature que moi, et non comme une chose. C'est pourquoi je lui dois un respect que je ne dois pas à une table.
\item Autrui est \textbf{autre que moi} : il n'est pas moi. Sa conscience m'est fermée. Je ne ressentirai jamais sa douleur exactement comme il la ressent, je ne verrai jamais le monde exactement de son point de vue. Entre lui et moi, il y a une distance irréductible.
\end{itemize}

Autrui est donc un \emph{paradoxe vivant} : le plus proche de moi (nous partageons la même humanité) et en même temps le plus inaccessible (sa conscience m'est cachée).

\begin{savaistu}[D'où vient le mot « alter ego » ?]
En latin, \emph{alter} signifie « l'autre des deux » : ce n'est pas n'importe quel autre, mais celui qui fait face au premier dans une relation à deux. \emph{Ego} signifie « moi ». L'expression \emph{alter ego}, employée déjà par les Romains pour désigner un ami intime, signifie donc littéralement « un autre moi-même ». Elle exprime parfaitement le statut philosophique d'autrui : un être qui est à la fois \emph{moi} (par sa nature de sujet) et \emph{autre} (par son existence séparée).
\end{savaistu}

\courssub{Autrui n'est ni une chose ni un simple objet : c'est un sujet}

Pour bien comprendre ce qu'est autrui, il faut le distinguer de ce qu'il n'est pas. Une \cle{chose} (une pierre, une chaise) n'a ni conscience, ni langage, ni droits : je peux l'utiliser, la déplacer, la détruire sans lui demander son avis. Un \cle{animal} possède une sensibilité, mais il ne me parle pas, ne me juge pas, ne m'assigne pas de devoirs au sens moral. Autrui, lui, est un \cle{sujet} : un être qui pense, qui dit « je », qui a ses propres projets et qui peut, à son tour, me percevoir et me juger.

\begin{center}
\begin{tabularx}{\linewidth}{|l|X|X|X|}
\hline
\rowcolor{popblueL} \textbf{Critère} & \textbf{La chose} & \textbf{L'animal} & \textbf{Autrui} \\
\hline
Corps visible & Oui & Oui & Oui \\
\hline
Conscience, vie intérieure & Non & Sensibilité sans conscience de soi réflexive & Oui, conscience de soi \\
\hline
Langage, parole & Non & Signaux limités & Oui, langage articulé \\
\hline
Droits et devoirs & Aucun & Protection limitée & Droits pleins (dignité, respect) \\
\hline
Peut me percevoir et me juger & Non & De façon élémentaire & Oui : je suis pour lui un objet \\
\hline
\end{tabularx}
\end{center}

Ce tableau fait apparaître l'essentiel : autrui partage avec la chose le fait d'avoir un corps visible, mais il partage avec \emph{moi} le fait d'être une conscience. C'est précisément ce mélange qui rend la relation à autrui unique et problématique.

\begin{attention}[Erreur fréquente : mal définir autrui]
Dans les copies, deux confusions reviennent souvent. Premièrement, confondre \emph{autrui} (un sujet précis, « l'autre homme » en face de moi) avec « les autres » au sens sociologique vague (la foule, l'opinion publique, la société). Deuxièmement, traiter autrui comme un simple objet extérieur, au même titre qu'un arbre ou qu'une pierre. Or tout le problème philosophique vient de ce qu'autrui n'est \emph{pas} un objet comme les autres : c'est un sujet. Une définition d'autrui qui oublie le mot « sujet » est une définition ratée.
\end{attention}

\courssub{Le double statut d'autrui : objet pour ma perception, sujet comme moi}

Reprenons la scène du marché. Pour mes yeux, le client qui s'approche est un \emph{objet} : un corps situé dans l'espace, que je peux décrire de l'extérieur (taille, vêtements, gestes, expressions). En ce sens, autrui fait partie du monde que je perçois, au même titre que les étals et les marchandises. Le philosophe peut dire : autrui est d'abord, pour moi, un \emph{objet de perception}.

Mais en même temps, je sais que ce corps n'est pas une machine vide : il est habité par une conscience. Ce client me voit, lui aussi ; il se fait une idée de moi ; il peut décider de partir, de mentir, de marchander. Autrement dit, je ne suis pas seulement celui qui le regarde : je suis aussi \emph{regardé par lui}. Devant autrui, je deviens moi-même un objet pour sa conscience. C'est ce renversement qui distingue radicalement la rencontre d'autrui de la perception d'une chose : une pierre ne me regarde jamais en retour.

\begin{popfigure}
\begin{center}
\begin{tikzpicture}[
  sphere/.style={circle, draw=popblue, thick, fill=popblueL, minimum size=2.4cm, align=center, font=\bfseries},
  lbl/.style={font=\small\itshape, text=popdark}
]
\node[sphere, align=center] (moi) at (0,0){Moi\\[2pt] \normalfont\small sujet conscient};
\node[sphere, draw=poporange, fill=poporangeL, align=center] (autrui) at (7.5,0){Autrui\\[2pt] \normalfont\small sujet conscient};
\draw[-{Stealth[length=3mm]}, very thick, popblue] (moi.30) to[bend left=18] node[lbl, above]{je perçois son corps (objet)} (autrui.150);
\draw[-{Stealth[length=3mm]}, very thick, poporange] (autrui.210) to[bend left=18] node[lbl, below]{il me perçoit et me juge (je deviens objet)} (moi.330);
\node[lbl, align=center] at (3.75,-2.3) {Double flèche : perception et reconnaissance mutuelles.\\ Chacun est à la fois sujet (qui voit) et objet (qui est vu).};
\end{tikzpicture}
\end{center}
\legende{Le double statut d'autrui : je le perçois comme un objet (son corps), mais il est un sujet qui, à son tour, me perçoit. La relation à autrui est donc réversible, contrairement à la perception d'une chose.}
\end{popfigure}

\begin{aretenir}[L'essentiel sur le statut d'autrui]
Autrui possède un \textbf{double statut} : il est \emph{objet} pour ma perception (je ne vois que son corps) et \emph{sujet} comme moi (il possède une conscience qui m'échappe et qui peut me saisir comme objet). Toute la philosophie d'autrui se joue dans cette dualité.
\end{aretenir}

\courssub{Le problème philosophique : l'accès à la conscience d'autrui}

Nous pouvons maintenant formuler rigoureusement le problème. Procédons comme dans une démarche d'investigation : partons d'une observation, dégageons la difficulté, puis formulons la question.

\textbf{Observation.} Je n'ai jamais perçu directement la conscience de qui que ce soit. Je vois des corps : des visages qui rougissent, des poings qui se serrent, des larmes qui coulent. Mais la colère elle-même, la douleur elle-même, la tristesse elle-même, je ne les vois pas. Quand mon voisin dit « j'ai mal », je n'éprouve pas sa douleur ; je crois seulement qu'il souffre.

\textbf{Difficulté.} De moi-même, j'ai une connaissance directe : je suis certain de penser, de douter, de souffrir, parce que je le vis de l'intérieur. D'autrui, je n'ai qu'une connaissance indirecte, par signes : je dois \emph{deviner} sa conscience à travers son corps. Mais alors, qu'est-ce qui me garantit que ces signes ne trompent pas ? Un acteur peut simuler la douleur ; un hypocrite peut cacher sa colère ; et, à la limite, rien ne me prouve de façon absolue que les corps qui m'entourent sont habités par des consciences plutôt que d'être de simples automates perfectionnés.

\textbf{Question.} Comment puis-je donc passer légitimement du corps perçu à la conscience supposée ? Autrement dit : \emph{comment connaissons-nous autrui ?}

\begin{exemplebox}[Le sourire du vendeur]
Une vendeuse du marché me sourit chaleureusement. Ce sourire signifie-t-il qu'elle est heureuse de me voir, ou est-ce une technique pour m'encourager à acheter ? Je ne peux pas trancher avec certitude, car je n'ai accès qu'au mouvement de ses lèvres, jamais à son intention. Cet exemple banal montre que toute connaissance d'autrui repose sur une \emph{interprétation} de signes corporels, donc sur un risque d'erreur.
\end{exemplebox}

Ce problème n'est pas un jeu d'esprit : il a des conséquences concrètes immenses. Si je peux me tromper sur ce qu'autrui pense et ressent, alors le malentendu, le conflit et le soupçon sont toujours possibles dans la vie sociale. Mais inversement, si je reconnais en autrui une conscience semblable à la mienne, alors je lui dois le respect, la parole et la justice. La question de la connaissance d'autrui débouche donc naturellement sur la question de nos \emph{rapports} avec lui.

\courssub{Problématique et annonce du plan}

\begin{methode}[Construire une problématique à partir du paradoxe d'autrui]
En dissertation, une bonne problématique sur autrui s'obtient en trois étapes :
\begin{enumerate}
\item \textbf{Dégager le paradoxe initial} : autrui est à la fois le plus proche (un semblable, un « autre moi ») et le plus inaccessible (sa conscience m'est cachée derrière son corps).
\item \textbf{Transformer le paradoxe en question} : comment puis-je connaître autrui si je n'ai accès qu'à son corps ? Et si je le reconnais comme sujet, quel type de relation m'unit à lui : conflit, domination, ou reconnaissance mutuelle ?
\item \textbf{Formuler la problématique en une phrase} qui articule les deux dimensions : connaissance et relation.
\end{enumerate}
Piège à éviter : ne jamais réduire la problématique à « qui est autrui ? », question de pure définition, sans enjeu. Il faut toujours montrer \emph{pourquoi} autrui pose problème.
\end{methode}

En suivant cette méthode, nous pouvons formuler la problématique de notre leçon :

\begin{aretenir}[Problématique de la leçon]
\textbf{Comment connaissons-nous autrui, et quelle est la nature de nos rapports avec lui ?} Autrement dit : par quels moyens puis-je accéder à la conscience d'un autre homme, dont je ne perçois que le corps ? Et une fois reconnu comme sujet, autrui est-il pour moi un obstacle qui limite ma liberté, ou la condition même de mon humanité ?
\end{aretenir}

Pour répondre à cette double question, notre leçon suivra deux grandes étapes :

\begin{itemize}
\item \textbf{Première partie : les modes de connaissance d'autrui.} Nous examinerons comment je passe du corps perçu à la conscience supposée : la perception, le raisonnement par analogie, l'empathie, le langage et la communication.
\item \textbf{Deuxième partie : l'intersubjectivité, ou la question des rapports avec autrui.} Nous étudierons la nature du lien qui unit les consciences entre elles : autrui est-il un rival qui menace ma liberté (Hegel, Sartre), ou un partenaire sans lequel je ne pourrais devenir pleinement humain ?
\end{itemize}

\begin{exoresolu}[Application : dégager le paradoxe d'une situation]
Énoncé. Un élève déclare : « Mon meilleur ami, je le connais parfaitement, je sais toujours ce qu'il pense. » Montrez que cette affirmation, bien que naturelle, contient un paradoxe philosophique, et formulez la question qu'elle soulève.
\tcblower
Solution détaillée. L'affirmation repose sur une expérience réelle : à force de fréquenter son ami, l'élève a appris à interpréter ses gestes, ses silences, ses expressions. Il y a donc une part de vérité : la connaissance d'autrui par les signes peut devenir très fine. Mais le paradoxe demeure : même dans l'amitié la plus profonde, l'élève n'a jamais accès directement à la conscience de son ami ; il ne fait que \emph{conjecturer} à partir du corps. Son ami peut toujours lui cacher une pensée, lui mentir ou le surprendre. Dire « je sais toujours ce qu'il pense » revient donc à nier la distance irréductible entre deux consciences. La question philosophique soulevée est alors : la connaissance d'autrui peut-elle jamais atteindre la certitude, ou reste-t-elle toujours une interprétation incertaine de signes ? C'est précisément le problème que la première partie de la leçon devra examiner.
\end{exoresolu}

\begin{exercice}[Pour s'entraîner]
À partir d'une scène de votre vie quotidienne (au quartier, au marché, en classe), décrivez en quelques lignes une situation où vous avez cru connaître ce que pensait ou ressentait quelqu'un. Identifiez : (1) les signes corporels perçus ; (2) l'état de conscience supposé ; (3) le risque d'erreur possible. Concluez en montrant en quoi cette situation illustre le paradoxe d'autrui.
\end{exercice}

\begin{corrige}[Exercice « Pour s'entraîner »]
Exemple de réponse attendue. Situation : en classe, je vois mon camarade baisser la tête et serrer son stylo après la distribution des notes. (1) Signes perçus : tête baissée, mâchoire contractée, silence, poing serré sur le stylo. (2) État de conscience supposé : je conclus qu'il est déçu ou honteux de sa note. (3) Risque d'erreur : il peut aussi être simplement fatigué, malade, ou préoccupé par un problème familial sans rapport avec la note ; il peut même jouer la déception pour qu'on le plaigne. Conclusion : cette scène illustre le paradoxe d'autrui, car je passe d'un corps perçu (signes extérieurs) à une conscience supposée (déception) sans jamais avoir d'accès direct à ce qu'il vit intérieurement. Ma connaissance d'autrui est donc une interprétation de signes, toujours révisable, jamais une certitude absolue.
\end{corrige}

\coursec{La connaissance d'autrui par la perception et l'analogie}

Nous avons vu dans l'introduction qu'autrui se présente à nous comme une énigme : il est à la fois un corps visible dans le monde et une conscience invisible, un «~je~» qui n'est pas le mien. La problématique de la leçon est donc posée~: comment puis-je connaître autrui, alors que je n'ai jamais accès directement à sa conscience~? Nous allons maintenant examiner le premier mode de connaissance d'autrui, le plus spontané et le plus immédiat~: la \cle{perception} de son corps, prolongée par un raisonnement particulier, le raisonnement par \cle{analogie}. Notre hypothèse de départ est celle du sens commun~: je connais autrui parce que je le vois. Il faudra tester cette hypothèse, en mesurer la force, mais aussi les limites.

\courssub{La perception directe du corps d'autrui}

Commençons par l'expérience la plus simple. Je marche dans une rue de Yaoundé et je croise un passant. Que perçois-je exactement~? Je vois un corps qui se déplace, un visage tourné vers moi, des gestes, une démarche, une posture. J'entends une voix, des paroles. Tout cela est donné à mes sens~: c'est le domaine de la \cle{perception}.

\begin{definition}[La perception]
La perception est l'acte par lequel un sujet prend conscience du monde extérieur par l'intermédiaire de ses sens (vue, ouïe, toucher, etc.). Percevoir autrui, c'est d'abord percevoir son \cle{corps}~: son visage, ses gestes, ses attitudes, ses paroles.
\end{definition}

L'intuition première est donc la suivante~: autrui m'est donné \emph{corporellement}. Je ne rencontre jamais une «~conscience pure~» flottant dans l'air~; je rencontre toujours un corps animé, expressif, qui sourit, qui fronce les sourcils, qui serre les poings ou qui ouvre les bras. Le corps d'autrui n'est pas un corps comme les autres~: contrairement à une table ou à un arbre, il est \emph{expressif}, c'est-à-dire qu'il semble porter à sa surface des significations intérieures. Un visage rougi, des larmes, un rire~: autant de \cle{signes} que je lis spontanément comme l'expression d'une vie psychique.

\begin{exemplebox}[Lecture spontanée des émotions]
Au marché Mokolo à Yaoundé, je vois une vendeuse qui sourit en accueillant une cliente~: je conclus immédiatement qu'elle est contente. Un chauffeur de taxi-moto klaxonne et crie après un automobiliste~: je conclus qu'il est en colère. Un enfant pleure devant son école~: je conclus qu'il est triste ou qu'il a mal. Dans chaque cas, je passe instantanément d'un signe corporel perçu (sourire, cri, larme) à un état de conscience supposé (joie, colère, tristesse).
\end{exemplebox}

\begin{popfigure}
\begin{center}
\begin{tikzpicture}[scale=1.05]
% Visage 1 : joie
\draw[thick,popblue] (0,0) circle (1);
\fill[popdark] (-0.35,0.25) circle (0.07);
\fill[popdark] (0.35,0.25) circle (0.07);
\draw[thick,popdark] (-0.4,-0.35) .. controls (-0.15,-0.7) and (0.15,-0.7) .. (0.4,-0.35);
\draw[thick,popblue] (-0.5,0.55) .. controls (-0.35,0.7) .. (-0.2,0.55);
\draw[thick,popblue] (0.2,0.55) .. controls (0.35,0.7) .. (0.5,0.55);
\node[below,popdark] at (0,-1.15) {\small Joie};
% Visage 2 : colère
\draw[thick,poporange] (3.2,0) circle (1);
\fill[popdark] (2.85,0.15) circle (0.07);
\fill[popdark] (3.55,0.15) circle (0.07);
\draw[thick,popdark] (2.7,0.6) -- (3.05,0.42);
\draw[thick,popdark] (3.7,0.6) -- (3.35,0.42);
\draw[thick,popdark] (2.8,-0.5) .. controls (3.05,-0.25) and (3.35,-0.25) .. (3.6,-0.5);
\node[below,popdark] at (3.2,-1.15) {\small Colère};
% Visage 3 : tristesse
\draw[thick,poppurple] (6.4,0) circle (1);
\fill[popdark] (6.05,0.2) circle (0.07);
\fill[popdark] (6.75,0.2) circle (0.07);
\draw[thick,poppurple] (5.9,0.5) .. controls (6.05,0.62) .. (6.2,0.5);
\draw[thick,poppurple] (6.6,0.5) .. controls (6.75,0.62) .. (6.9,0.5);
\draw[thick,popdark] (6.0,-0.4) .. controls (6.25,-0.15) and (6.55,-0.15) .. (6.8,-0.4);
\draw[thick,popblue] (6.78,-0.05) -- (6.78,-0.35);
\fill[popblue] (6.78,-0.42) circle (0.06);
\node[below,popdark] at (6.4,-1.15) {\small Tristesse};
% flèches de lecture
\draw[-{Stealth},very thick,popgreen] (7.6,0) -- (8.6,0);
\node[align=left,popdark] at (10.4,0) {\small Lecture~: je devine\\ \small l'émotion intérieure};
\end{tikzpicture}
\end{center}
\legende{«~Photographie~» schématique commentée~: trois expressions du visage (joie, colère, tristesse). La bouche, les sourcils et les larmes sont des signes extérieurs que je lis spontanément comme l'expression d'un état de conscience. Mais que vois-je réellement~? Des traits, des courbes, des mouvements~: jamais l'émotion elle-même.}
\end{popfigure}

\courssub{La limite de la perception~: des signes, jamais la conscience}

Arrêtons-nous maintenant sur une observation décisive. Lorsque je vois mon voisin pleurer, que perçois-je \emph{au sens strict}~? Des larmes qui coulent, des traits contractés, des sanglots. Mais sa tristesse elle-même, ce qu'il \emph{vit} de l'intérieur, je ne la perçois pas. Je perçois des signes extérieurs, jamais la conscience elle-même.

\begin{attention}[Ce que la perception ne donne pas]
La perception ne me donne jamais accès à la conscience d'autrui. Je vois le corps, j'entends la voix, je touche la main~; mais la douleur, la joie, la pensée d'autrui restent invisibles en elles-mêmes. Confondre le signe perçu (la larme) avec la réalité vécue (la tristesse) est l'erreur la plus fréquente~: le signe n'est pas la chose signifiée.
\end{attention}

Cette limite est irréductible. Même avec les instruments les plus perfectionnés, un médecin à l'hôpital de Douala peut mesurer les battements du cœur, la tension, l'activité cérébrale d'un patient~: il mesure toujours des phénomènes corporels, jamais la douleur telle que le patient la ressent. La douleur vécue est \emph{privée}~: elle n'existe que pour celui qui l'éprouve. C'est ce qu'on appelle le problème de l'accès à la conscience d'autrui~: ma conscience m'est donnée de l'intérieur, celle d'autrui ne m'est donnée que de l'extérieur, par l'intermédiaire de son corps.

Il faut donc conclure que la perception seule est insuffisante. Si je m'en tenais à ce que je perçois rigoureusement, autrui ne serait qu'un corps animé parmi d'autres, une sorte de machine compliquée. Or je suis convaincu qu'il y a «~quelqu'un~» derrière ce visage. D'où me vient cette conviction~? C'est ici qu'intervient le second moment de notre analyse~: le raisonnement par analogie.

\courssub{Le raisonnement par analogie}

L'intuition est la suivante~: autrui me ressemble. Son corps est semblable au mien~: deux yeux, deux mains, une bouche, des gestes comparables. Or je sais, par expérience intérieure, que \emph{mon} corps est animé par une conscience~: quand je me blesse, je souffre~; quand je souris, c'est que je suis content. Je suppose donc que le corps semblable d'autrui est, lui aussi, animé par une conscience semblable.

\begin{definition}[L'analogie]
L'\cle{analogie} est un raisonnement qui conclut d'une ressemblance constatée (les corps) à une ressemblance supposée (les consciences). Parce que le corps d'autrui ressemble au mien et se comporte comme le mien dans des situations semblables, je juge probable qu'il est, comme le mien, le siège d'une conscience.
\end{definition}

Analysons rigoureusement la structure de ce raisonnement. Il comporte trois étapes~:

\begin{enumerate}
\item \textbf{Constat de ressemblance corporelle}~: le corps d'autrui est semblable au mien (forme, gestes, réactions).
\item \textbf{Expérience de mon propre cas}~: je sais que mon corps est uni à une conscience~; à tel comportement de mon corps (je pleure) correspond tel état de ma conscience (je suis triste).
\item \textbf{Inférence}~: puisque le corps d'autrui se comporte comme le mien, je conclus qu'il est probablement accompagné d'une conscience qui vit des états semblables.
\end{enumerate}

\begin{popfigure}
\begin{center}
\begin{tikzpicture}[node distance=0.5cm and 0.9cm]
\node[draw=popblue,very thick,fill=popblueL,rounded corners,align=center,minimum width=3.4cm,minimum height=1.1cm] (a) {\textbf{Corps d'autrui perçu}\\ \small visage, gestes, attitudes};
\node[draw=poporange,very thick,fill=poporangeL,rounded corners,align=center,minimum width=3.4cm,minimum height=1.1cm,right=of a] (b) {\textbf{Ressemblance}\\ \small avec mon propre corps};
\node[draw=poppurple,very thick,fill=poppurpleL,rounded corners,align=center,minimum width=3.4cm,minimum height=1.1cm,right=of b] (c) {\textbf{Inférence analogique}\\ \small «~comme moi, donc...~»};
\node[draw=popgreen,very thick,fill=popgreenL,rounded corners,align=center,minimum width=3.4cm,minimum height=1.1cm,right=of c] (d) {\textbf{Conscience supposée}\\ \small d'autrui (probable)};
\draw[-{Stealth},very thick,popdark] (a) -- (b);
\draw[-{Stealth},very thick,popdark] (b) -- (c);
\draw[-{Stealth},very thick,popdark] (c) -- (d);
\node[below=0.35cm of b,align=center,popmuted] {\small Ce que je constate};
\node[below=0.35cm of d,align=center,popmuted] {\small Ce que je conclus};
\end{tikzpicture}
\end{center}
\legende{Schéma en chaîne du raisonnement par analogie~: je pars d'un corps perçu (1), je constate sa ressemblance avec le mien (2), j'opère une inférence (3) et je conclus à l'existence probable d'une conscience semblable à la mienne (4). La conclusion porte la marque de la probabilité, non de la certitude.}
\end{popfigure}

\begin{exemplebox}[Exemple camerounais~: la douleur du voisin]
Dans un quartier de Bafoussam, mon voisin se blesse au pied en fendant du bois~; il crie, se tient la jambe, grimace. Je sais immédiatement qu'il souffre, et je cours chercher de l'aide. Comment le sais-je~? Je n'ai jamais senti \emph{sa} douleur. Mais je sais ce que la douleur est \emph{pour moi}~: quand je me blesse, je crie et je grimace de la même manière. Par analogie, je conclus que derrière ses cris il y a une souffrance vécue, semblable à celle que je connais de l'intérieur. Toute la solidarité de voisinage, si vivante dans nos communautés camerounaises, repose sur ce raisonnement silencieux et instantané.
\end{exemplebox}

\courssub{Critique de l'analogie~: une probabilité, jamais une certitude}

Ce raisonnement est-il une preuve~? Examinons-le avec la rigueur du philosophe. Une démonstration certaine est un raisonnement dont la conclusion s'impose nécessairement (comme en mathématiques). Or l'analogie est une \cle{inférence}~: elle passe du connu à l'inconnu, du visible à l'invisible. Une telle inférence ne produit jamais qu'une conclusion \emph{probable}.

\begin{attention}[L'analogie ne prouve pas]
L'analogie ne \emph{prouve} pas l'existence de la conscience d'autrui~; elle la rend seulement hautement probable. Dans une dissertation, il faut donc éviter d'écrire «~l'analogie démontre qu'autrui a une conscience~»~: il faut écrire «~l'analogie fonde une présomption raisonnable, mais révisable~».
\end{attention}

Deux objections classiques montrent cette faiblesse.

\textbf{Première objection~: le problème de la simulation.} Si je ne perçois que des signes extérieurs, rien ne garantit que ces signes expriment une émotion réelle. Un signe peut être produit volontairement, sans la conscience correspondante.

\begin{exemplebox}[L'acteur de théâtre]
Un acteur de théâtre à Douala joue une scène de deuil~: il pleure, sa voix tremble, son visage exprime une tristesse bouleversante. Le spectateur qui s'en tiendrait au seul raisonnement analogique conclurait~: «~il est triste~». Or l'acteur ne ressent peut-être rien~: il \emph{simule}. L'analogie seule ne permet pas de distinguer l'émotion vraie de l'émotion jouée, puisque les signes extérieurs sont identiques dans les deux cas.
\end{exemplebox}

\textbf{Seconde objection~: le problème du mensonge.} Autrui peut non seulement simuler, mais aussi mentir~: me montrer un visage amical tout en me haïssant, affirmer qu'il va bien alors qu'il souffre. Le mot même de «~personne~» vient du latin \emph{persona}, le masque de théâtre~: autrui peut toujours se présenter à moi sous un masque. L'analogie, qui raisonne à partir des apparences, est impuissante à percer le masque.

\begin{exoresolu}[Analyser la portée du raisonnement analogique]
\textbf{Énoncé.} Un élève affirme~: «~Je vois mon camarade rire, donc je sais avec certitude qu'il est heureux.~» Ce raisonnement est-il valide~? Analysez sa structure et dégagez sa limite.
\tcblower
\textbf{Solution détaillée.} 1) Ce raisonnement est une inférence analogique~: il part d'un signe perçu (le rire) et conclut à un état de conscience (le bonheur), en s'appuyant sur la ressemblance entre le corps du camarade et le mien («~quand je ris, c'est que je suis heureux~»). 2) Or la conclusion ne suit pas avec nécessité~: le camarade peut rire par politesse, par nervosité, ou pour cacher une peine~; il peut simuler. 3) La limite est donc celle de toute analogie~: elle établit une probabilité, non une certitude. 4) Formulation correcte~: «~Je vois mon camarade rire, donc il est très probable qu'il soit heureux, mais je ne peux pas en avoir la certitude absolue.~»
\end{exoresolu}

\courssub{La thèse de Husserl~: autrui est donné par «~apprésentation~»}

Le philosophe allemand Edmund \cle{Husserl} (1859-1938), fondateur de la phénoménologie, a approfondi cette analyse dans ses \emph{Méditations cartésiennes} (1931). Sa question est précise~: comment autrui peut-il m'être donné comme \emph{autre conscience}, alors que je ne peux jamais vivre sa vie intérieure~?

\begin{aretenir}[Husserl, \emph{Méditations cartésiennes}]
Selon Husserl, autrui m'est donné par \cle{apprésentation}~: son corps, présent devant moi, m'\emph{indique} une conscience qui, elle, n'est jamais présentée directement. Quand je vois une maison, je ne vois que sa façade, mais j'anticipe qu'elle a un arrière~: l'arrière est «~apprésenté~» avec la façade. De même, le corps d'autrui est présenté, et sa conscience est apprésentée avec lui — sauf que, contrairement à l'arrière de la maison, je ne pourrai \emph{jamais} faire le tour pour vérifier~: la conscience d'autrui reste à jamais inaccessible à mon expérience directe.
\end{aretenir}

La thèse de Husserl confirme et précise ce que nous avons découvert~: la conscience d'autrui n'est pas une invention de mon imagination, car elle m'est réellement \emph{donnée} avec son corps (autrui n'est pas un fantôme)~; mais elle ne m'est jamais donnée \emph{en personne}. Autrui est là, devant moi, et pourtant son intériorité me demeure voilée. C'est cette structure paradoxale — présence du corps, absence de la conscience — qui fait toute l'originalité de l'expérience d'autrui.

\begin{savaistu}[De la phénoménologie à la vie quotidienne]
L'analyse de Husserl éclaire des situations très concrètes. Lors d'une consultation à l'hôpital général de Yaoundé, le médecin demande au patient~: «~Où avez-vous mal~? Sur une échelle de 0 à 10, quelle est l'intensité de votre douleur~?~» Pourquoi cette question~? Parce que la douleur du patient est apprésentée, jamais présentée~: le médecin ne peut pas la mesurer directement comme il mesure la température. Il doit faire confiance à la parole du patient, c'est-à-dire s'en remettre à un témoignage sur une conscience inaccessible. La médecine elle-même bute sur le mystère d'autrui.
\end{savaistu}

\courssub{Conséquence~: une connaissance toujours indirecte et perfectible}

Tirons le bilan de cette section. La connaissance d'autrui par la perception et l'analogie possède trois caractères essentiels.

\begin{enumerate}
\item Elle est \textbf{indirecte}~: je n'accède à la conscience d'autrui que par l'intermédiaire de son corps, de ses gestes, de ses paroles — autant de médiateurs.
\item Elle est \textbf{probable}~: fondée sur l'analogie, elle ne produit jamais de certitude absolue~; la simulation et le mensonge restent toujours possibles.
\item Elle est \textbf{perfectible}~: parce qu'elle est incertaine, elle peut être corrigée, approfondie, affinée. La fréquentation prolongée, l'écoute, le dialogue permettent de mieux connaître autrui, sans jamais atteindre une transparence totale.
\end{enumerate}

\begin{methode}[Mobiliser cette analyse dans une dissertation]
Pour utiliser correctement cette section dans un devoir~:
\begin{enumerate}
\item Commencer par décrire l'expérience spontanée~: je perçois le corps expressif d'autrui (visage, gestes, attitudes).
\item Dégager la limite~: je ne perçois que des signes, jamais la conscience elle-même.
\item Expliquer le recours à l'analogie~: ressemblance des corps, inférence vers une conscience semblable (donner un exemple concret).
\item Critiquer~: l'analogie donne une probabilité, non une certitude (simulation, mensonge).
\item Conclure provisoirement avec Husserl~: autrui est apprésenté à travers son corps~; sa conscience reste inaccessible directement.
\end{enumerate}
\end{methode}

\begin{exercice}[Entraînement]
1) Définissez le raisonnement par analogie et décrivez ses trois étapes. 2) Pourquoi dit-on que la conscience d'autrui est «~privée~»~? 3) Expliquez, à l'aide de l'exemple de l'acteur de théâtre, pourquoi l'analogie ne produit pas de certitude. 4) Que signifie le terme «~apprésentation~» chez Husserl~? En quoi ce concept dépasse-t-il le simple raisonnement par analogie~?
\end{exercice}

\begin{corrige}[Exercice]
1) L'analogie est un raisonnement qui conclut d'une ressemblance constatée (les corps) à une ressemblance supposée (les consciences). Étapes~: constat de la ressemblance corporelle~; expérience de mon propre cas (mon corps est uni à ma conscience)~; inférence vers une conscience semblable chez autrui. 2) La conscience d'autrui est «~privée~» parce qu'elle n'existe que pour celui qui la vit~: je ne peux ni voir ni mesurer la douleur, la joie ou la pensée d'autrui de l'intérieur~; je n'en perçois que les signes extérieurs. 3) L'acteur produit les mêmes signes extérieurs (larmes, voix tremblante) que l'homme réellement triste. Puisque les signes sont identiques, l'analogie, qui raisonne à partir des seuls signes, ne peut distinguer l'émotion vraie de l'émotion simulée~: sa conclusion reste donc seulement probable. 4) Chez Husserl, l'apprésentation désigne le mode de donation d'autrui~: son corps est présenté directement, sa conscience est co-donnée avec lui mais sans jamais être présentée en personne. Ce concept dépasse l'analogie car il montre qu'autrui n'est pas le résultat d'un raisonnement que je ferais après coup~: il m'est donné d'emblée comme conscience autre, quoique inaccessible, dans l'expérience perceptive elle-même.
\end{corrige}

Ainsi, la perception et l'analogie me donnent bien une connaissance d'autrui, mais une connaissance fragile, indirecte et toujours révisable. Faut-il alors renoncer à connaître autrui~? Non~: il existe d'autres voies, plus riches, que le seul regard porté sur un corps. La sympathie, l'écoute de la parole, l'échange dans le dialogue ouvrent un accès d'un autre type à la vie d'autrui. C'est ce que nous examinerons dans la section suivante, consacrée à l'empathie, au langage et à la communication.

\coursec{Empathie, langage et communication}

Dans la section précédente, nous avons montr\'e que la perception et l'analogie constituent les premiers modes de connaissance d'autrui~: je vois le corps de l'autre, et j'inf\`ere, par ressemblance avec le mien, qu'il abrite une conscience semblable \`a la mienne. Mais cette inf\`erence reste fragile~: elle me donne autrui comme un \emph{objet suppos\'e pensant}, non comme un sujet dont je partagerais l'int\'eriorit\'e. Or, dans la vie quotidienne, je ne me contente pas de deviner autrui~: je \emph{ressens avec lui}, et il me \emph{parle}. Deux modes de connaissance plus profonds doivent donc \^etre examin\'es~: l'\cle{empathie} et le \cle{langage}. Notre probl\`eme est le suivant~: ces deux modes me donnent-ils enfin un acc\`es direct \`a la conscience d'autrui, ou bien autrui demeure-t-il, m\^eme quand il me parle, partiellement ferm\'e \`a moi~?

\courssub{L'empathie~: ressentir avec autrui}

\textbf{Intuition.} Un ami perd son p\`ere. Avant m\^eme qu'il prononce un mot, en voyant son visage ferm\'e et ses \'epaules affaiss\'ees, je sens en moi-m\^eme une tristesse qui n'est pas la mienne et qui pourtant m'envahit. Je ne devine pas sa douleur comme je devine le contenu d'une bo\^ite ferm\'ee~: je l'\emph{\'eprouve} avec lui, \`a distance. Cette exp\'erience universelle est le point de d\'epart de la notion d'empathie.

\begin{definition}[L'empathie]
L'\cle{empathie} (en allemand \emph{Einf\"uhlung}, litt\'eralement «~sentiment \`a l'int\'erieur de~») est la capacit\'e de comprendre les \'etats affectifs d'autrui en se projetant en lui, c'est-\`a-dire en se mettant \`a sa place, \textbf{sans pour autant se confondre avec lui}. Dans l'empathie, je ressens la tristesse de l'autre \emph{comme sienne}~: je sais que c'est lui qui souffre, non moi.
\end{definition}

La notion a \'et\'e th\'ematis\'ee par la philosophe Edith Stein (disciple de Husserl) dans sa th\`ese \emph{Le probl\`eme de l'empathie} (1917)~: l'empathie est l'acte par lequel la conscience d'autrui m'est donn\'ee non pas directement (je ne suis jamais dans la t\^ete de l'autre), mais de mani\`ere \emph{analogique et v\'ecue} --- je vis l'\'emotion d'autrui «~\'a travers~» la mienne, sans la confondre avec la mienne.

\begin{attention}[Ne pas confondre empathie et sympathie]
Erreur fr\'equente dans les copies~: assimiler empathie et sympathie. Or~:
\begin{itemize}
\item la \cle{sympathie} (du grec \emph{sun-pathos}, «~souffrir avec~» au sens d'affinit\'e) d\'esigne l'\textbf{attirance affective}~: j'\emph{aime} autrui, je suis bien dispos\'e \`a son \'egard~;
\item l'\cle{empathie} d\'esigne une \textbf{op\'eration de connaissance}~: je \emph{comprends} autrui de l'int\'erieur, que je l'aime ou non.
\end{itemize}
Preuve qu'elles sont distinctes~: un bon m\'edecin ou un enqu\^eteur doit faire preuve d'empathie envers quelqu'un qu'il n'aime pas~; inversement, je peux aimer quelqu'un sans comprendre ce qu'il vit. L'empathie est un \emph{mode de connaissance}, la sympathie est un \emph{sentiment}.
\end{attention}

\begin{exemplebox}[La solidarit\'e villageoise lors des fun\'erailles au Cameroun]
Dans de nombreuses communaut\'es camerounaises, les fun\'erailles rassemblent tout le village ou tout le quartier~: voisins, parent\'e \'eloign\'ee, associations (\emph{tontines}, \emph{njangi}). Les participants pleurent avec la famille, chantent, veillent le corps, partagent les repas du deuil. Ce n'est pas seulement un devoir social~: chacun \emph{se met \`a la place} des endeuill\'es --- «~demain, ce pourrait \^etre moi~» --- et ressent avec eux la perte. Cette participation affective partag\'ee est une forme collective d'empathie~: elle permet \`a chacun de \emph{comprendre de l'int\'erieur} la douleur de l'autre, et elle all\`ege cette douleur en la r\'epartissant. L'empathie appara\^it ainsi comme un ciment du vivre-ensemble camerounais.
\end{exemplebox}

\begin{savaistu}[Les neurones miroirs]
La psychologie et les neurosciences contemporaines apportent un compl\'ement biologique \`a la notion philosophique d'empathie. Dans les ann\'ees 1990, l'\'quipe du neurophysiologiste italien Giacomo Rizzolatti a d\'ecouvert chez le singe (puis chez l'homme) des neurones dits «~\cle{miroirs}~»~: ces neurones s'activent \`a la fois quand le sujet accomplit une action ou \'eprouve une \'emotion \emph{et} quand il voit un cong\'en\`ere accomplir cette action ou \'eprouver cette \'emotion. Voir la grimace de douleur d'autrui active donc en moi les m\^emes circuits que ma propre douleur. Cette d\'ecouverte sugg\`ere que l'empathie a un fondement naturel~: nous serions biologiquement «~c\^abl\'es~» pour ressentir avec autrui. (Compl\'ement utile pour enrichir une copie du Bac, sans remplacer l'analyse philosophique.)
\end{savaistu}

\courssub{Le langage~: autrui me dit ce qu'il pense}

\textbf{Intuition.} Perception, analogie et empathie me font \emph{deviner} ou \emph{ressentir} autrui de l'ext\'erieur. Mais autrui poss\`ede un pouvoir extraordinaire~: il peut \emph{parler}. Quand mon voisin me dit «~j'ai peur~», je n'ai plus besoin d'inf\'erer sa peur~: il me la communique lui-m\^eme. Le langage semble \^etre l'acc\`es privil\'egi\'e, voire direct, \`a la conscience d'autrui.

\begin{definition}[Le langage comme mode de connaissance d'autrui]
Le \cle{langage} est le syst\`eme de signes (mots, phrases) par lequel autrui exprime volontairement ses pens\'ees, ses d\'esirs et ses sentiments. Par la parole, autrui cesse d'\^etre un objet que j'observe pour devenir un \textbf{sujet qui se r\'ev\`ele lui-m\^eme}~: il me dit ce qu'il pense et ce qu'il ressent. C'est le seul mode de connaissance o\`u autrui participe activement \`a sa propre connaissance par moi.
\end{definition}

La sup\'eriorit\'e du langage sur les autres modes est \'evidente~: la perception ne m'atteint que le corps, l'analogie ne me donne qu'une hypoth\`ese, l'empathie qu'une \'emotion partag\'ee mais muette. La parole, elle, peut exprimer des pens\'ees abstraites, des souvenirs, des projets --- des contenus invisibles et inaccessibles \`a tout regard. Mais cette sup\'eriorit\'e est-elle totale~? Non, et c'est l\`a le point d\'ecisif pour la dissertation.

\textbf{Les limites du langage.} Trois obstacles emp\^echent la parole d'\^etre une fen\^etre transparente sur la conscience~:
\begin{enumerate}
\item \textbf{Le mensonge}~: la parole peut vouloir dire le contraire de ce qu'elle dit. Autrui peut feindre, dissimuler, manipuler. Le menteur utilise pr\'ecis\'ement la confiance que nous accordons au langage pour nous tromper.
\item \textbf{Le malentendu}~: m\^eme de bonne foi, autrui peut mal choisir ses mots, et moi mal les interpr\'eter. Les mots n'ont pas exactement le m\^eme sens pour tous~: le mot «~respect~» n'\'evoque pas la m\^eme chose pour un \'el\`eve de Premi\`ere et pour un notable traditionnel.
\item \textbf{L'indicible}~: certaines exp\'eriences r\'esistent \`a toute expression. La douleur profonde, le deuil, l'amour intense «~se disent mal~»~: celui qui souffre le plus est souvent celui qui se tait. Comme le note le philosophe Vladimir Jank\'el\'evitch, il y a des v\'erit\'es de l'int\'eriorit\'e que la parole trahit en les formulant.
\end{enumerate}

\begin{attention}[Pi\`ege classique de la dissertation]
Ne jamais \'ecrire que «~le langage donne un acc\`es transparent \`a la pens\'ee d'autrui~». C'est l'erreur la plus fr\'equente. La parole est un \emph{m\'edium}~: elle passe entre la conscience d'autrui et moi, et comme tout interm\'ediaire, elle peut d\'eformer, filtrer ou trahir. Dire n'est pas montrer~: entre ce qu'autrui pense et ce qu'il dit, il y a toujours un \'ecart possible.
\end{attention}

\begin{exemplebox}[Le proverbe camerounais~: «~la bouche dit ce que le c\oe ur ne pense pas~»]
La sagesse populaire camerounaise conna\^it parfaitement cette limite du langage. Le proverbe «~\emph{la bouche dit ce que le c\oe ur ne pense pas}~» signifie que les paroles peuvent \^etre en d\'ecalage complet avec les intentions r\'eelles~: flatterie int\'eress\'ee, promesse \'electorale non tenue, d\'eclaration d'amiti\'e hypocrite. Le proverbe invite \`a la prudence herméneutique~: pour conna\^itre autrui, il ne suffit pas d'\'ecouter ce qu'il dit, il faut croiser la parole avec les actes, le regard et le contexte. La philosophie rejoint ici la sagesse du terroir~: la parole est un indice pr\'ecieux, mais un indice \emph{faillible}.
\end{exemplebox}

\courssub{La communication non verbale~: gestes, silence et regard}

Le langage ne se r\'eduit pas aux mots. Une grande partie de ce qu'autrui me communique passe par le \cle{langage non verbal}~: gestes, postures, mimiques, ton de la voix, regard, et m\^eme le silence. Ce langage du corps est souvent plus difficile \`a truquer que la parole~: le menteur contr\^ole ses mots, moins facilement ses tremblements ou son regard fuyant.

Dans les cultures camerounaises, la communication non verbale poss\`ede une richesse et des codes pr\'ecis~:
\begin{itemize}
\item \textbf{Le regard}~: dans plusieurs traditions, un cadet ne soutient pas longtemps le regard d'un a\^in\'e --- signe de respect, non de dissimulation. Un observateur ignorant ce code interpr\'etera \`a tort ce regard baiss\'e comme de la g\^ene ou de la culpabilit\'e.
\item \textbf{Le silence}~: il peut signifier le consentement, la r\'eflexion, la d\'esapprobation polie ou le deuil. Dans une palabre, le silence de l'a\^in\'e p\`ese souvent plus que les discours.
\item \textbf{Les gestes}~: donner ou recevoir un objet des deux mains, ou de la main droite, marque le respect~; hocher la t\^ete, taper dans les mains, prosterner certains saluts sont des messages codifi\'es.
\end{itemize}

Mais cette richesse est aussi une \textbf{ambigu\"it\'e}~: les signes non verbaux n'ont de sens que dans une culture donn\'ee, et ils se pr\^etent \`a de multiples interpr\'etations. Un m\^eme silence peut vouloir dire «~oui~», «~non~» ou «~je r\'efl\'echis~». Le langage non verbal compl\`ete donc la parole sans \^etre plus fiable qu'elle~: il exige, lui aussi, un travail d'interpr\'etation.

\courssub{Bilan~: une connaissance toujours combin\'ee et jamais totale}

Aucun des quatre modes \'etudi\'es ne donne \`a lui seul un acc\`es total \`a autrui. La perception atteint le corps mais pas la conscience~; l'analogie suppose la conscience sans la garantir~; l'empathie fait ressentir l'\'emotion sans livrer la pens\'ee~; le langage exprime la pens\'ee mais peut tromper. La connaissance d'autrui est donc toujours une \textbf{combinaison}~: je regarde autrui, je l'\'ecoute, je ressens avec lui, j'interpr\`ete ses gestes, et je croise ces indices pour former un jugement toujours r\'evisable. Autrui n'est jamais pour moi une \'evidence transparente~: il est un \emph{myst\`ere partiellement d\'echiffrable}.

\begin{popfigure}
\begin{center}
\begin{tikzpicture}
\node[circle, draw=popblue, fill=popblueL, thick, align=center, font=\bfseries, minimum size=2.6cm] (C) at (0,0) {Connaissance\\ d'autrui};
\node[rectangle, rounded corners, draw=poporange, fill=poporangeL, thick, align=center, minimum width=3cm] (P) at (90:3.4) {Perception\\ (je vois son corps)};
\node[rectangle, rounded corners, draw=popgreen, fill=popgreenL, thick, align=center, minimum width=3cm] (A) at (0:4.2) {Analogie\\ (il me ressemble,\\ donc il pense)};
\node[rectangle, rounded corners, draw=poppurple, fill=poppurpleL, thick, align=center, minimum width=3cm] (E) at (270:3.4) {Empathie\\ (je ressens avec lui)};
\node[rectangle, rounded corners, draw=poppink, fill=poppinkL, thick, align=center, minimum width=3cm] (L) at (180:4.2) {Langage\\ (il me dit ce\\ qu'il pense)};
\draw[-{Stealth}, very thick, poporange] (C) -- (P);
\draw[-{Stealth}, very thick, popgreen] (C) -- (A);
\draw[-{Stealth}, very thick, poppurple] (C) -- (E);
\draw[-{Stealth}, very thick, poppink] (C) -- (L);
\end{tikzpicture}
\end{center}
\legende{Les quatre modes de connaissance d'autrui~: aucun ne suffit seul, la connaissance d'autrui r\'esulte de leur combinaison.}
\end{popfigure}

\begin{center}
\begin{tabularx}{\linewidth}{|l|X|X|}
\hline
\rowcolor{popblueL} \textbf{Mode} & \textbf{Avantage} & \textbf{Limite} \\
\hline
Perception & Acc\`es imm\'ediat au corps, aux expressions, aux comportements & N'atteint que l'ext\'erieur~: la conscience reste invisible \\
\hline
Analogie & Permet de supposer une conscience semblable \`a la mienne & Simple inf\'erence~: rien ne garantit que l'autre pense comme moi \\
\hline
Empathie & Fait ressentir de l'int\'erieur l'\'emotion d'autrui & Reste muette sur les pens\'ees~; risque de projection (je pr\^ete \`a autrui mes propres sentiments) \\
\hline
Langage & Autrui exprime lui-m\^eme pens\'ees, souvenirs, projets & Mensonge, malentendu, indicible~: la parole peut tromper \\
\hline
\end{tabularx}
\end{center}

\begin{aretenir}[L'essentiel de la section]
\begin{itemize}
\item L'\cle{empathie} (\emph{Einf\"uhlung})~: comprendre les \'etats affectifs d'autrui en se projetant en lui, sans se confondre avec lui. \`A distinguer de la sympathie (aimer autrui).
\item Le \cle{langage}~: autrui me dit ce qu'il pense~; c'est l'acc\`es privil\'egi\'e \`a sa conscience, mais il a trois limites~: le mensonge, le malentendu, l'indicible.
\item La \cle{communication non verbale} (gestes, silence, regard) enrichit la connaissance d'autrui mais reste ambigu\"e et d\'epend des codes culturels.
\item Bilan~: aucun mode ne donne un acc\`es total~; conna\^itre autrui exige de combiner perception, analogie, empathie et parole.
\end{itemize}
\end{aretenir}

\begin{exoresolu}[Analyse de situation~: la parole suffit-elle \`a conna\^itre autrui~?]
\textbf{\'Enonc\'e.} Pendant la campagne pour l'\'election du d\'el\'egu\'e de classe, un candidat d\'eclare~: «~Je serai toujours \`a votre \'ecoute, je d\'efendrai vos int\'er\^ets.~» Une fois \'elu, il ne tient aucune de ses promesses. Un \'el\`eve conclut~: «~On ne peut jamais conna\^itre autrui, puisque sa parole peut tromper.~» En mobilisant les notions de la le\c con, discutez cette conclusion en trois arguments.
\tcblower
\textbf{Solution d\'etaill\'ee.}
\begin{enumerate}
\item \textbf{La conclusion contient une part de v\'erit\'e.} Le langage n'est pas un acc\`es transparent \`a la pens\'ee d'autrui~: le candidat a menti, et sa parole disait le contraire de ses intentions. Comme le rappelle le proverbe «~la bouche dit ce que le c\oe ur ne pense pas~», entre la conscience et la parole il y a toujours un \'ecart possible. L'\'el\`eve a donc raison de refuser la na\"ivet\'e qui croit tout ce qu'on lui dit.
\item \textbf{Mais la g\'en\'eralisation est abusive.} D'un cas de mensonge, l'\'el\`eve conclut que la parole trompe \emph{toujours}. Or le langage reste le mode privil\'egi\'e de connaissance d'autrui~: sans lui, aucune pens\'ee abstraite, aucun projet ne pourrait \^etre communiqu\'e. Le mensonge n'est possible que parce que la parole dit vrai le plus souvent~: on ne peut tromper que celui qui a de bonnes raisons de croire.
\item \textbf{La connaissance d'autrui combine plusieurs modes.} L'erreur des \'electeurs n'a pas \'et\'e d'\'ecouter la parole du candidat, mais de s'en tenir \`a elle. Ils auraient d\^u croiser ses d\'eclarations avec son comportement pass\'e (perception), ses gestes et son regard (communication non verbale), et l'empathie envers ses motivations r\'eelles. Conclusion~: la parole peut tromper, mais autrui n'est pas pour autant inconnaissable --- il est connaissable \emph{partiellement}, par la combinaison prudente de tous les modes de connaissance.
\end{enumerate}
\end{exoresolu}

\begin{exercice}[Pour s'entra\^iner]
\begin{enumerate}
\item D\'efinissez l'empathie et distinguez-la soigneusement de la sympathie, en illustrant chaque notion par un exemple de la vie scolaire.
\item «~Autrui est plus pr\'esent \`a moi dans son silence que dans ses discours.~» Discutez cette affirmation en vous appuyant sur l'\'etude du langage et de la communication non verbale.
\item Sujet de type Bac~: \emph{La parole me permet-elle de conna\^itre autrui~?} R\'edigez l'introduction et le plan d\'etaill\'e de la dissertation.
\end{enumerate}
\end{exercice}

\begin{corrige}[Exercice]
\textbf{\'El\'ements de correction.}
\begin{enumerate}
\item Empathie~: se projeter en autrui pour comprendre ses \'etats affectifs sans se confondre avec lui (ex.~: comprendre le stress d'un camarade avant un expos\'e parce qu'on se met \`a sa place). Sympathie~: attirance affective, aimer autrui (ex.~: appr\'ecier un camarade pour sa gentillesse). On peut \'eprouver de l'empathie pour quelqu'un qu'on n'aime pas, et aimer sans comprendre.
\item Attendu~: valeur du silence et du non verbal (le silence peut exprimer le deuil, le respect, le consentement~; le corps trahit ce que la parole cache) mais limites (le silence est ambigu, il d\'epend des codes culturels)~; conclusion nuanc\'ee~: le silence compl\`ete la parole sans la remplacer.
\item Introduction~: amorce (nous connaissons autrui d'abord par ce qu'il dit), probl\'ematique (la parole est-elle un acc\`es transparent ou un indice faillible~?), annonce du plan. Plan possible~: I. La parole, acc\`es privil\'egi\'e \`a la conscience d'autrui (autrui se r\'ev\`ele lui-m\^eme~; sup\'eriorit\'e sur la perception et l'analogie). II. Mais la parole peut tromper (mensonge, malentendu, indicible~; proverbe camerounais). III. Conna\^itre autrui exige de croiser parole, gestes, empathie et actes (connaissance combin\'ee, jamais totale).
\end{enumerate}
\end{corrige}

\coursec{L'intersubjectivité : définition et formes}

Nous avons vu, dans la section précédente, que l'empathie, le langage et la communication sont les voies par lesquelles j'accède à la conscience d'autrui : je devine ses émotions, je comprends ses paroles, j'échange avec lui. Mais ces moyens de connaissance ne sont pas de simples outils : ils révèlent quelque chose de plus profond. Ils montrent que ma conscience et celle d'autrui ne vivent pas chacune dans sa tour d'ivoire, mais qu'elles entretiennent entre elles des rapports permanents. C'est précisément ce que désigne le concept d'\cle{intersubjectivité}. Il ne s'agit plus seulement de savoir \emph{comment je connais autrui}, mais de comprendre \emph{ce qui se passe entre nous} : quelles formes prennent nos rapports réciproques, et pourquoi ces rapports sont-ils constitutifs de l'existence humaine ?

\courssub{Définition de l'intersubjectivité}

Commençons par l'intuition. Observez une cour de récréation dans un lycée de Douala ou de Yaoundé : des élèves discutent, se disputent parfois, s'entraident pour les devoirs, se saluent, se reconnaissent. Aucun d'eux ne vit isolé : chacun est pris dans un tissu de relations où sa conscience rencontre d'autres consciences. Ce tissu de relations, c'est l'intersubjectivité.

\begin{definition}[L'intersubjectivité]
L'\cle{intersubjectivité} désigne l'ensemble des \textbf{rapports réciproques} entre deux ou plusieurs \textbf{consciences} (ou sujets) qui se \textbf{reconnaissent mutuellement comme sujets}, c'est-à-dire comme des êtres capables de penser, de vouloir et d'agir librement. Le préfixe \emph{inter-} signifie « entre » : l'intersubjectivité est littéralement ce qui se passe \emph{entre les sujets}.
\end{definition}

Trois éléments de cette définition méritent d'être soulignés :

\begin{itemize}
\item La \textbf{réciprocité} : le rapport va dans les deux sens. Je reconnais autrui comme sujet, et autrui me reconnaît comme sujet. Si je traite l'autre comme un simple objet (un outil, un meuble), il n'y a plus intersubjectivité mais domination ou indifférence.
\item La \textbf{conscience} : l'intersubjectivité n'existe qu'entre des êtres doués de conscience. Ma relation à ma houe ou à ma moto n'est pas intersubjective ; ma relation à mon voisin de champ l'est.
\item La \textbf{reconnaissance mutuelle} : chaque sujet accorde à l'autre le même statut qu'il revendique pour lui-même, celui de sujet libre et pensant.
\end{itemize}

\begin{attention}[Ne pas confondre intersubjectivité et simple cohabitation]
L'intersubjectivité ne signifie pas seulement « vivre ensemble » ou « être nombreux au même endroit ». Des passagers entassés dans un bus d'une agence de voyage qui ne s'adressent jamais la parole cohabitent sans véritable intersubjectivité. Il faut qu'il y ait \textbf{échange réel} et \textbf{réciprocité entre consciences} : se parler, se reconnaître, coopérer, s'affronter même. Même le conflit est une forme d'intersubjectivité, car celui que je combats, je le reconnais comme un adversaire, donc comme un sujet.
\end{attention}

\courssub{L'intersubjectivité, structure fondamentale de l'existence humaine}

Pourquoi les philosophes accordent-ils tant d'importance à ce concept ? Parce que l'intersubjectivité n'est pas un accident de la vie humaine : elle en est la \textbf{structure fondamentale}. L'homme vit \emph{toujours} avec et parmi autrui.

Faisons une simple enquête sur nous-mêmes, comme le ferait un chercheur. Essayez de répertorier une seule journée de votre vie en supprimant tout ce qui vient d'autrui : la langue que vous parlez (apprise des autres), les vêtements que vous portez (fabriqués par d'autres), le nom que vous portez (donné par vos parents), les règles que vous suivez, les savoirs que vous apprenez à l'école, la monnaie que vous utilisez au marché... Il ne resterait presque rien. Cette expérience de pensée montre que \textbf{je ne deviens un sujet que par et avec autrui}. C'est ce qu'a souligné le philosophe allemand \cle{Husserl} : ma conscience est d'emblée ouverte aux autres consciences, et le monde lui-même est un monde \emph{commun}, partagé. De même, \cle{Merleau-Ponty} montre que dès la perception, je vis dans un monde où les autres sont déjà présents.

\begin{aretenir}[L'essentiel]
L'intersubjectivité est une \textbf{structure fondamentale de l'existence humaine} : nul ne naît, ne parle, ne pense et n'agit seul. Être un sujet, c'est toujours déjà être \emph{en relation} avec d'autres sujets. L'isolement absolu (l'enfant sauvage, le naufragé seul sur une île) est une situation limite qui confirme la règle : même Robinson conserve la langue, les outils et les habitudes reçus de la société.
\end{aretenir}

Mais si les rapports entre consciences sont permanents, ils ne prennent pas toujours la même forme. Tantôt nous coopérons, tantôt nous nous reconnaissons, tantôt nous nous affrontons, tantôt nous nous aimons. Distinguons donc les \textbf{quatre grandes formes de l'intersubjectivité}.

\courssub{Première forme : la communication et l'échange}

La forme la plus quotidienne de l'intersubjectivité est la \cle{communication} : l'échange de paroles, de signes, de biens et de services entre sujets.

\begin{itemize}
\item \textbf{La parole} : dialoguer, c'est reconnaître l'autre comme quelqu'un qui peut comprendre, répondre, argumenter. La conversation de quartier, la palabre traditionnelle sous l'arbre à palabres, la réunion de famille sont des lieux d'intersubjectivité par la parole.
\item \textbf{Le commerce} : acheter et vendre au marché de Mokolo ou de Mvog-Mbi suppose un accord entre consciences : le vendeur et l'acheteur se reconnaissent mutuellement comme partenaires libres d'accepter ou de refuser le prix.
\item \textbf{La coopération au travail} : le travail en équipe, le coup de main collectif pour construire une maison ou moissonner un champ (ce qu'on appelle parfois le \emph{travail communautaire} dans les villages) montre que l'action commune suppose l'accord des volontés.
\end{itemize}

Dans toutes ces situations, autrui n'est pas un obstacle mais un \textbf{partenaire} : l'échange enrichit chacun, car je reçois de l'autre des informations, des biens ou des forces que je n'aurais pas seul.

\courssub{Deuxième forme : la reconnaissance}

La deuxième forme est plus profonde : la \cle{reconnaissance}. Je n'attends pas seulement d'autrui des paroles ou des biens : j'attends qu'il me \textbf{reconnaisse comme sujet libre et égal}, c'est-à-dire qu'il me traite comme une personne dont la dignité est égale à la sienne.

Le philosophe allemand \cle{Hegel} a montré que la conscience de soi ne peut pas se suffire : pour me savoir sujet libre, j'ai besoin qu'une autre conscience libre me le confirme. Être reconnu, c'est exister pleinement aux yeux des autres. Inversement, être méprisé, humilié ou ignoré, c'est être nié dans son humanité. C'est pourquoi la lutte pour la reconnaissance est au cœur de la vie sociale : les revendications de dignité, d'égalité et de respect (par exemple la revendication de l'égalité entre hommes et femmes, ou le refus du mépris envers les jeunes) sont des luttes pour la reconnaissance.

\begin{exemplebox}[La reconnaissance dans la vie scolaire]
Quand un professeur félicite publiquement un élève pour son travail, il le reconnaît comme sujet capable et méritant : l'élève se sent exister, sa confiance grandit. À l'inverse, l'élève constamment humilié devant la classe se sent nié comme sujet. La reconnaissance est donc un besoin aussi fondamental que la communication.
\end{exemplebox}

\courssub{Troisième forme : le conflit}

L'intersubjectivité n'est pas toujours harmonieuse. Sa troisième forme est le \cle{conflit} : la lutte, la rivalité, la domination. Les consciences peuvent s'affronter parce qu'elles désirent la même chose (rivalité pour un bien, une place, un honneur) ou parce que chacune veut imposer sa volonté à l'autre.

Hegel décrit la \textbf{lutte pour la reconnaissance} : deux consciences qui veulent chacune être reconnues sans reconnaître l'autre s'engagent dans un combat qui aboutit au rapport de \textbf{maître et esclave} — l'une domine, l'autre est dominé. Le conflit montre que l'intersubjectivité peut se dégrader : quand la réciprocité se brise, l'un des sujets est réduit au rang d'objet ou d'instrument.

\begin{attention}[Le conflit reste une forme d'intersubjectivité]
Ne dites pas que le conflit est « l'absence de rapport avec autrui ». Au contraire : on ne se bat qu'avec quelqu'un qu'on reconnaît comme adversaire, donc comme sujet. Le conflit est une intersubjectivité \emph{dégradée}, où la réciprocité est rompue, mais il suppose toujours la rencontre des consciences. L'indifférence totale serait pire que le conflit.
\end{attention}

\courssub{Quatrième forme : l'amour et l'amitié}

La quatrième forme est la plus haute : l'\cle{amour} et l'\cle{amitié}, c'est-à-dire l'\textbf{union volontaire des consciences}. Ici, les sujets ne se contentent ni d'échanger, ni de se reconnaître, ni de s'affronter : ils veulent s'unir, partager leur vie, faire du bien à l'autre pour lui-même.

Dans l'amitié véritable, chacun veut le bien de l'autre autant que le sien propre ; dans l'amour, deux consciences libres choisissent de s'appartenir mutuellement sans se posséder comme des choses. Le mariage traditionnel camerounais, qui unit non seulement deux personnes mais deux familles, illustre bien cette union volontaire des consciences élargie à toute une communauté.

\begin{popfigure}
\begin{center}
\begin{tikzpicture}[scale=1]
% Deux sujets au centre
\draw[fill=popblueL,draw=popblue,thick] (-1.6,0) circle (0.55);
\node at (-1.6,0.15) {\small Sujet};
\node at (-1.6,-0.2) {\small A};
\draw[fill=poppinkL,draw=poppink,thick] (1.6,0) circle (0.55);
\node at (1.6,0.15) {\small Sujet};
\node at (1.6,-0.2) {\small B};
% Flèche réciproque centrale
\draw[<->,thick,popdark] (-1.0,0) -- (1.0,0);
\node[below,popdark] at (0,-0.15) {\scriptsize réciprocité};
% Quatre formes autour
\node[draw=popgreen,fill=popgreenL,rounded corners,align=center,thick] (com) at (0,2.3) {\small \textbf{Communication}\\ \scriptsize parole, commerce, coopération};
\node[draw=popgold,fill=popgoldL,rounded corners,align=center,thick] (rec) at (-4.3,0) {\small \textbf{Reconnaissance}\\ \scriptsize être reconnu libre et égal};
\node[draw=poporange,fill=poporangeL,rounded corners,align=center,thick] (conf) at (0,-2.3) {\small \textbf{Conflit}\\ \scriptsize lutte, rivalité, domination};
\node[draw=poppurple,fill=poppurpleL,rounded corners,align=center,thick] (am) at (4.3,0) {\small \textbf{Amour -- Amitié}\\ \scriptsize union volontaire des consciences};
% Liens
\draw[->,popgreen,thick] (com.south) -- (0,0.6);
\draw[->,popgold,thick] (rec.east) -- (-2.2,0);
\draw[->,poporange,thick] (conf.north) -- (0,-0.6);
\draw[->,poppurple,thick] (am.west) -- (2.2,0);
\end{tikzpicture}
\end{center}
\legende{Les quatre formes de l'intersubjectivité autour de deux sujets en interaction. La flèche centrale indique la réciprocité : chaque forme suppose que A et B se rapportent l'un à l'autre comme des consciences.}
\end{popfigure}

\courssub{Exemples camerounais : l'intersubjectivité en actes}

Les quatre formes ne sont pas des abstractions : elles structurent la vie sociale camerounaise quotidienne.

\textbf{La tontine} (ou \emph{djangui}, \emph{njangi}) est le chef-d'œuvre de l'intersubjectivité coopérative : un groupe de personnes cotise régulièrement une somme, et chacun reçoit à tour de rôle la totalité de la cagnotte. Ce système repose entièrement sur la \textbf{confiance mutuelle} : je donne mon argent aujourd'hui parce que je suis certain que les autres feront de même quand mon tour viendra.

\begin{exemplebox}[La tontine, coopération fondée sur la confiance]
Dans un quartier de Bafoussam, douze commerçantes se réunissent chaque dimanche pour leur tontine : chacune verse \num{10000} francs, et l'une d'elles reçoit \num{120000} francs pour investir dans son commerce. Aucune loi ne les y oblige : seule la parole donnée et la reconnaissance réciproque garantissent le système. La tontine illustre parfaitement une intersubjectivité de \textbf{coopération} où autrui est la condition de ma propre réussite économique.
\end{exemplebox}

\textbf{Les débats politiques} illustrent le conflit et la reconnaissance : les partis s'affrontent (conflit, rivalité pour le pouvoir), mais dans le cadre d'élections et d'institutions, chacun reconnaît l'autre comme adversaire légitime et les citoyens comme des sujets égaux dont la voix compte (reconnaissance).

\textbf{La famille élargie} illustre la solidarité intersubjective : au Cameroun, la famille ne se réduit pas aux parents et aux enfants ; elle inclut oncles, tantes, cousins, et chacun se sent responsable des autres (scolarisation d'un neveu, accueil d'un cousin en ville). C'est une intersubjectivité de solidarité où l'existence de chacun est liée à celle de tous.

\begin{popfigure}
\begin{center}
\begin{tikzpicture}
\begin{axis}[
    ybar, bar width=0.9cm,
    width=0.85\linewidth, height=6cm,
    ymin=0, ymax=80,
    ylabel={Proportion de la population adulte (\%)},
    symbolic x coords={Tontine, Association de quartier, Groupement religieux, Coop\'erative agricole},
    xtick=data,
    x tick label style={rotate=15,anchor=east,font=\scriptsize},
    ytick={0,20,40,60,80},
    nodes near coords, nodes near coords style={font=\scriptsize},
    every axis plot/.append style={fill=popblue,draw=popdark}
]
\addplot coordinates {(Tontine,62) (Association de quartier,45) (Groupement religieux,58) (Coop\'erative agricole,30)};
\end{axis}
\end{tikzpicture}
\end{center}
\legende{L'intersubjectivité coopérative en chiffres : proportion estimée de Camerounais adultes membres d'un groupe de solidarité (données indicatives reconstituées à des fins pédagogiques, en \%). La tontine arrive en tête : plus de six adultes sur dix y participent.}
\end{popfigure}

Ce graphique montre que la coopération intersubjective n'est pas un idéal abstrait : la majorité des Camerounais adultes participe à au moins un groupe fondé sur la confiance et la réciprocité. La vie sociale réelle confirme la thèse philosophique : l'homme vit toujours avec et parmi autrui.

\courssub{Méthode et entraînement}

\begin{methode}[Analyser une situation concrète d'intersubjectivité]
Pour appliquer la notion à une situation de la vie courante (exigence du savoir-faire au Probatoire et au Baccalauréat technique), procédez en quatre étapes :
\begin{enumerate}
\item \textbf{Identifier les sujets en présence} : qui est en relation avec qui ?
\item \textbf{Déterminer la forme dominante} : s'agit-il de communication et d'échange, de reconnaissance, de conflit, ou d'amour et d'amitié ?
\item \textbf{Vérifier la réciprocité} : les sujets se reconnaissent-ils mutuellement comme sujets, ou l'un est-il réduit au rang d'objet (domination) ?
\item \textbf{Conclure et justifier} : rédigez une conclusion qui montre ce que la situation révèle de la nature des rapports humains (autrui comme partenaire, obstacle ou condition).
\end{enumerate}
\end{methode}

\begin{exoresolu}[Analyser une situation concrète]
\textbf{Énoncé.} Dans un atelier de mécanique de Douala, un jeune apprenti travaille depuis deux ans sans être payé ; le patron ne lui apprend rien et le traite comme un simple porteur d'outils. Analysez cette situation à l'aide de la notion d'intersubjectivité.
\tcblower
\textbf{Solution détaillée.}
\begin{enumerate}
\item \emph{Sujets en présence} : le patron (maître d'apprentissage) et l'apprenti.
\item \emph{Forme dominante} : la situation devrait relever de la coopération au travail et de la communication (transmission d'un savoir-faire), mais elle est en réalité dominée par le \textbf{conflit} et la \textbf{domination}.
\item \emph{Réciprocité} : elle est rompue. Le patron ne reconnaît pas l'apprenti comme un sujet libre et égal : il le réduit au rang d'instrument, d'objet utile. On retrouve le rapport maître-esclave décrit par Hegel.
\item \emph{Conclusion} : cette situation montre une intersubjectivité dégradée. Pour qu'elle devienne authentique, il faudrait que le patron reconnaisse l'apprenti comme sujet (salaire, formation, respect) : la reconnaissance réciproque est la condition d'un rapport de travail juste.
\end{enumerate}
\end{exoresolu}

\begin{exercice}[À vous de jouer]
Analysez les trois situations suivantes en identifiant la forme d'intersubjectivité dominante et en justifiant votre réponse : (a) des élèves qui préparent ensemble un exposé en se répartissant les tâches ; (b) deux candidats à l'élection du délégué de classe qui se disputent violemment ; (c) une mère qui veille chaque soir sur son enfant malade.
\end{exercice}

\begin{corrige}[Exercice]
(a) \emph{Communication et coopération} : les élèves échangent et unissent leurs efforts en se reconnaissant comme partenaires égaux. (b) \emph{Conflit} : rivalité pour un même bien (le poste), mais les deux candidats se reconnaissent comme adversaires, donc comme sujets — l'intersubjectivité demeure, dégradée. (c) \emph{Amour} : union volontaire des consciences ; la mère veut le bien de l'enfant pour lui-même, sans rien attendre en échange.
\end{corrige}

\begin{aretenir}[Bilan de la section]
L'intersubjectivité est l'ensemble des rapports réciproques entre consciences qui se reconnaissent comme sujets. Structure fondamentale de l'existence humaine, elle prend quatre formes principales : la \textbf{communication} et l'échange, la \textbf{reconnaissance}, le \textbf{conflit}, et l'\textbf{amour} ou l'amitié. Ces formes, illustrées par la tontine, les débats politiques et la famille élargie au Cameroun, posent une dernière question : autrui est-il un obstacle à ma liberté, ou au contraire sa condition ? C'est l'objet de la section suivante, avec Hegel et Sartre.
\end{aretenir}

\coursec{Autrui : obstacle ou condition de ma liberté ? (Hegel, Sartre)}

\courssub{Transition : de l'intersubjectivité à la question de la liberté}

Après avoir analysé les modes de connaissance d'autrui (perception, analogie, empathie, langage) et défini l'intersubjectivité comme espace de constitution mutuelle du sens, nous devons maintenant affronter une tension majeure : \textbf{la présence d'autrui est-elle une menace pour ma liberté ou sa condition nécessaire ?} L'intersubjectivité n'est pas une donnée harmonieuse ; elle est le théâtre d'enjeux où se joue le statut même du sujet. Deux figures philosophiques majeures — Sartre et Hegel — éclairent ce débat sous des angles radicalement différents : pour Sartre, autrui est d'abord celui qui me \emph{limite} ; pour Hegel, il est celui qui me \emph{constitue}.

\courssub{La thèse du conflit : Sartre et le regard qui objectifie}

\begin{definition}[Le regard sartrien]
Pour Jean-Paul Sartre (\emph{L'Être et le Néant}, 1943), la relation à autrui n'est pas d'abord communication ou empathie, mais \cle{conflit}. Le \textbf{regard} d'autrui me découvre comme un \textbf{objet} dans son monde : il me dépossède de ma subjectivité (mon être-pour-soi) pour me figer en chose (être-en-soi).
\end{definition}

\begin{experience}[L'expérience du trou de serrure (Sartre)]
\textbf{Situation :} Je suis penché sur un trou de serrure, j'écoute, j'épie. Je suis tout entier absorbé par mon projet ; je \emph{suis} ma curiosité, ma liberté de choisir ce projet. Je n'ai pas de « moi » explicite, je suis conscience pure.
\textbf{Rupture :} J'entends des pas dans le couloir. Quelqu'un me regarde.
\textbf{Transformation :} Sous ce regard, je \textbf{deviens} un « voyeur ». Je suis saisi comme objet : « Je suis vu, donc je suis ». Ma liberté est niée ; je suis figé dans une étiquette (voyeur) que je n'ai pas choisie à cet instant. Je ressens la \textbf{honte} (ou la fierté si je veux être vu comme tel).
\textbf{Conclusion :} Autrui est le \emph{médiateur} par lequel je deviens un objet pour moi-même. Sans autrui, pas de « Moi » objectivable, mais aussi pas d'aliénation.
\end{experience}

\begin{propriete}[Structure du regard chez Sartre]
\begin{enumerate}
    \item \textbf{Objectivation :} Le regard me constitue comme \emph{être-pour-autrui}. Je ne suis plus le centre du monde, je deviens un point dans le monde d'autrui.
    \item \textbf{Aliénation :} Ma liberté est « volée » par la liberté d'autrui. Je suis \emph{jugé} sans pouvoir me défendre.
    \item \textbf{Conflit irréductible :} Soit je subis le regard (je suis objet), soit je tente de \emph{re-regarder} autrui pour le réduire à mon tour en objet (lutte pour la subjectivité). « L'enfer, c'est les autres » (\emph{Huis clos}, 1944) signifie : l'enfer, c'est l'impossibilité d'échapper au regard d'autrui qui me juge et me fige dans une image.
\end{enumerate}
\end{propriete}

\begin{attention}[Ne pas caricaturer Sartre]
La formule « L'enfer, c'est les autres » ne signifie pas qu'il faut haïr les gens ou vivre en ermite. Elle décrit la \textbf{structure} du rapport à autrui dans la vie quotidienne (l'être-pour-autrui). Sartre cherche une issue (la \emph{liberté concrète}, l'engagement, le « nous » du groupe en fusion) mais constate que dans l'immédiateté de la rencontre, le conflit est premier. Ne pas confondre la \emph{description} de la relation (ce qui est) et la \emph{morale} (ce qu'il faut faire).
\end{attention}

\begin{aretenir}[Sartre : autrui comme limite]
\begin{itemize}
    \item Autrui me \textbf{limite} : il borne ma liberté par son regard.
    \item Autrui me \textbf{constitue} malgré tout : sans lui, pas de « Moi » social, pas de responsabilité explicite.
    \item La relation authentique (amour, amitié, engagement commun) tente de \textbf{dépasser} l'objectivation par la reconnaissance réciproque des libertés, mais reste fragile.
\end{itemize}
\end{aretenir}

\courssub{La thèse de la reconnaissance : Hegel et la dialectique du Maître et de l'Esclave}

\begin{definition}[Reconnaissance (Hegel)]
Pour Hegel (\emph{Phénoménologie de l'Esprit}, 1807), la conscience de soi ne peut exister \emph{que} si elle est \textbf{reconnue} par une autre conscience de soi. « La conscience de soi existe en et pour soi lorsqu'elle existe pour une autre ; c'est-à-dire qu'elle n'existe qu'en étant reconnue. » Autrui n'est pas un obstacle : il est la \textbf{condition de possibilité} de mon identité.
\end{definition}

\begin{experience}[La dialectique du Maître et de l'Esclave (Hegel)]
\textbf{1. Rencontre :} Deux consciences de soi se font face. Chacune veut affirmer sa certitude d'elle-même comme universelle. Elles entrent en \textbf{lutte à mort} pour la reconnaissance pure (risquer sa vie prouve que je ne suis pas attaché à la vie biologique, mais à ma liberté).
\textbf{2. Issue :} L'une cède par peur de la mort (l'Esclave), l'autre risque tout (le Maître).
\textbf{3. Reconnaissance unilatérale (insatisfaisante) :}
\begin{itemize}
    \item Le \textbf{Maître} est reconnu par l'Esclave. Mais l'Esclave n'est pas une conscience \emph{libre} (il a cédé). La reconnaissance reçue par le Maître vaut donc \textbf{zéro} : on n'est vraiment reconnu que par un égal libre.
    \item L'\textbf{Esclave} n'est pas reconnu par le Maître (qui le nie). Mais par le \textbf{travail} (transformation de la nature), l'Esclave se \emph{reconnaît lui-même} dans l'objet produit. Il accède à l'indépendance et à la vérité de la conscience de soi.
\end{itemize}
\textbf{4. Dépassement :} La vraie reconnaissance est \textbf{réciproque} et \textbf{symétrique} (citoyenneté, droit, morale). Elle suppose l'égalité.
\end{experience}

\begin{popfigure}
\begin{tikzpicture}[
    node distance=1.5cm and 2.5cm,
    box/.style={draw=popdark, thick, rounded corners, fill=popcream, minimum width=2.8cm, minimum height=1.2cm, align=center, font=\small},
    arrow/.style={-Stealth, thick, popdark},
    labelarrow/.style={midway, fill=popcream, font=\scriptsize, inner sep=2pt}
]
% Nœuds
\node[box, align=center] (c1){Conscience A\\ (Désir de reconnaissance)};
\node[box, right=of c1, align=center] (c2){Conscience B\\ (Désir de reconnaissance)};
\node[box, below=of c1, align=center] (lutte){Lutte à mort\\ pour la reconnaissance};
\node[box, below=of lutte, align=center] (res){Résolution :\\ Maître / Esclave};
\node[box, left=of res, xshift=-1cm, align=center] (maitre){MAÎTRE\\ Reconnu, mais par un esclave\\ $\rightarrow$ Dépendance};
\node[box, right=of res, xshift=1cm, align=center] (esclave){ESCLAVE\\ Non reconnu par le Maître\\ Travail $\rightarrow$ Vérité de soi};
\node[box, below=of maitre, yshift=-0.5cm, align=center] (echec){Reconnaissance\\ UNILATÉRALE\\ \textbf{INSATISFAISANTE}};
\node[box, below=of esclave, yshift=-0.5cm, align=center] (reussite){Reconnaissance\\ RÉCIPROQUE\\ (Droit, Éthique)\\ \textbf{Condition de liberté}};
% Flèches
\draw[arrow] (c1) -- (c2) node[labelarrow, above] {Face à face};
\draw[arrow] (c1) -- (lutte);
\draw[arrow] (c2) -- (lutte);
\draw[arrow] (lutte) -- (res);
\draw[arrow] (res) -- (maitre);
\draw[arrow] (res) -- (esclave);
\draw[arrow] (maitre) -- (echec);
\draw[arrow] (esclave) -- (reussite);
\draw[arrow, popgreen] (echec) -- ++(0,-0.5) -| (reussite) node[labelarrow, near end, pos=0.75] {Dépassement dialectique};
\end{tikzpicture}
\legende{Schéma de la dialectique Maître/Esclave chez Hegel : la lutte pour la reconnaissance aboutit d'abord à une impasse (reconnaissance unilatérale) ; seule la reconnaissance réciproque (égalité) fonde la liberté véritable.}
\end{popfigure}

\begin{propriete}[Apports de Hegel]
\begin{enumerate}
    \item \textbf{Intersubjectivité constitutive :} Je ne suis sujet \emph{que} par l'autre. Pas de « Je » sans « Tu ».
    \item \textbf{Le travail comme médiation :} L'esclave, en transformant le monde, s'objective et se prouve à lui-même sa liberté. Le travail humanise.
    \item \textbf{Critique de la domination :} Le maître est \emph{dépendant} de l'esclave. La domination nie la condition de la reconnaissance (l'égalité). Elle est donc autodestructrice.
\end{enumerate}
\end{propriete}

\begin{savaistu}[Pour aller plus loin : Levinas et le visage d'autrui]
Emmanuel Levinas (\emph{Totalité et Infini}, 1961) prolonge ce débat : la relation à autrui n'est ni conflit (Sartre) ni dialectique (Hegel), mais \textbf{événement éthique}. Le \textbf{visage} d'autrui — sa nudité, sa vulnérabilité — m'adresse un commandement : \textbf{« Tu ne tueras point »}. Je suis \emph{responsable} d'autrui avant tout contrat et toute réciprocité. Cette ouverture, qui dépasse le strict programme de Première, pourra enrichir une conclusion de dissertation.
\end{savaistu}

\courssub{Synthèse comparative : conflit ou reconnaissance ?}

\begin{center}
\begin{tabularx}{\linewidth}{|>{\centering\arraybackslash}X|>{\centering\arraybackslash}X|>{\centering\arraybackslash}X|}
\hline
\rowcolor{popblueL}
\textbf{Dimension} & \textbf{Sartre (Conflit)} & \textbf{Hegel (Reconnaissance)} \\
\hline
\textbf{Nature du lien} & Ontologique : \emph{être-pour-autrui} (aliénation) & Dialectique : \emph{reconnaissance mutuelle} (constitution) \\
\hline
\textbf{Autrui est...} & Une \textbf{menace} pour ma subjectivité (le regard) & Un \textbf{partenaire nécessaire} (le miroir) \\
\hline
\textbf{Point clé} & Le \textbf{regard} (honte, objectivation) & La \textbf{lutte} et le \textbf{travail} (histoire) \\
\hline
\textbf{Liberté} & Menacée par la liberté d'autrui (conflit) & Réalisée \emph{dans} la reconnaissance réciproque (droit) \\
\hline
\textbf{Issue} & Dépassement par le \emph{nous} (engagement commun) & Dépassement par l'\textbf{État de droit} (reconnaissance universelle) \\
\hline
\textbf{Limite de la thèse} & Risque de guerre de tous contre tous & Risque d'aliénation dans l'institution \\
\hline
\end{tabularx}
\end{center}

\courssub{Application : les conflits communautaires au Cameroun comme échecs de la reconnaissance mutuelle}

\begin{exemplebox}[Le « vivre-ensemble » camerounais à l'épreuve de la philosophie]
\textbf{Contexte :} Le Cameroun traverse des tensions intercommunautaires (crise dans les régions du Nord-Ouest et du Sud-Ouest, tensions ethno-régionales, discours de haine sur les réseaux sociaux).
\textbf{Analyse sartrienne :} Chaque groupe regarde l'autre comme \textbf{objet} (« les sécessionnistes », « les francophones dominateurs »...). Le regard fige l'autre dans une \textbf{« essence »} réductrice (étiquette, stéréotype). C'est l'enfer sartrien : « L'enfer, c'est les autres » quand l'autre me nie ma liberté de me définir moi-même.
\textbf{Analyse hégélienne :} La crise révèle un \textbf{déficit de reconnaissance}. Les groupes ne se reconnaissent pas comme \emph{égaux en dignité} et en droits. La lutte pour la reconnaissance (revendications, manifestations) tourne à la « lutte à mort » symbolique ou réelle. Le « Maître » (groupe dominant) croit être reconnu, mais sans la reconnaissance libre de l'« Esclave » (groupe marginalisé), cette reconnaissance est vide. Le travail commun (développement, dialogue) est bloqué.
\textbf{Perspective :} Le « vivre-ensemble » ne peut être un simple slogan. Il exige :
\begin{enumerate}
    \item \textbf{Sortir du regard objectivant} (Sartre) : éducation à la nuance, refus des étiquettes.
    \item \textbf{Instituer la reconnaissance réciproque} (Hegel) : décentralisation effective, justice, équité dans le partage des ressources (reconnaissance institutionnelle).
    \item \textbf{Réhabiliter la responsabilité éthique} : protection des minorités, sanction des discours de haine, humanisation du débat public.
\end{enumerate}
\end{exemplebox}

\courssub{Méthode : dissertation — opposer, analyser, dépasser}

\begin{methode}[Traiter un sujet sur « Autrui et liberté » (ex. : « Autrui est-il une limite ou une condition de ma liberté ? »)]
\begin{enumerate}
    \item \textbf{Analyser les termes :} « Autrui » (autre sujet), « limite » (obstacle, borne), « condition » (ce qui rend possible), « liberté » (autonomie, pouvoir de choisir).
    \item \textbf{Problématiser :} Pourquoi la question se pose-t-elle ? Parce que l'expérience vécue oscille entre conflit (honte, oppression) et besoin vital (amitié, droit, amour). La liberté semble à la fois \emph{menacée} par autrui et \emph{réalisée} par lui.
    \item \textbf{Plan dialectique (thèse / antithèse / synthèse) :}
    \begin{itemize}
        \item \textbf{Thèse (Sartre) :} Autrui est une \textbf{limite}. Le regard objectifie, aliène, fige. Preuves : la honte, « l'enfer, c'est les autres ». La liberté pure serait solitude.
        \item \textbf{Antithèse (Hegel) :} Autrui est la \textbf{condition}. Sans reconnaissance, pas de conscience de soi, pas de droit, pas de liberté effective. La liberté \emph{concrète} est sociale (travail, droit, morale).
        \item \textbf{Synthèse :} Autrui est \textbf{les deux à la fois} : limite (je ne suis pas tout-puissant, je suis fini et vulnérable) et condition éthique et juridique. La vraie liberté n'est pas l'indépendance absolue, mais la \textbf{responsabilité assumée} dans un monde partagé. La loi protège la dignité de chacun.
    \end{itemize}
    \item \textbf{Conclure :} La liberté n'est pas « faire ce qu'on veut » (caprice), mais \textbf{s'autodéterminer dans un monde partagé}. Autrui est le \emph{miroir} (Hegel) et le \emph{juge} (Sartre) qui m'obligent à devenir un sujet libre et responsable.
\end{enumerate}
\end{methode}

\begin{exoresolu}[Sujet type Probatoire : « Reconnaître autrui, est-ce limiter sa propre liberté ? »]
\textbf{Énoncé :} « Reconnaître autrui, est-ce limiter sa propre liberté ? »
\tcblower
\textbf{Solution détaillée :}
\begin{enumerate}
    \item \textbf{Introduction :} Accroche (ex. : « Je ne suis pas seul au monde »). Définitions : reconnaître = accorder un statut, admettre l'existence et la valeur de l'autre ; liberté = autonomie, pouvoir d'agir. Problème : la reconnaissance semble imposer une contrainte (je dois tenir compte de l'autre), mais l'absence de reconnaissance nie l'humanité commune. Annonce du plan dialectique.
    \item \textbf{Thèse : oui, reconnaître autrui limite ma liberté immédiate (Sartre).} Reconnaître l'autre comme sujet, c'est accepter qu'il me regarde, me juge, résiste à mes projets. C'est renoncer à la toute-puissance du « pour-soi » solitaire. Exemple : la loi (reconnaissance juridique) m'interdit de voler, de tuer. C'est une limite \emph{extérieure}. Sartre : le regard d'autrui me fige, me fait honte, limite ma transparence à moi-même.
    \item \textbf{Antithèse : non, c'est la condition de ma liberté réelle (Hegel).} Sans reconnaissance, je ne suis pas un « Je » constitué. L'enfant non reconnu ne développe pas d'autonomie. Le citoyen sans droits n'a pas de liberté effective (seulement de la force). Le travail (Hegel) et le droit transforment la nature et les rapports sociaux pour \emph{étendre} mon pouvoir d'agir. La reconnaissance mutuelle (contrat, amitié, amour) \emph{augmente} ma capacité d'action (coopération).
    \item \textbf{Synthèse : reconnaître, c'est transformer une limite extérieure en limite intérieure choisie.} La limite n'est plus subie (contrainte), elle est \emph{consentie} comme loi morale : obéir à la loi que je me donne en tant qu'être raisonnable qui reconnaît les autres comme des fins. La responsabilité pour autrui devient la \emph{source} de ma subjectivité éthique : je suis libre \emph{parce que} je réponds de l'autre. La limite devient \textbf{horizon de sens}.
    \item \textbf{Conclusion :} Reconnaître autrui limite ma \emph{puissance} arbitraire mais fonde ma \emph{liberté} véritable : passer de l'indépendance solitaire à l'autonomie dans le droit et l'éthique.
\end{enumerate}
\end{exoresolu}

\begin{attention}[Pièges à éviter en dissertation]
\begin{itemize}
    \item \textbf{Confondre liberté et puissance :} Ne pas écrire « autrui me limite donc je ne suis pas libre ». La liberté n'est pas l'absence d'obstacles.
    \item \textbf{Opposer bêtement Sartre et Hegel :} Sartre n'est pas « anti-social », Hegel n'est pas « naïf ». Sartre cherche la réciprocité (le « nous ») ; Hegel montre que la reconnaissance unilatérale (maître/esclave) \emph{échoue}.
    \item \textbf{Généraliser sans exemples :} Toujours illustrer (le regard dans la rue, un procès, le vote, un conflit familial, la crise dans le Nord-Ouest et le Sud-Ouest).
    \item \textbf{Oublier la problématique :} Une copie qui juxtapose deux auteurs sans question directrice ne dépasse pas la moyenne.
\end{itemize}
\end{attention}

\begin{savaistu}[Kwame Anthony Appiah et la « conversation »]
Le philosophe ghanéen Kwame Anthony Appiah (\emph{Cosmopolitisme}, 2006) propose une voie pratique : la reconnaissance ne demande pas l'accord sur toutes les valeurs, mais la \textbf{conversation}. « La conversation n'a pas pour but de nous mettre d'accord, mais de nous habituer les uns aux autres. » Au Cameroun, le « vivre-ensemble » ne suppose pas l'uniformité, mais l'acceptation de la différence par le dialogue continu. C'est une éthique de la \textbf{cohabitation} qui rejoint la sagesse traditionnelle du \emph{njangi} (tontine) ou de la \emph{palabre} : on parle, on se reconnaît, on gère le conflit sans le nier.
\end{savaistu}

\coursec{Applications, méthode et entraînement au Probatoire}

Nous avons vu, avec Hegel et Sartre, qu'autrui est à la fois un obstacle à ma liberté (son regard m'objective) et la condition de mon existence comme conscience (je n'existe pleinement que reconnu par lui). Ce paradoxe n'est pas une simple curiosité théorique : il travaille notre vie quotidienne, nos conflits, nos échanges sur les réseaux sociaux et jusqu'à la politique nationale. Cette dernière section a pour but de vous rendre capables d'\emph{appliquer} les notions de la leçon à des situations concrètes, puis de construire une dissertation complète sur autrui, comme on vous le demandera au Probatoire (échéance de Première technique) et plus tard au Baccalauréat technique.

\courssub{Application 1 : la dispute de quartier et la dialectique de la reconnaissance}

Commençons par une situation que tout le monde connaît. Dans un quartier de Yaoundé, deux voisins se disputent violemment : l'un accuse l'autre d'avoir déplacé la borne qui sépare leurs parcelles. Chacun crie, chacun veut avoir le dernier mot, aucun n'écoute l'autre. La dispute dégénère et il faut l'intervention du chef de quartier.

Comment la philosophie nous aide-t-elle à comprendre cette scène banale ? Appliquons la \cle{dialectique du maître et de l'esclave} de Hegel. Chaque voisin ne se bat pas seulement pour quelques mètres de terrain : il se bat pour être \emph{reconnu}. Celui qui cède aura le sentiment d'être traité comme un être insignifiant, réduit au rang de chose. Chacun veut donc imposer sa vérité à l'autre, c'est-à-dire obtenir que l'autre reconnaisse sa supériorité, quitte à risquer le conflit physique. C'est exactement la « lutte à mort pour la reconnaissance » décrite par Hegel dans \emph{La Phénoménologie de l'Esprit}.

Mais la dialectique montre aussi l'impasse : si l'un des deux écrase l'autre et le réduit au silence, la reconnaissance obtenue ne vaut plus rien, car elle vient d'un être qu'on a abaissé. On ne peut être reconnu que par quelqu'un qu'on reconnaît soi-même comme libre et digne. La solution hégélienne passe donc par la \cle{reconnaissance réciproque} : l'intervention du chef de quartier, qui fait parler les deux parties et les amène à un compromis, institue un rapport où chacun reconnaît l'autre comme un interlocuteur légitime. Le droit et la médiation sont les formes civilisées de cette reconnaissance mutuelle.

\begin{exemplebox}[Analyse guidée]
Situation : deux élèves se disputent la place de chef de classe et se dénigrent devant les autres.
\begin{itemize}
\item \textbf{Étape 1} — Identifier l'enjeu : ce n'est pas la place elle-même, mais la reconnaissance (être vu, respecté, estimé par le groupe).
\item \textbf{Étape 2} — Appliquer la thèse : selon Hegel, chaque conscience veut être reconnue sans reconnaître l'autre ; selon Sartre, chacun cherche à réduire l'autre au rang d'objet par le regard et la parole blessante.
\item \textbf{Étape 3} — Tirer la leçon : le conflit ne se résout que si chacun accepte de voir dans l'autre un sujet et non un rival à écraser.
\end{itemize}
\end{exemplebox}

\courssub{Application 2 : les réseaux sociaux, empathie ou illusion de connaissance ?}

Sur WhatsApp ou Facebook, un jeune Camerounais peut avoir des centaines de « contacts » et d'« amis ». Mais connaît-il vraiment autrui à travers ces écrans ? C'est une excellente application de la question des \cle{modes de connaissance d'autrui}.

D'un côté, les réseaux sociaux développent une forme d'\cle{empathie} : on partage les joies (mariages, réussites au Probatoire) et les peines (deuils, maladies) de personnes éloignées ; les groupes WhatsApp de quartier ou de village organisent la solidarité (cotisations, tontines, levées de fonds pour un enterrement). Le langage, nous l'avons vu, est la voie royale d'accès à la pensée d'autrui : par les messages, autrui me dit ce qu'il ressent, ce que la seule perception de son corps ne pourrait jamais me révéler.

De l'autre côté, cette connaissance est fragile et souvent illusoire, pour trois raisons :
\begin{enumerate}
\item \textbf{L'image fabriquée} : sur les réseaux, chacun se présente sous son meilleur jour. Je ne perçois plus le corps vivant d'autrui, mais un profil construit, filtré, retouché. L'analogie (« il est comme moi ») se trompe plus facilement.
\item \textbf{L'absence du regard} : pour Sartre, c'est le regard de l'autre, sa présence corporelle, qui me révèle que je suis vu et qui fonde le rapport concret à autrui. Derrière un écran, ce regard fait défaut ; on peut bloquer, supprimer, ignorer autrui d'un clic, ce qui le réduit à un objet manipulable.
\item \textbf{Les malentendus} : sans le ton de la voix ni les gestes, un message ironique peut être pris pour une insulte ; les conflits éclatent plus vite, comme le montrent les disputes violentes dans certains groupes de discussion.
\end{enumerate}

Conclusion nuancée : les réseaux sociaux sont un \emph{instrument} de communication, donc un mode possible de connaissance d'autrui, mais ils ne remplacent pas la rencontre réelle où le corps, le regard et la voix rendent l'intersubjectivité pleinement effective.

\begin{attention}[Piège fréquent]
Ne rédigez pas un texte « pour ou contre les réseaux sociaux » comme on le ferait en français. Il faut \emph{philosopher} : mobiliser les notions du cours (perception, analogie, empathie, langage, regard) pour montrer \emph{en quoi} la connaissance d'autrui est transformée par ces outils. Le sujet n'est jamais la technique elle-même, mais le rapport à autrui.
\end{attention}

\courssub{Application 3 : le vivre-ensemble au Cameroun, un problème d'intersubjectivité}

Le Cameroun compte plus de 250 groupes ethniques, deux grandes langues officielles (le français et l'anglais) et plusieurs religions. Comment faire \emph{une} nation avec une telle diversité ? C'est le sens du mot d'ordre de « \cle{vivre-ensemble} », repris par les discours officiels, l'éducation civique et les campagnes pour la paix.

Philosophiquement, le vivre-ensemble est un problème d'\cle{intersubjectivité} : comment des sujets différents, qui ne se comprennent pas spontanément (langues, coutumes, histoires différentes), peuvent-ils former un « nous » ? Les notions de la leçon éclairent directement cette question :
\begin{itemize}
\item Par l'\emph{analogie}, je reconnais dans l'autre Camerounais, malgré sa langue ou son ethnie différentes, un être semblable à moi, qui souffre, espère et aime comme moi.
\item Par l'\emph{empathie} et le \emph{langage}, le dialogue entre communautés (bilinguisme, brassage scolaire, mariages mixtes) construit une compréhension réciproque.
\item Par la \emph{reconnaissance réciproque} (Hegel), aucune communauté ne doit chercher à dominer les autres : la cohésion nationale suppose que chacun reconnaisse la dignité égale de tous, ce que traduit politiquement l'égalité des citoyens devant la loi.
\end{itemize}

\begin{savaistu}[Le vivre-ensemble, une philosophie politique]
Le concept de « vivre-ensemble », devenu central dans le discours national camerounais (journées du vivre-ensemble, discours sur la paix et l'unité nationale), est une application politique directe de l'intersubjectivité étudiée dans cette leçon : il affirme que la nation n'existe que si les consciences se reconnaissent mutuellement comme sujets égaux, au-delà des différences ethniques, linguistiques et religieuses. Philosopher sur autrui, c'est donc aussi réfléchir aux conditions de la paix au Cameroun.
\end{savaistu}

\courssub{Méthode : construire une problématique à partir du paradoxe d'autrui}

La \cle{problématique} est la question directrice de la dissertation. Elle naît d'une difficulté, d'un paradoxe contenu dans le sujet. Or autrui est un concept naturellement paradoxal, ce qui facilite le travail :
\begin{itemize}
\item Autrui est à la fois \textbf{proche} (un semblable, un autre moi-même) et \textbf{lointain} (sa conscience m'est à jamais cachée, je ne peux pas entrer dans sa tête).
\item Autrui est à la fois \textbf{semblable} (je le comprends par analogie avec moi) et radicalement \textbf{autre} (il me résiste, il me surprend, il me juge).
\item Autrui est à la fois \textbf{obstacle} à ma liberté (son regard m'enferme, dit Sartre) et \textbf{condition} de ma liberté (il me reconnaît et m'aide à devenir moi-même, dit Hegel).
\end{itemize}

\begin{methode}[Construire une problématique en quatre étapes]
\begin{enumerate}
\item \textbf{Dégager le sens commun} : que croit-on spontanément ? (Ex. : « autrui est un être comme moi, que je connais facilement ».)
\item \textbf{Trouver la difficulté} : qu'est-ce qui contredit cette évidence ? (Ex. : « pourtant je ne peux jamais éprouver directement sa pensée ; il m'arrive de me tromper sur lui ».)
\item \textbf{Formuler la tension} sous forme d'opposition : proche/lointain, semblable/autre, obstacle/condition.
\item \textbf{Écrire la question} qui résume la tension : « Autrui m'étant à la fois semblable et irréductiblement autre, puis-je prétendre le connaître ? »
\end{enumerate}
\end{methode}

\begin{popfigure}
\begin{center}
\begin{tikzpicture}[
  grow=down,
  level 1/.style={sibling distance=7.2cm, level distance=1.6cm},
  level 2/.style={sibling distance=3.6cm, level distance=1.9cm},
  every node/.style={draw=popblue, thick, rounded corners=2pt, fill=popblueL, text width=3.1cm, align=center, font=\small},
  root/.style={draw=poporange, fill=poporangeL, text width=5.2cm, font=\small\bfseries},
  edge from parent/.style={draw=popdark, thick, -{Stealth[length=2mm]}}
]
\node[root]{Qu'est-ce qu'autrui pour moi ?}
  child { node[align=center]{Autrui comme \textbf{connaissance}\\ (puis-je le connaître ?)}
    child { node {Perception et analogie : je devine sa pensée d'après son corps} }
    child { node {Empathie et langage : je partage ses émotions, il me parle} }
  }
  child { node[align=center]{Autrui comme \textbf{rapport}\\ (que m'est-il ?)}
    child { node {Sartre : obstacle, son regard m'objective} }
    child { node {Hegel : condition, je n'existe que reconnu} }
  };
\end{tikzpicture}
\end{center}
\legende{Arbre de la problématique : la question fondamentale « Qu'est-ce qu'autrui pour moi ? » se divise en deux branches (la connaissance d'autrui et le rapport à autrui), chacune portant les thèses étudiées dans la leçon.}
\end{popfigure}

\courssub{Méthode : rédiger une conclusion efficace}

La conclusion n'est pas un résumé mécanique. Elle doit remplir trois fonctions :
\begin{enumerate}
\item \textbf{Répondre clairement à la problématique} annoncée dans l'introduction. Si la question était « autrui est-il un obstacle à ma liberté ? », la conclusion doit trancher, par exemple : « autrui limite ma liberté par son regard, mais il la rend possible par la reconnaissance ; il est donc moins un obstacle qu'une condition paradoxale de mon existence libre ».
\item \textbf{Dégager la portée} de cette réponse : qu'avons-nous gagné à philosopher ? (Nous comprenons mieux nos conflits, notre besoin d'estime, notre vie sociale.)
\item \textbf{Ouvrir sur une question nouvelle}, proche mais distincte : la justice (que devons-nous à autrui ?), l'amour (aimer autrui, est-ce le connaître autrement ?), la communication (le langage suffit-il à abolir la distance entre les consciences ?).
\end{enumerate}

\begin{exemplebox}[Exemple de conclusion rédigée]
« Ainsi, autrui m'apparaît d'abord comme un être fermé, dont je ne perçois que le corps, et même comme un rival qui menace ma liberté par son regard. Pourtant, l'analogie, l'empathie et surtout le langage rendent sa pensée accessible, et la reconnaissance réciproque fait de lui non plus un obstacle, mais la condition de mon humanité. Autrui est donc l'autre sans lequel je ne serais pas moi-même. Reste alors une question : si autrui m'est si nécessaire, quels devoirs ai-je envers lui ? C'est le problème de la justice. »
\end{exemplebox}

\courssub{Entraînement au Probatoire : sujets types et plan de dissertation}

\begin{methode}[Structure type d'une dissertation sur autrui]
Pour la plupart des sujets sur autrui, un plan dialectique en trois parties est efficace :
\begin{itemize}
\item \textbf{I. Thèse} : autrui m'échappe. Sa conscience est invisible ; je ne perçois que des signes (corps, gestes, paroles) ; son regard m'objective et limite ma liberté (Sartre).
\item \textbf{II. Antithèse} : autrui m'est pourtant accessible. L'analogie, l'empathie et le langage me donnent prise sur sa vie intérieure ; sans lui, pas de langage, pas de culture, pas de moi-même.
\item \textbf{III. Synthèse} : la reconnaissance réciproque comme solution. Autrui n'est ni un objet à connaître totalement ni un ennemi à fuir : c'est un sujet à reconnaître, et c'est dans cette reconnaissance mutuelle que se construisent la liberté et le vivre-ensemble (Hegel).
\end{itemize}
\end{methode}

\begin{popfigure}
\begin{center}
\begin{tikzpicture}[
  box/.style={draw=popblue, thick, rounded corners=3pt, fill=popblueL, text width=4.6cm, align=left, font=\small, inner sep=5pt},
  lab/.style={font=\small\bfseries, text=popdark},
  arr/.style={-{Stealth[length=2.5mm]}, very thick, draw=poporange}
]
\node[box, align=center] (intro) at (0,0){\textbf{Introduction}\\ Accroche (dispute, regard) ; définition d'autrui ; problématique : obstacle ou condition ?};
\node[box, align=center] (p1) at (0,-2.3){\textbf{I. Autrui m'échappe}\\ Conscience invisible ; le regard de Sartre m'enferme ; conflit des consciences.};
\node[box, align=center] (p2) at (0,-4.6){\textbf{II. Autrui m'est accessible}\\ Analogie, empathie, langage : je comprends et je partage sa vie intérieure.};
\node[box, align=center] (p3) at (0,-6.9){\textbf{III. La reconnaissance réciproque}\\ Hegel : je n'existe que reconnu ; autrui, condition de ma liberté et du vivre-ensemble.};
\node[box, fill=poporangeL, draw=poporange, align=center] (ccl) at (0,-9.2){\textbf{Conclusion}\\ Réponse à la problématique + ouverture (justice, amour, communication).};
\draw[arr] (intro) -- (p1);
\draw[arr] (p1) -- (p2);
\draw[arr] (p2) -- (p3);
\draw[arr] (p3) -- (ccl);
\end{tikzpicture}
\end{center}
\legende{Schéma d'un plan de dissertation rédigé sur le sujet « Autrui est-il un obstacle à ma liberté ? » : plan dialectique en trois parties, de l'introduction problématisée à la conclusion ouverte.}
\end{popfigure}

\begin{exemplebox}[Barème type au Probatoire]
La copie de dissertation est généralement évaluée selon des critères voisins de ceux-ci :
\begin{center}
\begin{tabularx}{\linewidth}{|X|c|}
\hline
\rowcolor{popblueL} \textbf{Critère d'évaluation} & \textbf{Points} \\
\hline
Pertinence de la problématique et de l'introduction & 4 pts \\
\hline
Exactitude des références (auteurs, thèses, exemples) & 4 pts \\
\hline
Qualité de l'argumentation (articulation des parties, démonstration) & 8 pts \\
\hline
Qualité de l'expression (langue, clarté, présentation) & 4 pts \\
\hline
\textbf{Total} & \textbf{20 pts} \\
\hline
\end{tabularx}
\end{center}
On remarque que l'argumentation pèse le plus lourd (8 points) : ce n'est pas ce que vous récitez qui compte, mais la manière dont vous construisez votre réponse.
\end{exemplebox}

\begin{attention}[L'erreur qui coûte le plus cher]
L'erreur la plus fréquente consiste à \emph{réciter les thèses} (Hegel a dit ceci, Sartre a dit cela) sans jamais les articuler à la problématique ni les illustrer par des exemples. Une copie qui aligne les noms d'auteurs sans répondre à la question posée est une copie hors sujet. Chaque référence doit être introduite, expliquée en une phrase, reliée au problème, et accompagnée d'un exemple concret (dispute de quartier, regard d'un supérieur, échange sur WhatsApp).
\end{attention}

\begin{exoresolu}[Sujet type : « Peut-on connaître autrui ? »]
Tracez le plan détaillé de ce sujet et rédigez la problématique.
\tcblower
\textbf{Analyse du sujet} : « connaître » signifie ici accéder à la pensée et aux sentiments d'autrui ; « peut-on » invite à examiner à la fois la possibilité et les limites de cette connaissance.

\textbf{Problématique} : autrui m'est donné comme un semblable, pourtant sa conscience reste invisible ; dès lors, puis-je prétendre le connaître vraiment, ou ma connaissance d'autrui reste-t-elle une simple supposition ?

\textbf{Plan dialectique} :
\begin{itemize}
\item \textbf{I. Autrui est inconnaissable en lui-même.} Je ne perçois que son corps ; sa conscience ne tombe sous aucun de mes sens. Je peux me tromper (hypocrisie, malentendu, mensonge). Le regard de Sartre me montre un autrui qui m'échappe et me juge sans que je puisse le saisir.
\item \textbf{II. Pourtant, une connaissance d'autrui est possible.} Par l'analogie, je lui attribue des sentiments semblables aux miens ; par l'empathie, je partage sa joie ou sa douleur ; par le langage, il me dit directement ce qu'il pense — ce que je ne pourrais jamais deviner autrement.
\item \textbf{III. Connaître autrui, c'est le reconnaître.} La connaissance d'autrui n'est pas une science exacte mais une relation : elle progresse dans le dialogue et la reconnaissance réciproque (Hegel). On connaît autrui moins en l'observant qu'en vivant avec lui.
\end{itemize}
\textbf{Conclusion attendue} : oui, on peut connaître autrui, mais imparfaitement et autrement qu'on ne connaît les choses ; ouverture : l'amour donne-t-il une connaissance plus profonde d'autrui ?
\end{exoresolu}

\begin{exercice}[Sujets d'entraînement]
Traitez l'un des trois sujets suivants (introduction problématisée + plan détaillé ; la rédaction complète pourra être faite en devoir) :
\begin{enumerate}
\item « Autrui est-il un obstacle à ma liberté ? »
\item « Peut-on connaître autrui ? »
\item « Avons-nous besoin d'être reconnus ? »
\end{enumerate}
Pour le sujet 3, pensez à mobiliser la dialectique hégélienne de la reconnaissance et des exemples camerounais (recherche de prestige, réseaux sociaux, honneur dans la famille).
\end{exercice}

\begin{corrige}[Exercice — indications pour le sujet 3]
« Avons-nous besoin d'être reconnus ? »
\begin{itemize}
\item \textbf{Problématique} : la reconnaissance semble un besoin universel (estime, honneur, statut) ; mais est-elle indispensable à l'existence humaine, ou peut-on s'en passer ?
\item \textbf{I.} Oui : selon Hegel, la conscience de soi n'existe que reconnue par une autre conscience ; l'enfant privé de tout regard bienveillant ne se construit pas ; les conflits sociaux naissent du mépris (exemples : humiliation, disputes de quartier, quête de « likes »).
\item \textbf{II.} Nuance : la recherche excessive de reconnaissance rend dépendant du regard d'autrui (Sartre : l'enfer, c'est les autres) ; le sage peut viser une estime de soi indépendante du jugement d'autrui.
\item \textbf{III.} Synthèse : nous avons besoin d'être reconnus, mais d'une reconnaissance réciproque et non servile ; la dignité consiste à reconnaître autrui autant qu'à être reconnu par lui.
\end{itemize}
\end{corrige}

\courssub{Bilan de la leçon et lien avec la leçon suivante}

\begin{aretenir}[Bilan de la leçon « Autrui »]
\begin{itemize}
\item \textbf{Connaître autrui} : je n'accède pas directement à sa conscience ; je le connais par la perception de son corps, par l'analogie avec moi-même, par l'empathie et surtout par le langage.
\item \textbf{L'intersubjectivité} : le rapport entre sujets n'est pas un rapport entre choses ; il fonde la communication, la vie sociale et le vivre-ensemble.
\item \textbf{Le paradoxe d'autrui} : obstacle à ma liberté par son regard (Sartre), il en est aussi la condition par la reconnaissance réciproque (Hegel).
\item \textbf{Méthode} : toute dissertation sur autrui part de ce paradoxe, oppose les thèses, et conclut par la reconnaissance mutuelle, avec des exemples concrets.
\end{itemize}
\end{aretenir}

Cette leçon nous a montré que je ne puis être moi-même sans autrui : mon identité se construit dans le regard et la reconnaissance d'autrui. Mais si autrui m'est nécessaire, une nouvelle question se pose : que dois-je à autrui ? Mes rapports avec lui sont-ils seulement des rapports de connaissance et de reconnaissance, ou sont-ils aussi réglés par des devoirs ? C'est le passage de la question de la connaissance à celle de la morale, que la leçon suivante abordera en interrogeant les fondements de nos obligations envers les autres.

\newpage

\coursec{Activité d'intégration}

\courssub{Objectifs de l'activité}
\begin{itemize}
    \item \textbf{Mobiliser les savoirs} de la leçon : distinguer les modes de connaissance d'autrui (perception, analogie, empathie, langage), identifier les formes de l'intersubjectivité (conflit, reconnaissance, coopération) et confronter les thèses de \cle{Hegel}, \cle{Sartre} et \cle{Lévinas} sur le rapport à autrui.
    \item \textbf{Développer les savoir-faire} officiels : appliquer des notions abstraites à une situation concrète vécue au lycée (transfert), analyser des documents hétérogènes (image, texte, dialogue), construire un tableau comparatif rigoureux et rédiger une \cle{mini-dissertation} structurée (introduction problématisée, développement argumenté, conclusion).
    \item \textbf{Travailler l'argumentation orale et écrite} : justifier chaque interprétation par une référence précise au texte ou à l'image, écouter les autres groupes pour enrichir sa propre réflexion (intersubjectivité épistémique) et présenter collectivement une synthèse cohérente.
\end{itemize}

\courssub{Situation}
\textbf{Contexte :} Nous sommes au \emph{Lycée Technique de Bepanda} (Douala), un mardi de mars, 11 h 30. La sonnerie de la récréation retentit. La cour est bondée, bruyante, traversée de flux d'élèves en tenue kaki et bleu. Trois « documents » sont soumis à votre analyse philosophique.

\begin{exemplebox}[Document 1 : Photographie — « La file d'attente à la cantine »]
\emph{Description de la scène (photo prise par un élève du club presse) :} Au premier plan, trois élèves (deux garçons, une fille) se font face à moins de 50 cm. L'élève du centre (chemise kaki déboutonnée, brassard de délégué) a le bras tendu, paume ouverte vers l'avant, bloquant le passage. L'élève de gauche (fille, sac à dos rose) recule d'un pas, le buste penché en arrière, visage fermé, lèvres pincées. L'élève de droite (garçon, casquette à l'envers) pointe un doigt vers le délégué, bouche ouverte en plein cri. Autour, une dizaine d'élèves observent : certains rient (téléphones sortis), d'autres détournent le regard, deux tentent d'intervenir en posant une main sur l'épaule du délégué. Au fond, le mur de la cantine porte une affiche fanée : « \textbf{Respect – Discipline – Excellence} ».
\end{exemplebox}

\begin{exemplebox}[Document 2 : Dialogue rapporté — « La place du fond »]
\emph{Scène : salle d'informatique, poste n\textsuperscript{o} 12. Deux élèves de Première TMA (Maintenance), \textbf{Kamel} et \textbf{Esther}, se disputent l'unique poste fonctionnel avec écran plat.}
\begin{itemize}
    \item \textbf{Kamel} (s'installant) : « Dégage, c'est mon poste. J'ai réservé hier soir. »
    \item \textbf{Esther} (ne bougeant pas, les yeux dans les siens) : « Tu mens. Tu n'étais pas là. Et puis, tu ne me regardes même pas quand tu parles. Tu me parles comme si j'étais une chaise. »
    \item \textbf{Kamel} (évitant le regard, tapant sur le clavier) : « Ce n'est pas mon problème. Le prof a dit : premier arrivé, premier servi. Lève-toi. »
    \item \textbf{Esther} (se levant lentement, voix basse) : « Tu as gagné. Mais tu as perdu quelque chose. »
    \item \textbf{Kamel} (sans tourner la tête) : « J'ai gagné le poste. C'est tout ce qui compte. »
\end{itemize}
\end{exemplebox}

\begin{exemplebox}[Document 3 : Extrait commenté — Hegel, \textit{Phénoménologie de l'esprit} (1807), dialectique du maître et de l'esclave (traduction adaptée)]
\emph{« La conscience de soi existe en et pour soi lorsqu'elle existe \textbf{pour une autre} conscience de soi ; c'est-à-dire qu'elle n'existe qu'en étant reconnue. [...] L'individu qui n'a pas risqué sa vie ne peut être reconnu comme conscience de soi indépendante. [...] Le maître est le maître seulement par rapport à l'esclave ; l'esclave est l'esclave seulement par rapport au maître. Mais dans cette relation, c'est l'esclave qui, par le travail et la peur de la mort, transforme le monde et accède à une vérité que le maître, dans sa jouissance stérile, ne peut atteindre. La reconnaissance unilatérale est une reconnaissance inachevée. Seule la \textbf{reconnaissance mutuelle} permet à chacune des deux consciences d'être pleinement elle-même. »}

\emph{— Note de cours : Hegel montre que je ne deviens « Je » que si un autre « Je » me reconnaît comme tel. Le conflit initial (lutte à mort) est le passage obligé vers une intersubjectivité supérieure : la reconnaissance réciproque.}
\end{exemplebox}

\courssub{Consignes}
\emph{Travail en groupes de 4 ou 5 élèves (20 min de préparation, 10 min de restitution par groupe). Chaque groupe rend une fiche unique.}

\begin{enumerate}
    \item \textbf{Modes de connaissance d'autrui (Doc 1, 2, 3).} Pour chaque document, identifiez \textbf{le ou les modes principaux} mobilisés par les acteurs (ou l'auteur) pour connaître l'autre : \cle{perception sensible} (le corps, le visage, le regard), \cle{analogie} (j'infère son intériorité à partir de mon propre cas), \cle{empathie} (ressentir \emph{avec} l'autre), \cle{langage/communication} (parole, dialogue, silence signifiant). Justifiez par un élément précis du document.
    \item \textbf{Formes de l'intersubjectivité.} Qualifiez la relation intersubjective dominante dans chaque document : \cle{conflit} (lutte pour la reconnaissance, objectivation), \cle{reconnaissance} (mutuelle ou unilatérale, lutte pour la vie ou la mort chez Hegel), \cle{coopération} (action commune, entraide). Expliquez pourquoi la situation n'est pas réductible à une seule forme (exemple : le conflit chez Hegel \emph{contient} la promesse de la reconnaissance).
    \item \textbf{Tableau comparatif appliqué : Sartre / Hegel / Lévinas face à la situation de la cantine (Doc 1).} Complétez le tableau en 3 colonnes (Auteur / Thèse centrale sur autrui / Application à la scène de la cantine : le délégué, la fille, le garçon à la casquette).
    \item \textbf{Mini-dissertation collective (250 à 350 mots).} Sujet : \textbf{« Autrui est-il un obstacle ou une condition de ma liberté ? »} Structure imposée :
    \begin{enumerate}
        \item \textbf{Introduction} : accroche (ancrage dans la situation), définition des termes (autrui, liberté, obstacle/condition), \textbf{problématique} (formulée en question), \textbf{annonce du plan} (2 ou 3 parties).
        \item \textbf{Développement} : \textbf{un seul paragraphe argumenté} (thèse / antithèse / synthèse ou démarche dialectique) mobilisant \textbf{au moins deux auteurs} (Hegel, Sartre, Lévinas) et \textbf{un exemple précis} tiré des documents 1 ou 2.
        \item \textbf{Conclusion} : réponse synthétique à la problématique, ouverture (exemple : éthique de la reconnaissance au lycée, citoyenneté).
    \end{enumerate}
\end{enumerate}

\courssub{Ressources à mobiliser}
\begin{definition}[Modes de connaissance d'autrui]
\textbf{Perception sensible} : je vois le corps vivant, le visage, le regard de l'autre (Husserl, Merleau-Ponty). \textbf{Analogie} : argument par analogie (Mill) : « Comme moi, l'autre a un corps ; comme moi, il a une conscience. » \textbf{Empathie} (\emph{Einfühlung}) : saisie immédiate, non inférentielle, de l'affect d'autrui (Scheler, Stein). \textbf{Langage} : l'autre se donne à comprendre dans le dialogue ; la communication institue un « nous » (Gadamer, Ricœur).
\end{definition}

\begin{definition}[Intersubjectivité]
Espace de sens partagé entre sujets. Trois figures types au programme : \textbf{Conflit} (Sartre : le regard m'objective, je suis « pour-autrui » malgré moi) ; \textbf{Reconnaissance} (Hegel : je ne suis libre que si l'autre me reconnaît libre ; lutte à mort, puis travail, puis reconnaissance mutuelle) ; \textbf{Coopération / Éthique} (Lévinas : le visage de l'autre m'ordonne « Tu ne tueras point », responsabilité antérieure à ma liberté).
\end{definition}

\begin{aretenir}[Repères auteurs pour le tableau et la dissertation]
\begin{itemize}
    \item \textbf{Sartre} (\textit{L'Être et le Néant}, 1943 ; \textit{Huis clos}, 1944) : « L'enfer, c'est les autres. » Le \cle{regard} d'autrui me fige en objet (je suis vu), il aliène ma subjectivité. La liberté tente alors de reprendre l'initiative : fuir le regard ou objectiver l'autre à mon tour. Conflit originel.
    \item \textbf{Hegel} (\textit{Phénoménologie de l'esprit}, 1807) : autrui est \textbf{condition} de ma liberté. Pas de conscience de soi sans reconnaissance. Dialectique du maître et de l'esclave : le conflit initial (risque de la vie) débouche, par le travail, sur la reconnaissance mutuelle.
    \item \textbf{Lévinas} (\textit{Totalité et Infini}, 1961) : autrui n'est ni un obstacle ni un partenaire de contrat, mais l'\cle{Autre} absolu (le visage). Ma liberté se fonde sur ma \textbf{responsabilité} pour lui. L'éthique précède l'ontologie.
\end{itemize}
\end{aretenir}

\courssub{Démarche et corrigé commenté}
\begin{exoresolu}[Corrigé guidé de l'activité d'intégration]
\textbf{Étape 1 — Modes de connaissance d'autrui (Consigne 1)}
\begin{itemize}
    \item \textbf{Doc 1 (Photo)} : \cle{perception sensible} dominante. Les corps parlent : bras tendu (blocage), recul du buste (rejet), doigt pointé (accusation), regards (fuyants, fixes, moqueurs). L'affiche « Respect » agit comme un \textbf{signifiant visuel} (langage écrit) perçu mais ignoré. \emph{Analogie} implicite : chaque observateur projette sa propre réaction (« Et moi, qu'aurais-je fait ? »).
    \item \textbf{Doc 2 (Dialogue)} : \cle{langage/communication} (paroles explicites, mensonge, silence final) et \cle{empathie} \textbf{manquée} ou \textbf{refusée}. Esther réclame le regard (« Tu ne me regardes même pas ») : elle exige une \textbf{reconnaissance intersubjective} par le visage (Lévinas). Kamel fuit le regard : refus d'empathie, traitement d'Esther comme \textbf{objet} (moyen pour obtenir le poste). L'analogie est bloquée par l'égoïsme.
    \item \textbf{Doc 3 (Hegel)} : connaissance \cle{conceptuelle et dialectique}. Hegel ne perçoit pas un individu concret ; il \textbf{reconstruit logiquement} la structure nécessaire de la relation à autrui. La « connaissance » ici est philosophique : l'intersubjectivité est \textbf{condition de possibilité} de la conscience de soi.
\end{itemize}

\textbf{Étape 2 — Formes de l'intersubjectivité (Consigne 2)}
\begin{center}
\begin{tabularx}{\linewidth}{|l|X|X|}\hline
\rowcolor{popblueL}\textbf{Document} & \textbf{Forme dominante} & \textbf{Justification / Nuance} \\ \hline
Doc 1 (Photo) & \cle{Conflit} (lutte pour la place, reconnaissance unilatérale exigée par le délégué) & Présence latente d'une \textbf{coopération} naissante (les deux élèves qui interviennent) et d'une \textbf{reconnaissance} symbolique (l'affiche « Respect » comme idéal partagé, mais non réalisé). \\ \hline
Doc 2 (Dialogue) & \cle{Conflit} asymétrique (Kamel : force et droit formel ; Esther : appel éthique au visage) & Bascule vers une \textbf{reconnaissance ratée} : Esther sort du conflit par le haut (« Tu as gagné... mais tu as perdu »), esquissant une \textbf{dimension éthique} au sens de Lévinas (responsabilité pour l'autre, même vaincu). \\ \hline
Doc 3 (Hegel) & \cle{Lutte pour la reconnaissance} puis \textbf{reconnaissance mutuelle} (visée finale) & Le conflit (lutte à mort) est le \textbf{moment nécessaire} pour accéder à la vérité de l'intersubjectivité : la liberté ne s'affirme que \emph{par} et \emph{pour} l'autre. \\ \hline
\end{tabularx}
\end{center}

\textbf{Étape 3 — Tableau comparatif Sartre / Hegel / Lévinas appliqué à la cantine (Consigne 3)}
\begin{center}
\begin{tabularx}{\linewidth}{|l|X|X|}\hline
\rowcolor{popblueL}\textbf{Auteur} & \textbf{Thèse centrale sur autrui} & \textbf{Application à la scène (Doc 1)} \\ \hline
\textbf{Sartre} & Autrui = \textbf{obstacle} radical. Le regard me chosifie : je deviens « cet objet vu » (le délégué vu comme « tyran », la fille vue comme « victime »). Liberté = tenter de reprendre ma subjectivité (regarder à mon tour, fuir, ou nier l'autre). & Le délégué \textbf{regarde} la fille pour la figer en « celle qui doit obéir ». La fille \textbf{subit} le regard (recul, lèvres pincées). Le garçon à la casquette \textbf{renvoie le regard} (doigt pointé, cri) : tentative de \textbf{récupérer sa liberté} par la contre-objectivation (faire du délégué un objet « abusif »). L'enfer, c'est la file d'attente. \\ \hline
\textbf{Hegel} & Autrui = \textbf{condition} de ma liberté. Je ne suis « Je » que si un autre « Je » me reconnaît. Le conflit (ici : bras tendu contre recul) est l'\textbf{épreuve de la vérité} : chacun risque sa « vie sociale » (réputation, place). La solution n'est pas la fuite mais la \textbf{reconnaissance mutuelle} (négociation, règle acceptée par tous). & Le délégué exige une \textbf{reconnaissance unilatérale} (position de maître) : « Je décide, vous obéissez. » La fille et le garçon refusent cette place d'esclave. Les deux intervenants (mains sur l'épaule) tentent d'instaurer une \textbf{médiation} (règle partagée) pour passer du conflit à la \textbf{reconnaissance réciproque} : « Nous sommes tous élèves, nous avons tous faim, faisons la queue. » \\ \hline
\textbf{Lévinas} & Autrui = \textbf{visage} qui me commande « Tu ne tueras point ». Ma liberté se \textbf{fonde} sur ma responsabilité pour lui (l'éthique comme philosophie première). L'autre n'est pas un thème pour ma liberté : il m'assigne à répondre de lui. & Le \textbf{visage de la fille} (recul, lèvres pincées, souffrance muette) \textbf{ordonne} au délégué : « Ne m'écrase pas. » Le délégué \textbf{fait violence} à cet ordre éthique en transformant le visage en obstacle. Le garçon à la casquette, en pointant le doigt, \textbf{dénonce} la violation de l'éthique. La vraie grandeur du délégué serait de \textbf{baisser le bras} par responsabilité pour le visage de l'autre. \\ \hline
\end{tabularx}
\end{center}

\textbf{Étape 4 — Mini-dissertation type (Consigne 4)}

\emph{Voici une proposition de rédaction collective (environ 300 mots), respectant la structure imposée.}

\begin{quote}
\textbf{Introduction.} Dans la cour du Lycée Technique de Bepanda, un délégué bloque le passage : son bras tendu fige une camarade en « obstacle » à son désir de manger vite. Cette scène banale illustre le paradoxe central de l'intersubjectivité : l'autre est à la fois celui sans qui je ne suis rien (Hegel) et celui qui menace ma souveraineté (Sartre). \textbf{Autrui} désigne toute conscience distincte de la mienne ; \textbf{la liberté} est le pouvoir de se déterminer soi-même. \textbf{Problématique} : la relation à autrui est-elle structurellement aliénante (obstacle) ou constitutivement fondatrice (condition) de ma liberté ? \textbf{Plan} : nous montrerons d'abord que le regard d'autrui m'objective et menace ma liberté (Sartre), avant de démontrer que seule la reconnaissance mutuelle (Hegel) ou la responsabilité éthique (Lévinas) permet d'accéder à une liberté véritable.

\textbf{Développement.} Pour Sartre, dans \textit{L'Être et le Néant}, le \cle{regard} d'autrui est l'expérience originaire de mon aliénation : quand le délégué fixe la fille, il la constitue en « être-pour-autrui » ; elle devient un objet dans son monde et perd sa transcendance. C'est l'enfer : « L'enfer, c'est les autres. » La liberté de la fille semble anéantie ; elle ne peut que fuir ou tenter de retourner le regard (ce que fait le garçon à la casquette). Mais cette thèse du conflit permanent néglige que le délégué lui-même a \textbf{besoin} de la fille pour être « délégué » : sans élève à commander, son brassard ne signifie rien. C'est la leçon de \textbf{Hegel} : dans la dialectique du maître et de l'esclave, le maître qui nie l'esclave se nie lui-même, car il ne reçoit qu'une reconnaissance unilatérale, vide. La vraie liberté naît de la \textbf{reconnaissance mutuelle} : les deux intervenants qui posent la main sur l'épaule du délégué tentent d'instaurer cette universalité concrète (« Nous sommes tous égaux devant la faim »). Plus radicalement, \textbf{Lévinas} renverse la perspective : le visage nu de la fille (sa vulnérabilité) ne limite pas ma liberté, il la \textbf{fonde} en m'assignant la responsabilité de ne pas la blesser. Ma liberté n'est pas « faire ce que je veux » (Kamel au poste 12), mais \textbf{répondre} à l'appel de l'Autre.

\textbf{Conclusion.} Autrui n'est donc pas un simple obstacle à contourner (thèse sartrienne du regard objectivant), mais la \textbf{condition} de ma liberté : sans l'autre qui me reconnaît (Hegel) ou qui m'appelle à répondre (Lévinas), je reste enfermé dans une subjectivité solipsiste et stérile. La scène de la cantine le prouve : la file d'attente n'est pas une prison, c'est l'institution fragile qui nous rend libres \emph{ensemble}. \emph{Ouverture : éduquer à la citoyenneté, c'est apprendre à transformer le conflit en reconnaissance mutuelle, le regard qui fige en visage qui oblige.}
\end{quote}
\tcblower
\textbf{Commentaire pédagogique du corrigé :}
\begin{itemize}
    \item \textbf{Ancrage} : l'introduction part du Doc 1 (bras tendu, brassard) : excellente entrée en matière.
    \item \textbf{Problématique} : bien formulée, opposant « structurellement aliénante » et « constitutivement fondatrice ».
    \item \textbf{Argumentation} : le paragraphe unique articule les trois auteurs \textbf{sans les juxtaposer} : Sartre (constat), Hegel (dépassement dialectique), Lévinas (fondation éthique). L'exemple du « brassard sans élève » est une application précise de Hegel. L'exemple de Kamel et Esther (Doc 2) illustre la liberté comme responsabilité (Lévinas) contre la liberté comme caprice.
    \item \textbf{Conclusion} : synthèse claire, réouverture sur le sens politique (citoyenneté) du cours de philosophie en série technique.
\end{itemize}
\end{exoresolu}

\courssub{Bilan de l'activité}
Cette activité d'intégration a permis de \textbf{valider les acquis} de la leçon « Autrui » par une mise en situation authentique (lycée camerounais, documents réels ou imagés). Les points clés mis en évidence sont :

\begin{enumerate}
    \item \textbf{La pluralité des modes de connaissance} : les élèves ont constaté qu'une même situation (conflit à la cantine) mobilise simultanément la perception (corps), l'analogie (je me mets à sa place), le langage (cri, affiche) et l'empathie (refusée ou naissante). Aucun mode ne suffit seul.
    \item \textbf{La dynamique de l'intersubjectivité} : le passage du \textbf{conflit} (donné immédiat, Sartre) à la \textbf{reconnaissance} (travail dialectique, Hegel) ou à la \textbf{responsabilité} (exigence éthique, Lévinas) n'est pas automatique ; il demande une \textbf{médiation} (règle, loi, intervention d'un tiers, éducation). C'est le sens de l'institution scolaire.
    \item \textbf{L'articulation théorie/pratique} : le tableau comparatif a forcé à appliquer des concepts abstraits (regard, dialectique du maître et de l'esclave, visage) à des micro-événements concrets (bras tendu, doigt pointé, silence d'Esther). C'est l'exercice type de l'\cle{épreuve du Probatoire et du Baccalauréat technique} (dissertation ou commentaire de texte).
    \item \textbf{L'écriture philosophique} : la contrainte de la mini-dissertation (introduction codifiée, paragraphe dialectique unique, conclusion ouverte) a obligé à la \textbf{rigueur syntaxique et logique} : pas de « je pense que », mais « Sartre montre que... », « Hegel démontre que... », « On peut objecter que... ».
    \item \textbf{La dimension citoyenne} : in fine, l'activité montre que la philosophie en série technique n'est pas décorative : elle outille pour \textbf{penser le vivre-ensemble} dans l'atelier, sur le chantier, dans l'entreprise — là où la reconnaissance mutuelle (sécurité, qualité, respect des normes) est condition de l'efficacité technique et de la dignité humaine.
\end{enumerate}

\begin{attention}[Pièges fréquents observés lors de la correction]
\begin{itemize}
    \item \textbf{Confondre « obstacle » et « condition »} : dire « Autrui est un obstacle MAIS aussi une condition » sans articuler \emph{comment} l'obstacle est dépassé (Hegel) ou \emph{pourquoi} la relation est fondatrice (Lévinas) conduit à une copie faible.
    \item \textbf{Citer les auteurs sans les expliquer} : « Selon Sartre, l'enfer c'est les autres » est insuffisant. Il faut dire : « Selon Sartre, le regard me constitue en objet, ce qui nie ma transcendance. »
    \item \textbf{Oublier l'exemple concret} : une dissertation sans référence aux documents 1, 2 ou 3 est hors sujet au Probatoire et au Baccalauréat technique.
    \item \textbf{Mélanger Lévinas et Hegel} : Lévinas refuse la dialectique de la reconnaissance (le visage ne se « reconnaît » pas, il m'\emph{ordonne}). Ne pas écrire « Lévinas dit qu'il faut se reconnaître mutuellement ».
    \item \textbf{Orthographe des noms propres} : on écrit \textbf{Lévinas} (avec accent), \textbf{Hegel}, \textbf{Sartre} ; les titres d'œuvres sont en italique : \textit{L'Être et le Néant}, \textit{Phénoménologie de l'esprit}, \textit{Totalité et Infini}.
\end{itemize}
\end{attention}

\newpage

\coursec{Méthodes et exercices résolus}

\courssub{Méthode 1 : Appliquer les notions de l'unité à une situation concrète}

\begin{methode}[Analyser une situation de la vie courante avec les notions de la leçon]
Pour traiter une situation concrète (récit, fait de société, scène de la vie quotidienne camerounaise) à la lumière de la leçon sur autrui, suivre les étapes suivantes :
\begin{enumerate}
\item \textbf{Lire attentivement la situation} et repérer les éléments pertinents : qui sont les personnes en présence ? Quel type de rapport les unit (regard, parole, conflit, coopération) ?
\item \textbf{Identifier la notion mobilisable} : s'agit-il d'un problème de \cle{connaissance d'autrui} (comment je sais qu'autrui pense et souffre comme moi ?) ou d'un problème d'\cle{intersubjectivité} (quel type de relation établir avec autrui ?) ?
\item \textbf{Dégager le problème philosophique} caché dans la situation : le formuler sous forme de question (ex. : « Puis-je être certain que l'autre éprouve réellement de la douleur ? »).
\item \textbf{Mobiliser les thèses des auteurs} : l'argument par \cle{analogie} (je juge autrui d'après moi-même), l'\cle{empathie} (je ressens avec autrui), la dialectique du \cle{maître et de l'esclave} chez Hegel (la reconnaissance), le \cle{regard} chez Sartre (autrui me fige en objet).
\item \textbf{Appliquer explicitement la notion au cas} : montrer en quoi la thèse éclaire la situation, avec la formule : « Ici, ... illustre ... parce que ... ».
\item \textbf{Conclure} en répondant clairement à la question posée, sans oublier les limites éventuelles de l'analyse.
\end{enumerate}
\textbf{Formules à retenir} : « Cette situation relève de la notion de ... » ; « Selon Hegel/Sartre, ... » ; « On peut donc dire que ... ».
\end{methode}

\begin{exoresolu}[La dispute au marché]
\textbf{Énoncé.} À l'entrée du marché Mokolo à Yaoundé, deux commerçantes se disputent violemment une place d'étalage. Chacune crie, accuse l'autre et finit par appeler le chef du quartier pour trancher. Un passant remarque : « Chacune veut que l'autre reconnaisse son droit. » Identifiez la notion philosophique en jeu dans cette situation et montrez, en vous appuyant sur un auteur étudié, ce qu'elle révèle du rapport à autrui.
\tcblower
\textbf{Solution détaillée.}
\begin{enumerate}
\item \textbf{Repérage.} La situation met en scène deux sujets en conflit, chacun exigeant de l'autre la reconnaissance de son droit. Il ne s'agit pas seulement d'un différend matériel (une place), mais d'une exigence de \cle{reconnaissance}.
\item \textbf{Notion mobilisée.} Le conflit relève de l'\cle{intersubjectivité}, c'est-à-dire de la question des rapports entre sujets conscients.
\item \textbf{Problème philosophique.} Pourquoi le rapport à autrui prend-il spontanément la forme d'un conflit ? Que cherche réellement chaque sujet à travers la dispute ?
\item \textbf{Thèse de l'auteur.} Selon \textbf{Hegel}, dans la dialectique du maître et de l'esclave, la conscience de soi ne peut s'affirmer qu'en étant \emph{reconnue} par une autre conscience. Chaque conscience désire que l'autre la reconnaisse comme supérieure, ce qui engendre une « lutte à mort » pour la reconnaissance.
\item \textbf{Application.} Chaque commerçante ne veut pas seulement la place : elle veut que l'autre \emph{reconnaisse} son droit, donc sa valeur et son existence de sujet. Le recours au chef du quartier montre que la reconnaissance exige un tiers et des règles communes : la reconnaissance purement conflictuelle est instable, elle appelle une médiation.
\item \textbf{Conclusion.} Cette situation illustre la thèse hégélienne : autrui est d'abord un rival, mais le conflit révèle que j'ai \emph{besoin} d'autrui pour exister comme sujet reconnu. Autrui est donc à la fois obstacle et condition de mon affirmation.
\end{enumerate}
\end{exoresolu}

\courssub{Méthode 2 : Construire et justifier une réponse argumentée (question de réflexion)}

\begin{methode}[Rédiger et justifier une réponse à une question philosophique]
Pour répondre à une question de réflexion (type courte dissertation au Probatoire/Bac technique) :
\begin{enumerate}
\item \textbf{Définir les termes} du sujet dès la première phrase (autrui, intersubjectivité, conscience, liberté...).
\item \textbf{Dégager la problématique} : transformer le sujet en question problème (ex. : « Autrui est-il un obstacle ou une condition de ma liberté ? »).
\item \textbf{Annoncer un plan en deux ou trois moments} : thèse, antithèse, dépassement (ou deux thèses confrontées).
\item \textbf{Développer chaque moment} selon le schéma : idée $\rightarrow$ argument $\rightarrow$ référence à un auteur $\rightarrow$ exemple concret.
\item \textbf{Justifier chaque affirmation} : toute thèse doit être soutenue par un raisonnement, jamais simplement affirmée.
\item \textbf{Conclure} par une réponse nuancée et personnelle à la problématique.
\end{enumerate}
\textbf{Réflexes} : une idée par paragraphe ; des connecteurs logiques (« d'abord », « cependant », « en effet », « par conséquent ») ; jamais de nom d'auteur sans sa thèse.
\end{methode}

\begin{exoresolu}[Puis-je connaître autrui comme je me connais moi-même ?]
\textbf{Énoncé.} « Je ne peux jamais savoir ce que l'autre ressent vraiment. » Discutez cette affirmation en un développement structuré d'une vingtaine de lignes.
\tcblower
\textbf{Solution détaillée (corrigé rédigé).}

\emph{Introduction.} Autrui désigne l'autre homme, celui qui, comme moi, est un sujet pensant et sentant. Mais si je me connais moi-même de l'intérieur, par conscience immédiate, je n'accède à autrui que de l'extérieur, par son corps et ses conduites. Dès lors, la connaissance d'autrui est-elle possible, ou autrui demeure-t-il à jamais une énigme ?

\emph{Premier moment : le doute.} Il semble d'abord que je ne puisse jamais éprouver directement la pensée d'autrui. Je vois un homme grimacer et crier : je perçois des signes, non la douleur elle-même. Rien ne me garantit que derrière ces gestes il y a une conscience semblable à la mienne : un automate perfectionné produirait les mêmes signes. La connaissance d'autrui paraît donc toujours incertaine, indirecte.

\emph{Deuxième moment : les voies d'accès à autrui.} Cependant, cette incertitude n'est pas une ignorance totale. Par le raisonnement par \cle{analogie}, je juge autrui d'après moi-même : ayant moi-même mal lorsque je me blesse, je conclus que l'autre souffre lorsqu'il présente les mêmes signes. Par l'\cle{empathie}, je ressens spontanément avec autrui : la vue d'un enfant qui pleure éveille immédiatement en moi sa tristesse. Enfin, par le \cle{langage}, autrui me dit ce qu'il pense : la parole crée entre nous une communication des esprits. Ces trois voies ne me donnent pas la certitude absolue, mais une connaissance suffisante et vivante.

\emph{Conclusion.} Je ne connais donc pas autrui comme je me connais moi-même : sa conscience reste pour moi invisible. Mais l'analogie, l'empathie et le langage me relient à lui de manière réelle, quoique imparfaite. Autrui n'est ni transparent ni totalement opaque : il est ce proche que je ne cesse jamais tout à fait de découvrir.
\end{exoresolu}

\courssub{Méthode 3 : Distinguer les modes de connaissance d'autrui}

\begin{methode}[Identifier le mode de connaissance mobilisé dans un cas donné]
Trois modes de connaissance d'autrui sont au programme. Pour les distinguer et les reconnaître :
\begin{enumerate}
\item \textbf{La perception et l'analogie} : j'observe le corps et les comportements d'autrui, et je \emph{déduis} qu'il a une conscience semblable à la mienne. Réflexe : repérer dans l'énoncé des mots comme « voir », « observer », « ressembler », « en conclure ».
\item \textbf{L'empathie} : je \emph{ressens} directement avec autrui, sans raisonnement. Réflexe : repérer une émotion contagieuse, immédiate (rire, pleurs, souffrance partagée).
\item \textbf{Le langage et la communication} : autrui \emph{exprime} sa pensée par la parole, et je le comprends. Réflexe : repérer un dialogue, un récit, un aveu.
\end{enumerate}
\textbf{Formule à retenir} : « Ce cas relève de ... car le sujet accède à autrui par ... ». Toujours préciser la limite du mode identifié (l'analogie reste incertaine ; l'empathie peut se tromper ; la parole peut mentir).
\end{methode}

\begin{exoresolu}[Trois scènes, trois modes de connaissance]
\textbf{Énoncé.} Pour chacune des scènes suivantes, indiquez le mode de connaissance d'autrui en jeu et justifiez votre réponse :
(a) En voyant son voisin boiter après un match de foot, Ali conclut qu'il s'est fait mal.
(b) À l'annonce du décès d'un proche, la mère de Sandra éclate en sanglots et Sandra se met elle-même à pleurer.
(c) Lors d'une palabre familiale à Bafoussam, l'oncle explique longuement les raisons de sa décision, et chacun comprend son point de vue.
\tcblower
\textbf{Solution détaillée.}
\begin{enumerate}
\item \textbf{Scène (a) : la perception et l'analogie.} Ali \emph{observe} un comportement visible (boiter) et \emph{conclut} à un état intérieur invisible (la douleur), en se fondant sur sa propre expérience : « quand je boite, c'est que j'ai mal ». C'est le schéma typique du raisonnement par analogie : je juge autrui d'après moi-même. Limite : cette conclusion reste probable, non certaine (le voisin pourrait simuler).
\item \textbf{Scène (b) : l'empathie.} Sandra ne raisonne pas : elle \emph{ressent} immédiatement la tristesse de sa mère, au point de pleurer elle-même. L'émotion d'autrui se communique à elle directement, sans médiation du raisonnement. Limite : l'empathie peut être trompeuse ou excessive (on peut pleurer sans comprendre les raisons de l'autre).
\item \textbf{Scène (c) : le langage et la communication.} L'oncle \emph{exprime} ses raisons par la parole ; les auditeurs accèdent à sa pensée parce qu'il la dit. Le langage donne accès à des contenus de conscience (motifs, justifications) que ni la perception ni l'empathie ne peuvent atteindre. Limite : la parole peut dissimuler ou mentir.
\item \textbf{Conclusion.} Les trois scènes montrent que la connaissance d'autrui emprunte des voies complémentaires : observer, ressentir, écouter. Aucune n'est infaillible, mais leur combinaison rend la vie commune possible.
\end{enumerate}
\end{exoresolu}

\courssub{Méthode 4 : Mobiliser Hegel et Sartre sur le rapport à autrui}

\begin{methode}[Utiliser correctement les références philosophiques]
Pour exploiter les deux auteurs au programme sans les confondre :
\begin{enumerate}
\item \textbf{Hegel — la reconnaissance.} Thèse : la conscience de soi a besoin d'être reconnue par une autre conscience ; d'où la lutte pour la reconnaissance (dialectique du maître et de l'esclave). Idée-clé : autrui est \emph{condition} de ma conscience de moi.
\item \textbf{Sartre — le regard.} Thèse : sous le regard d'autrui, je deviens un objet (« l'enfer, c'est les autres ») ; le conflit est au cœur des relations. Idée-clé : autrui est d'abord \emph{obstacle} à ma liberté, car il me fige.
\item \textbf{Choisir l'auteur adapté} à la question : conflit et aliénation $\rightarrow$ Sartre ; reconnaissance et besoin d'autrui $\rightarrow$ Hegel ; sujet dialectique (obstacle \emph{et} condition) $\rightarrow$ les deux confrontés.
\item \textbf{Citer brièvement} : nom de l'auteur + sa thèse en une phrase + application à l'exemple. Ne jamais écrire une biographie.
\end{enumerate}
\end{methode}

\begin{exoresolu}[Autrui : obstacle ou condition de ma liberté ?]
\textbf{Énoncé.} Un élève déclare : « Les autres m'empêchent d'être libre ; je serais plus libre seul. » En confrontant les thèses de Hegel et de Sartre, montrez ce que cette opinion a de vrai et ce qu'elle oublie.
\tcblower
\textbf{Solution détaillée.}
\begin{enumerate}
\item \textbf{Définitions.} Autrui est l'autre sujet conscient ; la liberté est le pouvoir de se déterminer soi-même. Le problème : la présence d'autrui limite-t-elle ma liberté, ou la rend-elle possible ?
\item \textbf{Ce que l'opinion a de vrai — Sartre.} Selon Sartre, le regard d'autrui me saisit comme un objet : sous le regard du professeur, je cesse d'être pure liberté pour devenir « l'élève dissipé », « le tricheur ». Autrui me juge, me fige, m'aliène. En ce sens, l'élève a raison : autrui est un obstacle, car il m'enferme dans une image dont je ne suis pas maître.
\item \textbf{Ce que l'opinion oublie — Hegel.} Mais selon Hegel, sans autrui, je ne pourrais même pas me savoir libre. La conscience de soi n'existe que reconnue par une autre conscience : c'est dans la lutte pour la reconnaissance (maître et esclave) que le sujet prend conscience de sa valeur. Un homme absolument seul n'aurait ni langage, ni estime de soi, ni conscience d'être libre.
\item \textbf{Dépassement.} Autrui est donc dialectiquement les deux : il limite ma liberté par son regard (Sartre), mais il la constitue par la reconnaissance (Hegel). Ma liberté ne s'exerce jamais dans le vide : elle se joue \emph{parmi} les autres.
\item \textbf{Conclusion.} L'opinion de l'élève est donc à moitié vraie : autrui contraint ma liberté, mais il en est aussi la condition. Être libre, ce n'est pas être seul, c'est affirmer sa liberté dans le rapport aux autres.
\end{enumerate}
\end{exoresolu}

\courssub{Méthode 5 : Définir et illustrer l'intersubjectivité}

\begin{methode}[Traiter une question sur l'intersubjectivité]
\begin{enumerate}
\item \textbf{Définir} : l'\cle{intersubjectivité} désigne l'ensemble des relations entre sujets conscients, c'est-à-dire le fait que des consciences se rencontrent, communiquent et se reconnaissent mutuellement comme sujets.
\item \textbf{Distinguer les formes} du rapport à autrui : la \emph{communication} (échange de paroles), la \emph{coopération} (travail commun), l'\emph{amour et l'amitié} (union des personnes), mais aussi le \emph{conflit} (lutte, rivalité).
\item \textbf{Montrer l'enjeu} : l'intersubjectivité fonde la vie sociale, la morale (le respect d'autrui comme personne) et le langage.
\item \textbf{Illustrer par un exemple camerounais} : la palabre villageoise, le travail en association (tontine, travaux communautaires), le marché comme lieu d'échanges.
\item \textbf{Conclure} sur la double face du rapport à autrui : union et conflit.
\end{enumerate}
\end{methode}

\begin{exoresolu}[La tontine du quartier]
\textbf{Énoncé.} Dans un quartier de Douala, douze femmes se réunissent chaque mois pour une tontine : chacune cotise, et le total revient à tour de rôle à l'une d'elles. Les décisions se prennent après discussion, et les conflits se règlent par la parole. Montrez que cette pratique illustre l'intersubjectivité et précisez les formes du rapport à autrui qu'elle met en œuvre.
\tcblower
\textbf{Solution détaillée.}
\begin{enumerate}
\item \textbf{Rappel de la définition.} L'intersubjectivité est la relation entre sujets conscients qui se reconnaissent mutuellement comme sujets et construisent ensemble un monde commun.
\item \textbf{La coopération.} La tontine repose sur un projet commun : chacune cotise en confiance pour le bénéfice de toutes. Aucune ne pourrait atteindre seule ce résultat (réunir une somme importante) : autrui apparaît ici comme \emph{condition} de la réussite de chacune.
\item \textbf{La communication.} Les décisions « se prennent après discussion » : le langage est le lien intersubjectif par excellence. Par la parole, les sujets échangent leurs points de vue et construisent une volonté commune.
\item \textbf{La reconnaissance et la confiance.} Cotiser suppose de reconnaître les autres comme des personnes dignes de confiance, qui tiendront leur engagement. Chacune est traitée comme un sujet responsable, non comme une chose.
\item \textbf{Le conflit et sa régulation.} Même unie, la communauté connaît des tensions ; mais elles se règlent « par la parole », c'est-à-dire par une médiation intersubjective et non par la violence.
\item \textbf{Conclusion.} La tontine illustre les trois grandes formes positives de l'intersubjectivité : coopération, communication, reconnaissance. Elle montre que le rapport à autrui, loin d'être seulement un conflit, est le fondement de la vie sociale et de l'entraide.
\end{enumerate}
\end{exoresolu}

\courssub{Méthode 6 : Rédiger une introduction de dissertation sur autrui}

\begin{methode}[Construire une introduction complète en quatre temps]
\begin{enumerate}
\item \textbf{L'amorce} : partir d'un constat général et vrai (ex. : « L'homme ne vit jamais seul : dès la naissance, il est entouré d'autres hommes. »).
\item \textbf{La définition des termes} : définir brièvement le ou les mots-clés du sujet (autrui, connaissance, intersubjectivité, liberté).
\item \textbf{La problématique} : dégager la difficulté et la formuler en une question (ou deux questions liées).
\item \textbf{L'annonce du plan} : présenter en une phrase les deux ou trois moments du développement.
\end{enumerate}
\textbf{Réflexes} : introduction courte (5 à 8 lignes) ; interdiction de recopier le sujet tel quel ; la problématique doit être une vraie question, pas une affirmation.
\end{methode}

\begin{exoresolu}[Rédiger l'introduction d'un sujet de dissertation]
\textbf{Énoncé.} Sujet : « La présence d'autrui est-elle nécessaire à la vie humaine ? » Rédigez une introduction complète à ce sujet.
\tcblower
\textbf{Solution détaillée (corrigé rédigé).}

\emph{Amorce.} Nul homme ne vient au monde ni ne grandit seul : la famille, le village, l'école, le marché sont autant de lieux où l'existence humaine se déploie parmi les autres.

\emph{Définitions.} Autrui désigne l'autre homme, sujet conscient comme moi ; la vie humaine désigne non la simple survie biologique, mais l'existence proprement humaine, faite de langage, de conscience et de relations.

\emph{Problématique.} On pourrait croire qu'un homme pourrait vivre seul, à l'écart de tous, et que les autres ne sont qu'une contrainte. Mais sans autrui, qui m'apprendrait à parler, qui me reconnaîtrait, qui me donnerait conscience de moi-même ? La présence d'autrui est-elle alors un simple accident de l'existence, ou la condition même de l'humanité de l'homme ?

\emph{Annonce du plan.} Nous verrons d'abord que la tentation de la solitude semble rendre autrui superflu ; nous montrerons ensuite que la connaissance de soi, le langage et la reconnaissance exigent autrui ; nous conclurons enfin que vivre humainement, c'est vivre avec et pour les autres.

\emph{Justification de la démarche.} L'introduction suit les quatre temps exigés : amorce concrète, définitions précises, problématique formulée en question, plan dialectique en trois moments (thèse, antithèse, synthèse) clairement annoncé.
\end{exoresolu}

\begin{attention}[Les 5 erreurs qui coûtent des points au Bac]
\begin{enumerate}
\item \textbf{Confondre « autrui » et « les autres » au sens banal.} Autrui n'est pas n'importe quelle personne : c'est l'autre \emph{sujet conscient}, semblable à moi. Le définir simplement comme « mon prochain » ou « mon voisin » sans préciser qu'il est une conscience est une définition insuffisante.
\item \textbf{Citer un auteur sans sa thèse.} Écrire « selon Hegel » ou « comme dit Sartre » sans exposer la pensée de l'auteur (lutte pour la reconnaissance ; le regard qui me fige) ne vaut aucun point. Toute référence doit être expliquée en une phrase et reliée à l'argument.
\item \textbf{Affirmer sans justifier.} Dire « autrui est nécessaire » ou « autrui est un obstacle » sans argument ni exemple est une opinion, pas une démonstration. Chaque thèse doit suivre le schéma : idée $\rightarrow$ argument $\rightarrow$ auteur $\rightarrow$ exemple.
\item \textbf{Confondre les modes de connaissance.} Ne pas mélanger analogie (je \emph{déduis} la pensée d'autrui de ses comportements) et empathie (je \emph{ressens} avec autrui sans raisonnement). Attribuer à l'empathie un raisonnement, ou à l'analogie une émotion immédiate, est une erreur de notion.
\item \textbf{Oublier la problématique et la conclusion.} Beaucoup de copies exposent des notions sans répondre à la question posée. La problématique doit apparaître dans l'introduction sous forme de question, et la conclusion doit y répondre explicitement et de manière nuancée (obstacle \emph{et} condition, connaissance réelle \emph{mais} imparfaite).
\end{enumerate}
\end{attention}

\newpage

\coursec{Exercices}

\courssub{Je vérifie mes connaissances}

\begin{exercice}[QCM : Les modes de connaissance d'autrui]
Parmi les propositions suivantes, une seule est exacte. Recopiez la lettre correspondante et justifiez brièvement votre choix.
\begin{enumerate}
    \item La connaissance d'autrui par \emph{l'analogie} consiste à :
    \begin{itemize}
        \item[A.] Ressentir immédiatement la douleur de l'autre comme la mienne.
        \item[B.] Inférer que l'autre a des états mentaux semblables aux miens parce qu'il a un corps et des comportements semblables aux miens.
        \item[C.] Comprendre l'autre uniquement par le langage verbal.
        \item[D.] Nier l'existence de la conscience d'autrui.
    \end{itemize}
    \item Pour \textbf{Husserl}, la réduction phénoménologique appliquée à autrui permet de mettre en évidence :
    \begin{itemize}
        \item[A.] Que l'autre est un objet parmi d'autres dans le monde.
        \item[B.] La constitution de l'autre comme \emph{alter ego} par l'\emph{appresentation} (accouplement ou \emph{Paarung}).
        \item[C.] Que l'empathie est une illusion sentimentale.
        \item[D.] La supériorité de la connaissance scientifique sur l'intuition.
    \end{itemize}
    \item Le \emph{langage} comme mode de connaissance d'autrui se distingue de la perception parce qu'il :
    \begin{itemize}
        \item[A.] Ne permet pas d'accéder à la subjectivité de l'autre.
        \item[B.] Repose sur une convention et une intentionnalité communicative que la simple observation du corps ne donne pas.
        \item[C.] Est une forme d'analogie déguisée.
        \item[D.] Supprime toute ambiguïté dans la compréhension d'autrui.
    \end{itemize}
    \item L'\emph{intersubjectivité} désigne fondamentalement :
    \begin{itemize}
        \item[A.] Le conflit permanent entre les consciences.
        \item[B.] La communauté de sens et de monde partagée entre plusieurs sujets.
        \item[C.] La connaissance objective des sciences humaines.
        \item[D.] La solitude ontologique de l'individu.
    \end{itemize}
\end{enumerate}
\end{exercice}

\begin{exercice}[Vrai / Faux justifié : Intersubjectivité et liberté]
Pour chaque affirmation, indiquez VRAI ou FAUX et justifiez en une phrase en mobilisant un auteur ou une notion du cours.
\begin{enumerate}
    \item L'intersubjectivité se réduit à la simple coexistence physique d'individus dans un même espace.
    \item Pour \textbf{Hegel}, la reconnaissance mutuelle dans la dialectique du Maître et de l'Esclave est la condition nécessaire de la conscience de soi et de la liberté.
    \item Selon \textbf{Sartre}, dans \emph{L'Être et le Néant}, le regard d'autrui me constitue comme une liberté absolue en me dévoilant à moi-même.
    \item L'empathie (\emph{Einfühlung}) chez \textbf{Max Scheler} ou \textbf{Edith Stein} est une perception originaire de la vie psychique d'autrui, et non une inférence analogique.
\end{enumerate}
\end{exercice}

\begin{exercice}[Définitions précises]
Définissez rigoureusement les quatre notions suivantes en philosophie (2 à 3 lignes chacune, vocabulaire technique attendu) :
\begin{enumerate}
    \item L'analogie (comme argument pour l'existence d'autrui).
    \item L'appresentation (ou \emph{Paarung}) chez Husserl.
    \item La dialectique de la reconnaissance (Hegel).
    \item L'aliénation par le regard (Sartre).
\end{enumerate}
\end{exercice}

\begin{exercice}[Mise en relation : situations et notions]
Reliez chaque situation concrète au mode de connaissance d'autrui ou à la forme d'intersubjectivité qu'elle illustre le mieux. Justifiez le lien en une phrase.
\begin{center}
\begin{tabularx}{\linewidth}{|X|l|}\hline
\rowcolor{popblueL}\textbf{Situation} & \textbf{Notion / Auteur} \\ \hline
1. Je vois un enfant pleurer au marché Mokolo à Yaoundé ; je ressens immédiatement sa tristesse sans raisonner. & A. Analogie (Mill, Descartes) \\ \hline
2. Mon collègue me dit : « Je suis épuisé par ce chantier » ; je comprends sa fatigue par ses mots. & B. Empathie (Scheler, Stein) \\ \hline
3. Je sais que mon frère a mal aux dents parce qu'il se tient la joue et grimace, comme moi quand j'ai mal. & C. Langage / Communication (Ricœur) \\ \hline
4. Deux associations de quartiers (Bassa et Akwa) négocient un accord de gestion des déchets communs. & D. Intersubjectivité éthique / Reconnaissance (Hegel, Ricœur) \\ \hline
\end{tabularx}
\end{center}
\end{exercice}

\courssub{Je m'entraîne}

\begin{exercice}[Analyse d'une situation : La solidarité face aux inondations à Douala]
\begin{quote}
« Après de violentes pluies, le quartier de New Bell, à Douala, est inondé. Des habitants, sans concertation préalable, forment spontanément des chaînes humaines pour évacuer les eaux des maisons voisines, partagent leur nourriture avec ceux qui ont tout perdu, et organisent le nettoyage commun des rues. »
\end{quote}
\begin{enumerate}
    \item Identifiez la forme d'intersubjectivité manifestée ici (primordiale, éthique, communicative, conflictuelle ?). Justifiez.
    \item Expliquez en quoi cette situation illustre la thèse de \textbf{Paul Ricœur} sur la \emph{promesse} et la confiance comme fondements du lien social.
    \item Montrez comment l'action collective transforme la vulnérabilité individuelle en puissance commune (vous pouvez mobiliser la reconnaissance chez \textbf{Hegel} ou la solidarité chez \textbf{Durkheim}).
\end{enumerate}
\end{exercice}

\begin{exercice}[Commentaire d'argument : L'analogie et ses limites]
\begin{quote}
« Je ne perçois pas la douleur de l'autre, je perçois son corps qui se tord, son visage qui se contracte. Par analogie avec mon propre cas (quand mon corps se tord, j'ai mal), j'infère qu'il a mal. Mais cette inférence est-elle certaine ? »
\end{quote}
\begin{enumerate}
    \item Reformulez l'argument de l'analogie (structure : prémisses / conclusion).
    \item Identifiez deux objections majeures à cet argument (risque de solipsisme, unicité du cas de départ, asymétrie entre ce que je vis et ce que j'observe).
    \item En quoi la phénoménologie (Husserl, Scheler) tente-t-elle de dépasser l'analogie par les concepts d'\emph{appresentation} ou de \emph{perception originaire} ?
\end{enumerate}
\end{exercice}

\begin{exercice}[Débat contradictoire : Autrui, enfer ou salut ?]
Le sujet : \og \textbf{Autrui est-il un obstacle ou une condition de ma liberté ?} \fg
\begin{enumerate}
    \item Rédigez une thèse (Autrui comme obstacle/aliénation) en vous appuyant sur \textbf{Sartre} (\emph{le regard}, \emph{l'enfer, c'est les autres}) et sur l'expérience de la honte.
    \item Rédigez une antithèse (Autrui comme condition/révélation) en vous appuyant sur \textbf{Hegel} (\emph{reconnaissance}, dialectique Maître/Esclave) et/ou \textbf{Levinas} (\emph{visage}, responsabilité).
    \item Proposez une synthèse dialectique qui articule conflit et nécessité de l'autre pour l'accès à l'humanité.
\end{enumerate}
\end{exercice}

\begin{exercice}[Méthode : Construire une introduction de dissertation]
Sujet : \og \textbf{Peut-on connaître autrui ?} \fg
Rédigez une introduction complète (accroche, définition des termes, analyse du sujet, problématique, annonce du plan en trois parties) en respectant la méthode officielle (MINESEC). Votre plan doit mobiliser : 1. Les voies d'accès (perception, analogie, empathie, langage) ; 2. Les limites et l'opacité radicale ; 3. La connaissance par l'action commune et la reconnaissance (intersubjectivité).
\end{exercice}

\begin{methode}[Rappel : les étapes de l'introduction au Probatoire / Bac technique]
\begin{enumerate}
    \item \textbf{Accroche} : fait ou constat qui amène naturellement au sujet.
    \item \textbf{Définition des termes} : « connaître » (saisir par l'esprit, atteindre la vérité d'une chose) ; « autrui » (l'autre conscience, l'autre que moi, semblable et pourtant distinct).
    \item \textbf{Analyse du sujet} : dégager le présupposé et la difficulté (autrui est un sujet, non un objet ; peut-on connaître un sujet comme on connaît une chose ?).
    \item \textbf{Problématique} : une question principale, éventuellement déclinée.
    \item \textbf{Annonce du plan} : trois parties clairement formulées.
\end{enumerate}
\end{methode}

\courssub{Je me prépare au Bac}

\begin{exercice}[Dissertation type Baccalauréat Technique — Sujet 1]
\textbf{Sujet :} \og \textbf{Autrui est-il un obstacle à ma liberté ?} \fg

\begin{enumerate}
    \item Rédigez l'introduction complète (accroche, définitions, analyse, problématique, plan annoncé).
    \item Développez la \textbf{Première partie} (Thèse : Autrui comme limitation/aliénation de ma liberté). Deux arguments minimum avec exemples (Sartre, vie courante : surveillance, pression sociale, regard objectivant).
    \item Développez la \textbf{Deuxième partie} (Antithèse : Autrui comme condition de ma liberté). Deux arguments minimum (Hegel : reconnaissance ; Levinas : responsabilité qui fonde ma subjectivité éthique).
    \item Développez la \textbf{Troisième partie} (Synthèse : une liberté \emph{avec} et \emph{par} autrui). Articulation : la liberté concrète s'exerce dans l'intersubjectivité (Ricœur, vie démocratique).
    \item Rédigez la \textbf{Conclusion} (réponse synthétique, ouverture sur la citoyenneté ou les droits de la personne au Cameroun).
\end{enumerate}
\end{exercice}

\begin{exercice}[Explication de texte guidée — Sujet 2]
\textbf{Texte :} Extrait adapté de \emph{L'Être et le Néant}, Jean-Paul Sartre (le regard).
\begin{quote}
« C'est donc par le regard d'autrui que je deviens un objet. Mais cet objet, je ne le suis que pour autrui. [...] Je suis vu, cela veut dire : je perds ma maîtrise du monde. [...] L'enfer, c'est les autres. »
\end{quote}
\begin{enumerate}
    \item \textbf{Identification :} Auteur, ouvrage, thèse principale du texte en une phrase.
    \item \textbf{Analyse structurée :} Dégagez les trois moments de l'argumentation (1. L'objectivation par le regard ; 2. La perte de maîtrise / aliénation ; 3. La portée existentielle : l'enfer).
    \item \textbf{Explication de concepts :} Expliquez dans le contexte du texte : \emph{être-pour-autrui}, \emph{regard}, \emph{honte}, \emph{liberté}.
    \item \textbf{Discussion critique :} En quoi cette thèse s'oppose-t-elle à la dialectique hégélienne de la reconnaissance ? La conflictualité est-elle la seule vérité de la relation à autrui ? (Apportez des nuances : empathie, amour, travail commun.)
\end{enumerate}
\end{exercice}

\begin{exercice}[Question de synthèse et analyse de situation — Sujet 3]
\textbf{Sujet :} \og \textbf{Peut-on connaître autrui ? Mobilisez la perception, l'analogie, l'empathie, le langage et l'intersubjectivité.} \fg

\begin{enumerate}
    \item \textbf{Synthèse théorique (une page) :} Structurez votre réponse en trois temps :
    \begin{itemize}
        \item Les voies d'accès (analogie : inférence ; empathie : perception originaire ; langage : médiation symbolique).
        \item Les limites de chaque voie (asymétrie, projection, malentendu, opacité de la conscience).
        \item La connaissance par l'intersubjectivité agie (travail commun, reconnaissance : connaître en agissant \emph{avec}).
    \end{itemize}
    \item \textbf{Application camerounaise :} À partir de la situation suivante, analysez les modes de connaissance en jeu et les obstacles :
    \begin{quote}
    « Lors d'un conflit foncier entre deux villages de la région de l'Ouest, les chefs traditionnels et les notables s'assoient sous l'arbre à palabres. Ils ne se contentent pas d'échanger des arguments juridiques ; ils partagent la noix de cola, invoquent les ancêtres communs, rappellent les alliances matrimoniales passées. Peu à peu, la tension baisse, un compromis émerge. »
    \end{quote}
    \begin{itemize}
        \item Identifiez les formes d'intersubjectivité mobilisées (communicative, éthique, symbolique/rituelle).
        \item Expliquez pourquoi le \emph{langage} juridique seul (procès formel) n'aurait pas suffi, au regard de la philosophie de \textbf{Ricœur} (confiance, parole donnée).
        \item En quoi la \emph{reconnaissance mutuelle} (Hegel) est-elle la condition de la résolution pacifique ?
    \end{itemize}
\end{enumerate}
\end{exercice}



\newpage

\coursec{Corrigés des exercices}

\courssub{Je vérifie mes connaissances}

\begin{corrige}[Exercice 1 — QCM : Les modes de connaissance d'autrui]
\begin{enumerate}
    \item \textbf{B.} L'analogie est un raisonnement par ressemblance : j'observe le corps et les conduites d'autrui, semblables aux miens, et j'en infère qu'il éprouve, comme moi, des états de conscience (douleur, joie, pensée). C'est une \emph{inférence}, non une perception directe.
    \item \textbf{B.} Chez Husserl (\emph{Méditations cartésiennes}, V\textsuperscript{e} méditation), autrui se constitue pour moi comme \emph{alter ego} : son corps m'est perçu, mais sa vie psychique m'est seulement \emph{appresentée} (co-donnée, jamais donnée en propre) par accouplement avec mon propre corps.
    \item \textbf{B.} Le langage repose sur des signes conventionnels et sur l'intention de communiquer : autrui peut me dire ce qu'il pense et ressent, ce que la seule observation de son corps ne révèle pas. Mais il n'élimine pas le mensonge ni le malentendu (d'où le rejet de D).
    \item \textbf{B.} L'intersubjectivité est la relation entre sujets qui partagent un monde commun de significations ; elle ne se réduit ni au conflit (A) ni à la coexistence d'objets (C, D).
\end{enumerate}
\end{corrige}

\begin{corrige}[Exercice 2 — Vrai / Faux justifié : Intersubjectivité et liberté]
\begin{enumerate}
    \item \textbf{FAUX.} L'intersubjectivité suppose un partage de sens, d'intentions et de monde entre consciences ; des individus physiquement proches peuvent demeurer étrangers les uns aux autres (simple juxtaposition sans communication).
    \item \textbf{VRAI.} Chez Hegel (\emph{Phénoménologie de l'Esprit}), la conscience de soi n'accède à la certitude d'elle-même que par la reconnaissance d'une autre conscience : « la conscience de soi existe en et par soi, en cela que et parce qu'elle existe pour une autre conscience de soi ».
    \item \textbf{FAUX.} Pour Sartre, le regard d'autrui m'\emph{objective} : il me fige en « être-pour-autrui », me vole mon monde et suscite la honte ; il aliène ma liberté au lieu de la constituer comme absolue.
    \item \textbf{VRAI.} Scheler et Stein soutiennent que l'empathie est un acte originaire de saisie du vécu d'autrui (je perçois la joie \emph{dans} le sourire), distinct du raisonnement par analogie défendu par Mill.
\end{enumerate}
\end{corrige}

\begin{corrige}[Exercice 3 — Définitions précises]
\begin{enumerate}
    \item \textbf{L'analogie} : raisonnement par lequel, constatant que le corps et les comportements d'autrui ressemblent aux miens, et sachant que mes comportements sont liés à des états de conscience, je conclus par ressemblance qu'autrui possède lui aussi une conscience. C'est une inférence probable, non une évidence.
    \item \textbf{L'appresentation} : chez Husserl, mode de donation indirecte par lequel la vie psychique d'autrui est co-visée avec la perception de son corps, sans jamais être donnée en propre ; par accouplement (\emph{Paarung}) de son corps au mien, autrui se constitue comme \emph{alter ego}.
    \item \textbf{La dialectique de la reconnaissance} : chez Hegel, processus par lequel deux consciences s'affrontent (lutte à mort, rapport Maître/Esclave) pour être reconnues ; la conscience de soi et la liberté ne s'accomplissent que dans la reconnaissance \emph{mutuelle} et égale.
    \item \textbf{L'aliénation par le regard} : chez Sartre, le regard d'autrui me saisit comme objet (« être-pour-autrui »), me fige dans une identité extérieure, me dépossède de mon monde et de ma transcendance ; la honte est la conscience de cette objectivation.
\end{enumerate}
\end{corrige}

\begin{corrige}[Exercice 4 — Mise en relation : situations et notions]
\begin{itemize}
    \item \textbf{1 -- B (Empathie).} La tristesse est saisie immédiatement, sans raisonnement inférentiel : c'est la perception originaire du vécu d'autrui décrite par Scheler et Stein.
    \item \textbf{2 -- C (Langage).} La compréhension passe par des signes conventionnels porteurs d'une intention de communiquer : autrui me dit son vécu.
    \item \textbf{3 -- A (Analogie).} Je raisonne par ressemblance entre son comportement (grimace, main sur la joue) et le mien dans la même situation, et j'infère sa douleur.
    \item \textbf{4 -- D (Intersubjectivité éthique).} La négociation suppose la reconnaissance mutuelle des parties comme interlocuteurs légitimes et la construction d'un monde commun (gestion partagée de l'espace).
\end{itemize}
\end{corrige}

\courssub{Je m'entraîne}

\begin{corrige}[Exercice 5 — Analyse d'une situation : La solidarité face aux inondations à Douala]
\begin{enumerate}
    \item Il s'agit d'une intersubjectivité \textbf{éthique} (et secondairement communicative) : les habitants se reconnaissent mutuellement comme personnes vulnérables et responsables les uns des autres ; l'entraide spontanée manifeste un monde commun de valeurs (solidarité, entraide de quartier), non un simple voisinage physique.
    \item Chez Ricœur, le lien social repose sur la parole donnée et la confiance : ici, sans contrat écrit, chacun compte sur l'engagement implicite des autres (tenir la chaîne, partager la nourriture). La promesse, même tacite, institue une relation durable et fiable entre des personnes qui se reconnaissent.
    \item Isolé, chaque habitant est impuissant face à la montée des eaux ; unis dans l'action, ils produisent une puissance qu'aucun ne possède seul. Chez Hegel, cette coopération est une forme de reconnaissance : chacun se découvre sujet agissant à travers l'action commune. Chez Ricoeur, l'agir ensemble (le « faire ensemble ») construit une identité narrative partagée ; la reconnaissance mutuelle dans l'action commune transforme la vulnérabilité individuelle en force collective.
\end{enumerate}
\end{corrige}

\begin{corrige}[Exercice 6 — Commentaire d'argument : L'analogie et ses limites]
\begin{enumerate}
    \item \textbf{Structure de l'argument.} Prémisse 1 : quand mon corps se tord et grimace, j'éprouve de la douleur. Prémisse 2 : le corps d'autrui se tord et grimace comme le mien. Conclusion : autrui éprouve probablement de la douleur. C'est un raisonnement par ressemblance, du connu (mon cas) vers l'inconnu (le cas d'autrui).
    \item \textbf{Objections.} (a) L'inférence repose sur un cas unique (le mien) : généraliser à partir d'un seul exemple est logiquement fragile. (b) L'asymétrie est irréductible : je vis ma douleur de l'intérieur, je ne fais qu'observer les signes extérieurs de l'autre ; rien ne garantit qu'ils soient accompagnés d'une conscience (on peut simuler). Dès lors, le doute solipsiste (« et si j'étais seul à être conscient ? ») n'est pas logiquement réfuté par l'analogie.
    \item \textbf{Dépassement phénoménologique.} Husserl montre que la conscience d'autrui n'est pas inférée mais \emph{appresentée} : percevant son corps comme vivant, je co-vise immédiatement sa subjectivité, sans raisonnement. Scheler va plus loin : dans l'empathie, je perçois originairement la joie \emph{dans} le sourire, la douleur \emph{dans} les larmes ; l'expression n'est pas un indice d'un invisible, mais la présence même du vécu d'autrui.
\end{enumerate}
\end{corrige}

\begin{corrige}[Exercice 7 — Débat contradictoire : Autrui, enfer ou salut ?]
\begin{enumerate}
    \item \textbf{Thèse.} Pour Sartre, autrui est d'abord celui qui me regarde et me fige : sous son regard, je deviens objet, « être-pour-autrui » ; il me vole mon monde et mes possibles. La honte révèle cette objectivation : je me découvre tel que l'autre me voit, prisonnier d'une identité extérieure. D'où la formule de \emph{Huis clos} : « l'enfer, c'est les autres » — autrui limite et aliène ma liberté.
    \item \textbf{Antithèse.} Pour Hegel, sans autrui je ne suis même pas une conscience de soi : je n'accède à la certitude de moi-même que si une autre conscience me reconnaît ; la liberté s'accomplit dans la reconnaissance mutuelle. Pour Levinas, le visage d'autrui m'appelle à la responsabilité : loin de m'aliéner, il fonde ma subjectivité éthique et m'arrache à l'égoïsme.
    \item \textbf{Synthèse.} Le conflit décrit par Sartre est réel, mais il n'est pas le dernier mot : il est lui-même le signe que j'ai besoin d'autrui (on ne lutte que pour être reconnu). La liberté concrète n'est ni la solitude ni la fusion : elle s'exerce \emph{avec} autrui, dans des relations de reconnaissance réciproque, de parole et d'action commune (Ricœur). Autrui est à la fois le risque de l'aliénation et la condition de mon humanité.
\end{enumerate}
\end{corrige}

\begin{corrige}[Exercice 8 — Méthode : Construire une introduction de dissertation]
Nous vivons entourés d'autres hommes : nous leur parlons, travaillons avec eux, partageons leurs joies et leurs peines. Pourtant, je n'ai jamais accès directement à ce qu'ils pensent ou éprouvent : je ne perçois que des corps, des gestes, des paroles. Connaître, c'est saisir par l'esprit la nature ou la vérité d'une chose ; autrui désigne l'autre conscience, cet autre moi-même qui m'est semblable mais dont l'intériorité m'est fermée. Le sujet présuppose donc une difficulté : autrui n'est pas un objet parmi les objets, mais un sujet ; or connaître un sujet, est-ce possible de la même manière qu'on connaît une chose ? Dès lors se pose la question : \textbf{peut-on connaître autrui ?} Autrement dit : par quelles voies accédons-nous à la conscience d'autrui, ces voies sont-elles suffisantes, et faut-il chercher ailleurs que dans la seule observation la véritable rencontre d'autrui ? Nous examinerons d'abord les voies d'accès à autrui (perception, analogie, empathie, langage) ; nous en montrerons ensuite les limites et l'opacité radicale de la conscience d'autrui ; nous soutiendrons enfin qu'autrui se connaît véritablement dans l'action commune et la reconnaissance, c'est-à-dire dans l'intersubjectivité vécue.
\end{corrige}

\courssub{Je me prépare au Bac}

\begin{corrige}[Exercice 9 — Dissertation type Baccalauréat Technique — Sujet 1]
\textbf{1. Introduction.} Dès la naissance, je suis entouré d'autrui : famille, voisins, camarades, chefs. Mais cette présence constante est ambiguë : les autres m'aident à vivre, et pourtant ils me jugent, m'imposent des règles, me surveillent. Autrui désigne l'autre conscience, semblable à la mienne ; la liberté est le pouvoir d'agir selon sa propre volonté, sans contrainte extérieure. Le sujet demande si la présence d'autrui entrave ce pouvoir ou au contraire le rend possible. Problématique : \textbf{autrui est-il un obstacle à ma liberté ?} Nous verrons d'abord qu'autrui apparaît comme une limitation de ma liberté ; puis qu'il en est en réalité la condition ; enfin que la liberté concrète s'exerce avec et par autrui.

\textbf{2. Première partie : autrui comme obstacle.} \emph{Argument 1 (Sartre)} : le regard d'autrui m'objective ; sous le regard, je deviens « être-pour-autrui », figé dans une identité que je n'ai pas choisie ; la honte en témoigne. « L'enfer, c'est les autres » : autrui me vole mon monde. \emph{Argument 2 (vie courante)} : la pression sociale et la surveillance limitent mes choix — au quartier comme sur les réseaux sociaux, la crainte du qu'en-dira-t-on conduit à renoncer à des projets, à se conformer ; l'individu sacrifie sa singularité pour être accepté.

\textbf{3. Deuxième partie : autrui comme condition.} \emph{Argument 1 (Hegel)} : la conscience de soi n'existe que reconnue ; dans la dialectique du Maître et de l'Esclave, la liberté pleine exige la reconnaissance mutuelle et égale — sans autrui, je ne suis même pas un « je ». \emph{Argument 2 (Levinas)} : le visage d'autrui m'ordonne « tu ne tueras point » et me constitue comme sujet responsable ; loin d'aliéner, cette responsabilité fonde ma dignité et m'arrache à l'égoïsme. Ajoutons que c'est par autrui que j'apprends le langage, les savoirs, les techniques : ma liberté effective (pouvoir agir) dépend de ce que les autres m'ont transmis.

\textbf{4. Troisième partie : synthèse.} Le conflit avec autrui est réel, mais il révèle un besoin : on ne lutte que pour être reconnu. La liberté n'est pas l'isolement de l'individu sans limites ; elle s'accomplit dans des relations réciproques : parole donnée et confiance (Ricœur), lois communes, coopération. Une liberté qui nie autrui se détruit elle-même ; une liberté qui reconnaît autrui devient concrète et durable.

\textbf{5. Conclusion.} Autrui est un obstacle apparent mais une condition réelle de ma liberté : c'est avec et par les autres que je deviens libre. Ouverture : cette leçon fonde la citoyenneté camerounaise — le vivre-ensemble, le respect des droits de la personne et le dialogue entre communautés sont la forme politique de la reconnaissance mutuelle.

\textbf{Barème indicatif (sur 20) :} Introduction (3 pts) — Partie 1 (4 pts) — Partie 2 (4 pts) — Partie 3 (4 pts) — Conclusion + ouverture (3 pts) — Qualité de la langue, rigueur philosophique, exemples camerounais (2 pts).

\textbf{Points de vigilance :} Ne pas confondre liberté (pouvoir de choisir) et licence (faire ce qu'on veut) ; citer explicitement Hegel (reconnaissance) et Sartre (regard/honte) ; articuler les parties par des connecteurs logiques (certes, cependant, en réalité) ; l'ouverture doit être ancrée dans le contexte camerounais (citoyenneté, vivre-ensemble).
\end{corrige}

\begin{corrige}[Exercice 10 — Explication de texte guidée — Sujet 2]
\begin{enumerate}
    \item \textbf{Identification.} Il s'agit de Jean-Paul Sartre, \emph{L'Être et le Néant} (1943). Thèse : le regard d'autrui me constitue comme objet et aliène ma liberté, de sorte que la relation à autrui est fondamentalement conflictuelle.
    \item \textbf{Analyse structurée.} Moment 1 : « par le regard d'autrui je deviens un objet » — autrui me saisit de l'extérieur et me fige. Moment 2 : « je perds ma maîtrise du monde » — le monde, qui s'organisait autour de moi, se réorganise autour d'autrui ; je suis dépossédé. Moment 3 : « l'enfer, c'est les autres » — conclusion existentielle : dépendre du regard d'autrui, c'est subir une condamnation permanente.
    \item \textbf{Concepts.} \emph{Être-pour-autrui} : mon être tel qu'il existe dans la conscience de l'autre, identité figée que je ne contrôle pas. \emph{Regard} : non pas seulement les yeux, mais la présence d'autrui qui me saisit comme objet (je peux me sentir regardé sans témoin visible). \emph{Honte} : conscience vive d'être cet objet devant autrui. \emph{Liberté} : pour Sartre, la conscience est spontanéité et projet ; le regard d'autrui la menace en la fixant.
    \item \textbf{Discussion.} Chez Hegel, la lutte des consciences débouche sur la reconnaissance mutuelle, condition de la liberté ; chez Sartre, le conflit semble sans issue. Mais la conflictualité n'est pas la seule vérité : l'empathie (Scheler) donne accès au vécu d'autrui sans objectivation ; l'amour et l'amitié accueillent l'autre comme sujet ; le travail commun et la parole échangée (Ricœur) construisent un monde partagé. Sartre décrit une structure réelle mais unilatérale de la relation à autrui.
\end{enumerate}

\textbf{Barème indicatif (sur 20) :} Identification (2 pts) — Analyse structurée (6 pts) — Explication concepts (6 pts) — Discussion critique (6 pts).

\textbf{Points de vigilance :} Ne pas résumer le texte mais l'analyser (moments logiques) ; définir les concepts \emph{dans le contexte du texte} (pas une définition de dictionnaire) ; la discussion doit confronter Sartre à un autre auteur au programme (Hegel, Levinas, Ricœur, Scheler) ; nuancer « l'enfer, c'est les autres » (citation de \emph{Huis clos}, pièce de théâtre, mais thèse ontologique dans \emph{L'Être et le Néant}).
\end{corrige}

\begin{corrige}[Exercice 11 — Question de synthèse et analyse de situation — Sujet 3]
\textbf{1. Synthèse théorique.} \emph{Voies d'accès.} Par la perception, je saisis le corps d'autrui, ses gestes et ses expressions. Par l'analogie (Mill), j'infère sa conscience à partir de la ressemblance de nos comportements. Par l'empathie (Scheler, Stein), je perçois originairement son vécu dans son expression même. Par le langage, autrui me dit directement ce qu'il pense et ressent : la parole est la voie la plus riche. \emph{Limites.} L'analogie reste une inférence incertaine (asymétrie entre mon vécu et ses signes, possibilité de la simulation) ; l'empathie risque la projection (j'attribue à autrui mes propres sentiments) ; le langage peut mentir ou créer le malentendu. Une opacité subsiste : je n'aurai jamais la conscience d'autrui de l'intérieur. \emph{Dépassement.} Autrui se connaît moins en le regardant qu'en agissant avec lui : le travail commun, la parole donnée, la reconnaissance mutuelle (Hegel, Ricœur) font de la connaissance d'autrui une pratique partagée, où chacun se révèle dans l'action.

\textbf{2. Application.} \emph{Formes d'intersubjectivité.} La palabre mobilise l'intersubjectivité communicative (échange de paroles et d'arguments), éthique (reconnaissance de l'autre comme interlocuteur légitime) et symbolique (cola partagée, ancêtres invoqués : rites qui réactivent un monde commun de significations). \emph{Insuffisance du seul langage juridique.} Un procès formel tranche un droit mais ne restaure pas la confiance ; chez Ricœur, le lien social repose sur la parole donnée et l'estime mutuelle : les rites du partage recréent cette confiance, condition d'un accord accepté durablement. \emph{Reconnaissance mutuelle.} Chez Hegel, la conscience de soi n'existe que reconnue : en s'asseyant ensemble, en rappelant les alliances matrimoniales, chaque village reconnaît l'autre comme partenaire et non comme ennemi ; cette reconnaissance réciproque transforme le conflit en compromis, et la violence en parole.

\textbf{Barème indicatif (sur 20) :} Synthèse théorique structurée en 3 temps (8 pts) — Application : identification formes intersubjectivité (3 pts), insuffisance langage juridique/Ricœur (4 pts), reconnaissance mutuelle/Hegel (3 pts) — Qualité rédaction, vocabulaire technique, ancrage camerounais (2 pts).

\textbf{Points de vigilance :} La synthèse doit être structurée (pas une liste d'auteurs) ; pour l'application, ne pas faire du droit coutumier mais de la philosophie (Ricœur, Hegel) ; montrer que le rite (cola, ancêtres) a une fonction épistémique et éthique (créer de la confiance, reconnaître l'altérité) ; utiliser les termes techniques : intersubjectivité symbolique, parole donnée, reconnaissance, monde commun.
\end{corrige}

\newpage

\coursec{Fiche bilan}

\begin{popfigure}
\begin{tikzpicture}[
  mindmap,
  grow cyclic,
  every node/.style={concept, execute at begin node=\hskip0pt},
  root concept/.append style={
    concept color=popblue, minimum size=2.5cm, text=white, font=\bfseries\large
  },
  level 1 concept/.append style={
    sibling angle=360/6, minimum size=1.8cm, text=white, font=\bfseries\normalsize
  },
  level 2 concept/.append style={
    sibling angle=45, minimum size=1.2cm, text=white, font=\small
  },
  concept color=popblue!80
]
\node[root concept] {Autrui}
  child[concept color=poporange] { node[concept, align=center]{1. Définition\\ \& Problématique}
    child { node[concept] {Autrui = autre moi} }
    child { node[concept] {Problème : accès à l'esprit d'autrui} }
  }
  child[concept color=popgreen] { node[concept, align=center]{2. Modes de\\ connaissance}
    child { node[concept, align=center]{Perception\\ (corps, comportement)} }
    child { node[concept, align=center]{Analogie\\ (argument par analogie)} }
    child { node[concept] {Limites : solipsisme} }
  }
  child[concept color=poppurple] { node[concept, align=center]{3. Empathie,\\ Langage, Communication}
    child { node[concept, align=center]{Einfühlung\\ (Husserl, Stein)} }
    child { node[concept] {Langage : pont intersubjectif} }
    child { node[concept] {Communication : partage de sens} }
  }
  child[concept color=poppink] { node[concept] {4. Intersubjectivité}
    child { node[concept] {Définition : communauté de sujets} }
    child { node[concept] {Formes : empathie, dialogue, travail, conflit} }
    child { node[concept] {Husserl, Schutz, Merleau-Ponty} }
  }
  child[concept color=popgold] { node[concept] {5. Autrui \& Liberté}
    child { node[concept, align=center]{Hegel : reconnaissance,\\ dialectique maître/esclave} }
    child { node[concept, align=center]{Sartre : regard, objet, conflit,\\ "L'enfer, c'est les autres"} }
    child { node[concept] {Obstacle OU condition ?} }
  }
  child[concept color=popblue] { node[concept] {6. Méthode \& Bac}
    child { node[concept] {Dissertation : thèse/antithèse/synthèse} }
    child { node[concept] {Commentaire : expliquer, analyser, discuter} }
    child { node[concept] {Exemples concrets (CM)} }
  };
\end{tikzpicture}
\legende{Carte mentale de la leçon « Autrui » — Première Technique}
\end{popfigure}

\begin{bilan}

\courssub{Définitions clés}
\begin{itemize}
  \item \cle{Autrui} : Autre sujet conscient, distinct de moi, centre de perspective propre (\emph{alter ego}).
  \item \cle{Solipsisme} : Thèse selon laquelle seul mon esprit existe ; l'existence d'autrui est indémontrable.
  \item \cle{Argument par analogie} : Inférence du type « Je me comporte ainsi quand je pense → Autrui se comporte ainsi → Autrui pense ». (J. S. Mill).
  \item \cle{Empathie (Einfühlung)} : Captation vécue de l'expérience d'autrui sans confusion (Husserl, Stein) ; distincte de la sympathie (sentiment affectif).
  \item \cle{Intersubjectivité} : Structure originaire du monde commun ; constitution mutuelle des sujets (Husserl, Merleau-Ponty, Schutz).
  \item \cle{Reconnaissance (Hegel)} : Autrui est nécessaire pour que je me reconnaisse comme conscience de soi (Dialectique Maître/Esclave).
  \item \cle{Le Regard (Sartre)} : Autrui me fige en objet (\emph{pour-autrui}), m'aliène, mais fonde ma facticité et ma responsabilité.
\end{itemize}

\courssub{Thèses \& Auteurs à connaître (Tableau de synthèse)}
\begin{center}
\begin{tabularx}{\linewidth}{|l|X|X|}\hline
\rowcolor{popblueL}
\textbf{Auteur / Courant} & \textbf{Thèse centrale sur Autrui} & \textbf{Concept clé} \\ \hline
\cle{Argument par analogie} (J.S. Mill) & Connaissance par inférence probabiliste à partir de mon cas propre. & Analogie comportementale \\ \hline
\cle{Phenoménologie} (Husserl, Stein) & Apperception par analogie (\emph{Paarung}) + Empathie comme acte intentionnel originaire. & \emph{Einfühlung}, Constitution intersubjective \\ \hline
\cle{Merleau-Ponty} & Corps propre et corps d'autrui s'entrelacent ; intercorporéité primordiale. & Chiasme, Intercorporéité \\ \hline
\cle{Hegel} & Autrui est \textbf{condition} de ma liberté : je ne suis conscience de soi que par reconnaissance mutuelle. & Dialectique Maître/Esclave, Reconnaissance \\ \hline
\cle{Sartre} & Autrui est \textbf{obstacle} (conflit originel) : le Regard me chosifie. Mais aussi médiation nécessaire pour ma vérité. & Le Regard, Être-pour-autrui, Conflit \\ \hline
\cle{Buber} & Relation \emph{Je-Tu} (présence, réciprocité) vs \emph{Je-Ça} (instrumentalisation). & Dialogue, Relation \\ \hline
\end{tabularx}
\end{center}

\courssub{Méthodes express (Bac)}
\begin{enumerate}
  \item \textbf{Dissertation :} Analyser le sujet → Problématiser (ex: « Autrui : limite ou ouverture ? ») → Plan dialectique (Thèse : Autrui obstacle / Antithèse : Autrui condition / Synthèse : Dialectique reconnaissance/aliénation) → Exemples précis (regard dans la rue, dialogue familial, travail en équipe, conflit au Cameroun).
  \item \textbf{Commentaire :} 1. Contextualiser (auteur, œuvre, époque). 2. Résumer la thèse en une phrase. 3. Expliquer les arguments pas à pas (vocabulaire technique : \emph{analogie, empathie, objectivation, reconnaissance}). 4. Discuter critique (portée, limites, contre-exemple). 5. Conclure sur la portée philosophique.
  \item \textbf{Application concrète (Cameroun) :} « Le respect des aînés » (reconnaissance sociale), « La file d'attente au guichet » (intersubjectivité normative), « Le regard du policier / du chef » (aliénation sartrienne / pouvoir hégélien).
\end{enumerate}

\end{bilan}

\begin{aretenir}[Checklist avant le Bac]
\begin{itemize}
  \item $\square$ Savoir définir \textbf{Autrui}, \textbf{Solipsisme}, \textbf{Intersubjectivité}, \textbf{Empathie}, \textbf{Reconnaissance}.
  \item $\square$ Expliquer l'\textbf{argument par analogie} (schéma : Moi $\rightarrow$ Autrui) et ses limites (risque solipsisme).
  \item $\square$ Différencier \textbf{perception} (corps visible) et \textbf{apperception/empathie} (esprit invisible).
  \item $\square$ Citer \textbf{Husserl/Stein} : empathie comme acte intentionnel originaire (pas inférence).
  \item $\square$ Citer \textbf{Merleau-Ponty} : intercorporéité, corps comme « double » de l'autre.
  \item $\square$ Restituer la \textbf{dialectique Maître/Esclave (Hegel)} : lutte, risque de mort, travail, reconnaissance mutuelle.
  \item $\square$ Analyser le \textbf{Regard chez Sartre} : objectivation, honte, conflit, « L'enfer, c'est les autres » (\emph{Huis Clos}).
  \item $\square$ Opposer \textbf{Hegel (condition)} vs \textbf{Sartre (obstacle/condition ambiguë)}.
  \item $\square$ Maîtriser la structure type d'une \textbf{dissertation} (Intro : accroche, définitions, problème, thèse, plan ; Développement équilibré ; Conclusion : bilan + ouverture).
  \item $\square$ Avoir 2-3 \textbf{exemples camerounais} prêts (famille, administration, travail, rites) pour illustrer chaque thèse.
  \item $\square$ Savoir \textbf{commenter un texte} : thèse, arguments, concepts, discussion critique, conclusion.
\end{itemize}
\end{aretenir}

\findecours

\end{document}
